\documentclass[11pt]{article}
\usepackage[T1]{fontenc}
\usepackage[utf8]{inputenc}
\usepackage[margin=1in]{geometry}
\usepackage{amsmath,amssymb,amsthm,mathtools}
\usepackage{xurl}
\usepackage{enumitem}
\usepackage[colorlinks=true,linkcolor=blue,citecolor=blue,urlcolor=blue]{hyperref}
\numberwithin{equation}{section}
\newtheorem{theorem}{Theorem}[section]
\newtheorem{lemma}[theorem]{Lemma}
\newtheorem{proposition}[theorem]{Proposition}
\newtheorem{corollary}[theorem]{Corollary}
\newtheorem{counterexample}[theorem]{Counterexample}
\theoremstyle{definition}
\newtheorem{definition}[theorem]{Definition}
\newtheorem{example}[theorem]{Example}
\theoremstyle{remark}
\newtheorem{remark}[theorem]{Remark}
\DeclareMathOperator{\Var}{Var}
\DeclareMathOperator{\Prob}{P}
\DeclareMathOperator{\supp}{supp}
\newcommand{\E}{\mathbb E}
\newcommand{\N}{\mathbb N}
\newcommand{\Z}{\mathbb Z}
\newcommand{\1}{\mathbf 1}
\title{Zeta-Law Entropy, Gamma Curvature, and Laplace-Measure Theory}
\author{Pudim AI Project}
\date{}
\begin{document}
\maketitle

\begin{abstract}
We develop a unified framework for zeta-law entropy, gamma-function curvature, complete monotonicity, Bernstein and Stieltjes representations, and Laplace-measure comparison.  The common mechanism is to convert analytic inequalities into positivity statements for kernels, moments, determinants, and logarithmic derivatives.  The resulting structural theorems yield sharp monotonicity, convexity, endpoint, and parameter-range conclusions for externally stated problems in special-function theory.
\end{abstract}

\section{Introduction}
Many inequalities for special functions become transparent only after the expression is rewritten as a positive object: a probability moment, a Laplace kernel, a Stieltjes transform, a Bernstein derivative, or a determinant of moments.  This paper organizes a collection of solved source problems around that common mechanism.  The point is not that the applications are chronologically adjacent, but that the same proof grammar repeatedly turns analytic claims into positivity statements.
The main text is therefore arranged by method.  Each section introduces the relevant mechanism, states the reusable results, and then records the integrated applications that follow from them.  Stable APP labels are retained only as public source-problem annotations; they do not determine theorem numbering or the order of exposition.

\section{Motivation and main mechanisms}
The recurring strategy has three steps.  First, normalize the source expression until the object being tested is a kernel, moment sequence, or logarithmic derivative.  Second, prove positivity, complete monotonicity, or a finite determinant inequality in that normal form.  Third, transport the result back to the source problem, often obtaining a sharp endpoint, a monotonicity interval, or a counterexample.
This organization also explains the negative results.  When the positive-kernel route cannot hold, an endpoint expansion, a rational specialization, or an explicit numerical witness often identifies exactly where the proposed statement fails.

\section{Notation and preliminaries}
The Riemann zeta function is denoted by \(\zeta\), Euler's gamma function by \(\Gamma\), the digamma function by \(\psi=\Gamma'/\Gamma\), and Euler's constant by \(\gamma=-\psi(1)\).  Higher derivatives of \(\psi\) are the polygamma functions.  We write \(\N=\{1,2,\ldots\}\), and use \(\E\), \(\Prob\), and \(\Var\) for expectation, probability, and variance.
The Pochhammer, or rising factorial, is \((a)_m=a(a+1)\cdots(a+m-1)\), with \((a)_0=1\).  The notation \([z^m]F(z)\) denotes coefficient extraction.
A smooth function \(f:(0,\infty)\to\mathbb R\) is completely monotone if \((-1)^m f^{(m)}(x)\ge0\) for every \(m\ge0\) and \(x>0\).  A positive function is logarithmically completely monotone when its logarithm has the same alternating derivative sign pattern.  A nonnegative function \(B\) is a Bernstein function if \(B'\) is completely monotone.  A complete Bernstein function is a Bernstein function with the standard Stieltjes-density representation.  A Stieltjes function is a positive superposition of kernels \((x+t)^{-1}\).  A generalized Stieltjes function is a positive superposition of kernels \((x+t)^{-\lambda}\) for a fixed order \(\lambda>0\).

\section{Zeta-law kernels and tail partitions}
The first layer treats zeta-normalized expressions as probability, kernel, or tail-partition identities.  This viewpoint turns several inequalities into positivity statements for moments, finite shadows, and reciprocal partition functions.
\begin{proposition}[Qi-Agarwal divisor-polygamma \(f_{2}\) is completely monotone]
\label{res:qa-divisor-polygamma-f2-cm}
The second divisor-polygamma sum \(f_2(x)=[\psi'(x)]^2+\psi''(x)\) is completely monotone on \((0,\infty)\).
\end{proposition}
\begin{proof}
For every \(k\ge1\), the standard polygamma sign theorem gives
\[
(-1)^{k+1}\psi^{(k)}(x)
=
\int_0^\infty \frac{t^k e^{-xt}}{1-e^{-t}}\,dt,
\]
so \((-1)^{k+1}\psi^{(k)}\) is completely monotone on \((0,\infty)\).

If \(n\) is odd and \(km=n\), then \(k\) and \(m\) are both odd. Hence \(\psi^{(k)}\) itself is a positive completely monotone function, and
\[
\bigl[\psi^{(k)}(x)\bigr]^m
\]
is completely monotone by finite product closure of positive completely monotone functions. Summing over the finitely many divisor pairs \(km=n\) preserves complete monotonicity. Therefore \(f_n\) is completely monotone for every odd \(n\).

For \(n=2\), the divisor pairs are \((k,m)=(1,2)\) and \((2,1)\), so
\[
f_2(x)=\bigl[\psi'(x)\bigr]^2+\psi''(x).
\]
The same Qi--Agarwal source records the known sharp theorem that
\[
\bigl[\psi'(x)\bigr]^2+\lambda\psi''(x)
\]
is completely monotonic on \((0,\infty)\) if and only if \(\lambda\le1\). Taking \(\lambda=1\) gives \(f_2\in CM(0,\infty)\).

Thus the source's proposed even-parity clause, "\(f_{2\ell}\) is not completely monotonic for all \(\ell\in\mathbb N\)", is false already at \(\ell=1\). The source problem is solved negatively/correctively:

the odd half is true;
the stated even half is false;
the corrected even frontier begins at \(n=4\), not \(n=2\).
\end{proof}

\begin{lemma}[incomplete beta tail derivative has positive discrete Laplace representation for lambda...]
\label{res:incomplete-beta-tail-derivative-cm-certificate}
For \(b>0\) and \(0<\lambda\le1\), the derivative of \(I_{b,\lambda}(x)=B(b,\lambda)-B(b,\lambda;e^{-x})\) is completely monotone; for \(0<\lambda<1\), \(I_{b,\lambda}'(x)=\sum_{n=0}^\infty (1-\lambda)_n e^{-(b+n)x}/n!\).
\end{lemma}
\begin{proof}
For \(b>0\) and \(0<\lambda\le1\), define
\[
I_{b,\lambda}(x)=B(b,\lambda)-B(b,\lambda;e^{-x}).
\]
With the lower incomplete beta convention
\[
B(b,\lambda;z)=\int_0^z u^{b-1}(1-u)^{\lambda-1}\,du,
\]
one has
\[
I_{b,\lambda}(x)=\int_0^x e^{-bt}(1-e^{-t})^{\lambda-1}\,dt,
\]
and \(I_{b,\lambda}\) is a Bernstein function on \((0,\infty)\).

The substitution \(u=e^{-t}\) gives
\[
B(b,\lambda)-B(b,\lambda;e^{-x})
=\int_{e^{-x}}^1u^{b-1}(1-u)^{\lambda-1}\,du
=\int_0^x e^{-bt}(1-e^{-t})^{\lambda-1}\,dt.
\]
Therefore
\[
I_{b,\lambda}'(x)=e^{-bx}(1-e^{-x})^{\lambda-1}.
\]
For \(\lambda=1\), this is \(e^{-bx}\), a completely monotone function.

For \(0<\lambda<1\), put \(c=1-\lambda\). Then \(c\in(0,1)\), and
\[
(1-e^{-x})^{-c}
=\sum_{n=0}^\infty \frac{(c)_n}{n!}e^{-nx}
\]
with positive coefficients and locally uniform convergence on \((0,\infty)\). Thus
\[
I_{b,\lambda}'(x)
=\sum_{n=0}^\infty \frac{(c)_n}{n!}e^{-(b+n)x},
\]
the Laplace transform of the positive discrete measure
\[
\sum_{n=0}^\infty \frac{(c)_n}{n!}\delta_{b+n}.
\]
So \(I_{b,\lambda}'\) is completely monotone. Since \(I_{b,\lambda}(x)\ge0\), \(I_{b,\lambda}\in C^\infty(0,\infty)\), and \(I_{b,\lambda}'\) is completely monotone, \(I_{b,\lambda}\) is a Bernstein function.
\end{proof}

\begin{lemma}[log function signed derivative polynomial normal form and tail vanishing]
\label{res:log-function-signed-derivative-tail-vanishing}
For every \(f\in\mathcal F_{1,n}\), where \(f(x)=p(\log x)/x\), and every \(k\ge0\), the signed derivative \(L^k f\) with \(L=-d/dx\) has the form \(x^{-k-1}p_k(\log x)\) for a polynomial \(p_k\), and therefore \(L^k f(x)\to0\) as \(x\to\infty\).
\end{lemma}
\begin{proof}
We first prove that for each \(k\ge0\),
\[
L^k f(x)=x^{-k-1}p_k(\log x)
\]
for some polynomial \(p_k\).

For \(k=0\), this is the definition with \(p_0=p\).  Suppose
\[
L^k f(x)=x^{-k-1}p_k(\log x).
\]
Then
\[
\begin{aligned}
L^{k+1}f(x)
&=-\frac{d}{dx}\left(x^{-k-1}p_k(\log x)\right)\\
&=x^{-k-2}\left((k+1)p_k(\log x)-p_k'(\log x)\right).
\end{aligned}
\]
Thus \(p_{k+1}(y)=(k+1)p_k(y)-p_k'(y)\), which is again a polynomial.

Since every polynomial in \(\log x\) grows slower than every positive power of \(x\),
\[
x^{-k-1}p_k(\log x)\to0
\qquad (x\to\infty).
\]
Therefore
\[
\lim_{x\to\infty}L^k f(x)=0
\]
for every \(k\ge0\).

Assume \(f\in D_{k+1}\).  Then
\[
L^{k+1}f(x)\ge0
\]
on \((0,\infty)\).  Put
\[
g(x)=L^k f(x).
\]
By definition of \(L\),
\[
-g'(x)=L^{k+1}f(x)\ge0.
\]
Using \(g(x)\to0\) as \(x\to\infty\), integrate from \(x\) to \(R\):
\[
g(x)-g(R)=\int_x^R L^{k+1}f(t)\,dt.
\]
Letting \(R\to\infty\) gives
\[
L^k f(x)=g(x)=\int_x^\infty L^{k+1}f(t)\,dt\ge0.
\]
Hence \(f\in D_k\).  Therefore
\[
D_{k+1}\subseteq D_k
\]
for all \(k\ge0\) and all \(n\).
\end{proof}

\section{Gamma and polygamma curvature}
Gamma, digamma, beta, and polygamma questions are organized by logarithmic derivatives and Laplace kernels.  The guiding principle is that complete-monotonicity claims should be reduced to sign-controlled integral kernels or endpoint obstructions.
\begin{theorem}[Simon gamma quotient complement is CM for alpha in (0,1)]
\label{res:simon-gamma-quotient-complement-cm}
For every \(0<\alpha<1\), the function \(1-\Gamma(x+\alpha)/(\Gamma(x)x^\alpha)\) is completely monotone on \((0,\infty)\).
\end{theorem}
\begin{proof}
For \(0<\alpha<1\), put
\[
F_\alpha(x)=\frac{\Gamma(x+\alpha)}{\Gamma(x)x^\alpha},\qquad x>0.
\]
Then \(F_\alpha\) is a Bernstein function. In fact, \(1-F_\alpha\) is completely monotone.

The stronger route asserting that \(F_\alpha\) is a complete Bernstein function is false.

Let \(\beta=1-\alpha\in(0,1)\). Since \(\Gamma(x+1)=x\Gamma(x)\), the beta integral gives
\[
F_\alpha(x)
=x^\beta\frac{\Gamma(x+\alpha)}{\Gamma(x+1)}
=\frac{x^\beta}{\Gamma(\beta)}
\int_0^\infty e^{-xu}(e^u-1)^{-\alpha}\,du.
\]
Also
\[
1=\frac{x^\beta}{\Gamma(\beta)}
\int_0^\infty e^{-xu}u^{-\alpha}\,du.
\]
Thus
\[
1-F_\alpha(x)
=\frac{x^\beta}{\Gamma(\beta)}
\int_0^\infty e^{-xu}\left(u^{-\alpha}-(e^u-1)^{-\alpha}\right)\,du.
\]

Define
\[
J_\alpha(t)=\int_0^t (t-u)^{\alpha-1}(e^u-1)^{-\alpha}\,du.
\]
By Laplace convolution,
\[
\mathcal L[J_\alpha](x)
=\Gamma(\alpha)x^{-\alpha}
\int_0^\infty e^{-xu}(e^u-1)^{-\alpha}\,du.
\]
Therefore
\[
F_\alpha(x)
=\frac{x}{\Gamma(\alpha)\Gamma(1-\alpha)}\mathcal L[J_\alpha](x).
\]

After the substitution \(u=tv\),
\[
J_\alpha(t)=
\int_0^1
(1-v)^{\alpha-1}
\left(\frac{t}{e^{tv}-1}\right)^\alpha\,dv.
\]
For fixed \(v\in(0,1]\), the function
\[
t\mapsto \frac{t}{e^{tv}-1}
\]
is decreasing because
\[
\frac{d}{dt}\frac{t}{e^{tv}-1}
=
\frac{e^{tv}-1-tv e^{tv}}{(e^{tv}-1)^2}\le0.
\]
The integrand is dominated by \((1-v)^{\alpha-1}v^{-\alpha}\), which is integrable on \((0,1)\). Hence \(J_\alpha\) is decreasing. Moreover,
\[
J_\alpha(0+)
=\int_0^1(1-v)^{\alpha-1}v^{-\alpha}\,dv
=\Gamma(\alpha)\Gamma(1-\alpha),
\]
and \(J_\alpha(t)\to0\) as \(t\to\infty\).

Let \(\mu_\alpha=-dJ_\alpha\). This is a positive finite measure with total mass
\[
\mu_\alpha([0,\infty))=\Gamma(\alpha)\Gamma(1-\alpha).
\]
Stieltjes integration by parts gives, for \(x>0\),
\[
\int_0^\infty e^{-xt}\,\mu_\alpha(dt)
=
\Gamma(\alpha)\Gamma(1-\alpha)-x\mathcal L[J_\alpha](x).
\]
Using the formula for \(F_\alpha\),
\[
1-F_\alpha(x)
=
\frac{1}{\Gamma(\alpha)\Gamma(1-\alpha)}
\int_0^\infty e^{-xt}\,\mu_\alpha(dt).
\]
Thus \(1-F_\alpha\) is completely monotone.

Consequently,
\[
F_\alpha'(x)
=
\frac{1}{\Gamma(\alpha)\Gamma(1-\alpha)}
\int_0^\infty t e^{-xt}\,\mu_\alpha(dt),
\]
so \(F_\alpha'\) is completely monotone. Since \(F_\alpha(x)>0\), \(F_\alpha\) is a Bernstein function.

For \(t>0\), differentiating the \(v\)-integral under the same endpoint domination gives \(\mu_\alpha(dt)=m_\alpha(t)\,dt\), where
\[
m_\alpha(t)
=
\alpha\int_0^1
(1-v)^{\alpha-1}
\left(\frac{t}{e^{tv}-1}\right)^\alpha
\left(
\frac{v e^{tv}}{e^{tv}-1}-\frac1t
\right)\,dv.
\]
The bracket is positive, because with \(a=tv>0\),
\[
\frac{a e^a}{e^a-1}>1.
\]
Therefore the Laplace-density route in the argument is also closed.

The case \(\alpha=1/2\) is the specialization
\[
F_{1/2}(x)=\frac{\Gamma(x+1/2)}{\Gamma(x)\sqrt{x}},
\]
and the same complement-CM representation proves that \(F_{1/2}\) is Bernstein.

For \(z=-r+i0\) with \(r\in(\alpha,1)\), take the upper-lip branch of \(z^{-\alpha}\). Then
\[
F_\alpha(-r+i0)
=
\frac{\Gamma(\alpha-r)}{\Gamma(-r)}r^{-\alpha}e^{-i\pi\alpha}.
\]
Here \(\Gamma(\alpha-r)<0\) and \(\Gamma(-r)<0\), so the real gamma ratio is positive. Hence
\[
\operatorname{Im} F_\alpha(-r+i0)<0.
\]
A complete Bernstein function has the Pick property and maps the upper half-plane into itself. Therefore \(F_\alpha\) is not a complete Bernstein function for \(0<\alpha<1\).
\end{proof}

\begin{proposition}[Baricz V\_q positive Laplace representation in q plus one]
\label{res:baricz-vq-q-laplace-normal-form}
For every \(x>0\) and \(q>-1\), Baricz's function satisfies \(V_q(x)=\pi^{-1/2}\int_0^\infty s^{-1/2}e^{-x^2s}(1+s)^{-(q+1)}\,ds\).
\end{proposition}
\begin{proof}
Put \(u=t^2-x^2\). Then \(t=(u+x^2)^{1/2}\) and
\[
dt=\frac{du}{2(u+x^2)^{1/2}}.
\]
Since \(e^{x^2}e^{-t^2}=e^{-u}\), the factor \(2\) in the definition cancels the \(1/2\) from \(dt\), giving
\[
V_q(x)=\frac1{\Gamma(q+1)}
\int_0^\infty e^{-u}u^q(u+x^2)^{-1/2}\,du.
\]
This integral is finite for \(q>-1\): near \(0\) the integrand is \(O(u^q)\), and at infinity it is exponentially decaying.

For \(r>0\),
\[
r^{-1/2}=\frac1{\sqrt\pi}\int_0^\infty s^{-1/2}e^{-rs}\,ds.
\]
With \(r=u+x^2\), Tonelli's theorem applies because the integrand is nonnegative. Thus
\[
\begin{aligned}
V_q(x)
&=\frac1{\Gamma(q+1)\sqrt\pi}
\int_0^\infty\int_0^\infty
e^{-u}u^q s^{-1/2}e^{-(u+x^2)s}\,ds\,du\\
&=\frac1{\Gamma(q+1)\sqrt\pi}
\int_0^\infty s^{-1/2}e^{-x^2s}
\left(\int_0^\infty u^q e^{-(1+s)u}\,du\right)\,ds\\
&=\frac1{\sqrt\pi}\int_0^\infty
s^{-1/2}e^{-x^2s}(1+s)^{-(q+1)}\,ds.
\end{aligned}
\]
The last step uses
\[
\int_0^\infty u^q e^{-(1+s)u}\,du
=\frac{\Gamma(q+1)}{(1+s)^{q+1}}.
\]

This proves the claim.

Let \(a=q+1>0\). The normal form becomes
\[
V_{a-1}(x)=\frac1{\sqrt\pi}\int_0^\infty
s^{-1/2}e^{-x^2s}(1+s)^{-a}\,ds.
\]
For each integer \(n\ge0\), differentiation under the integral gives
\[
(-1)^n\frac{d^n}{da^n}V_{a-1}(x)
=\frac1{\sqrt\pi}\int_0^\infty
(\log(1+s))^n s^{-1/2}e^{-x^2s}(1+s)^{-a}\,ds\ge0.
\]
The differentiation is justified locally uniformly in \(a>0\): near \(s=0\), \((\log(1+s))^n s^{-1/2}=O(s^{n-1/2})\), and as \(s\to\infty\), the factor \(e^{-x^2s}\) dominates every logarithmic and polynomial factor.

Thus \(a\mapsto V_{a-1}(x)\) is completely monotone on \((0,\infty)\). Equivalently, after the change of variables \(y=\log(1+s)\), it is a positive Laplace transform in \(a\).

This proves the claim.

For \(a>0\), write
\[
F_x(a)=V_{a-1}(x)
=\int_{(0,\infty)} e^{-ay}\,d\nu_x(y),
\]
where \(\nu_x\) is the positive pushforward of
\[
\frac1{\sqrt\pi}s^{-1/2}e^{-x^2s}\,ds
\]
under \(y=\log(1+s)\). This measure is not concentrated at one point because it has positive density on \(s>0\), hence on \(y>0\).

For \(a_1\ne a_2\) and \(0<\lambda<1\), strict Holder inequality gives
\[
F_x(\lambda a_1+(1-\lambda)a_2)
<
F_x(a_1)^\lambda F_x(a_2)^{1-\lambda}.
\]
Strictness holds because \(e^{-a_1y}\) and \(e^{-a_2y}\) are not proportional \(\nu_x\)-almost everywhere unless \(\nu_x\) is supported at a single \(y\).

Thus \(a\mapsto F_x(a)\) is strictly log-convex on \((0,\infty)\). Since \(a=q+1\), the shift preserves strict log-convexity, so \(q\mapsto V_q(x)\) is strictly log-convex on \((-1,\infty)\) for every fixed \(x>0\).

This proves the claim.
\end{proof}

\begin{proposition}[Chen Choi Gurland Euler product sum normal form]
\label{res:chen-choi-gurland-euler-sum-normal-form}
For \(F(x)=\Gamma(1/x)\Gamma(3/x)/\Gamma(2/x)^2\), \(\log F(x)=\log(4/3)+\sum_{m=1}^{\infty}A(mx)\), with locally uniform convergence of all derivative series on compact subintervals of \((0,\infty)\).
\end{proposition}
\begin{proof}
Let
\[
A(u)=2\log(u+2)-\log(u+1)-\log(u+3).
\]
Then \(A(u)>0\) for \(u>0\), since
\[
A(u)=\log\frac{(u+2)^2}{(u+1)(u+3)}
\]
and the numerator exceeds the denominator by \(1\).
Moreover
\[
A'(u)=\frac2{u+2}-\frac1{u+1}-\frac1{u+3}
=-\frac{2}{(u+1)(u+2)(u+3)}.
\]
Using
\[
\int_0^\infty e^{-ut}e^{-at}\,dt=\frac1{u+a},
\]
we get
\[
-A'(u)=\int_0^\infty e^{-ut}e^{-t}(1-e^{-t})^2\,dt.
\]
Since \(A(u)\to0\) as \(u\to\infty\), integration from \(u\) to infinity gives
\[
A(u)=\int_0^\infty e^{-ut}\frac{e^{-t}(1-e^{-t})^2}{t}\,dt.
\]
The density is positive, locally integrable at \(0\), and exponentially decaying at infinity. Hence \(A\) is strictly completely monotone.

Weierstrass' product gives
\[
\log\Gamma(z)=-\log z-\gamma z-
\sum_{m=1}^{\infty}\left(\log(1+z/m)-z/m\right).
\]
Put \(z=1/x\). In
\[
\log\Gamma(z)+\log\Gamma(3z)-2\log\Gamma(2z),
\]
the linear terms cancel. The logarithmic terms give \(\log(4/3)\). The summand at index \(m\) becomes
\[
\log\frac{(1+2z/m)^2}{(1+z/m)(1+3z/m)}
=\log\frac{(mx+2)^2}{(mx+1)(mx+3)}=A(mx).
\]
Thus
\[
\log F(x)=\log\frac43+\sum_{m=1}^{\infty}A(mx).
\]
For every compact \(x\in[\varepsilon,M]\) and every \(n\ge0\),
\[
\frac{d^n}{dx^n}A(mx)=m^nA^{(n)}(mx)=O_{n,\varepsilon}(m^{-2}),
\]
because the three-point second difference in \(A^{(n)}\) is \(O((mx)^{-n-2})\). Therefore the series may be differentiated termwise locally uniformly.

For \(n=0\), \(\log F(x)>0\) follows from \(\log(4/3)>0\) and \(A(mx)>0\). For \(n\ge1\),
\[
(-1)^n(\log F)^{(n)}(x)=
\sum_{m=1}^{\infty}m^n(-1)^nA^{(n)}(mx)>0.
\]
Therefore \(\log F\) is strictly completely monotone on \((0,\infty)\), proving Chen-Choi Conjecture 1.
\end{proof}

\begin{proposition}[Wakrim W symbol Bernstein exact beta range zero to one]
\label{res:wakrim-w-symbol-exact-bf-range}
For \(0<\alpha<1\) and \(\beta\ge0\), the Wakrim W-symbol \(\Phi_{\alpha,\beta}(s)=s^\alpha(1+(1-\alpha)s^{\alpha-1})^{-\beta}\) is a Bernstein function on \((0,\infty)\) if and only if \(0\le\beta\le1\).
\end{proposition}
\begin{proof}
Set \(a=1-\alpha\in(0,1)\). Then
\[
\Phi_{\alpha,\beta}(s)
=s^{1-a}\left(1+a s^{-a}\right)^{-\beta}
=\frac{s^{1-a+a\beta}}{(s^a+a)^\beta}.
\]
Write
\[
p=1-a+a\beta=\alpha+(1-\alpha)\beta.
\]
Differentiating gives
\[
\Phi_{\alpha,\beta}'(s)
=s^{p-1}(s^a+a)^{-\beta-1}
\left[p(s^a+a)-\beta a s^a\right].
\]
Since \(p-\beta a=1-a=\alpha\),
\[
\Phi_{\alpha,\beta}'(s)
=s^{-a(1-\beta)}(s^a+a)^{-\beta-1}
\left[\alpha s^a+a(\alpha+a\beta)\right].
\]
With \(y=s^a\), this becomes
\[
\Phi_{\alpha,\beta}'(s)=H(s^a),
\]
where
\[
H(y)=
y^{\beta-1}(y+a)^{-\beta}
\left(\alpha+\frac{a^2\beta}{y+a}\right).
\]

For \(0<\beta\le1\), each factor in \(H\) is completely monotone on \((0,\infty)\): \(y^{\beta-1}=y^{-(1-\beta)}\), \((y+a)^{-\beta}\), and the positive sum \(\alpha+a^2\beta/(y+a)\). Therefore \(H\) is completely monotone by product and positive-sum closure.

The map \(s\mapsto s^a\) is Bernstein for \(0<a<1\). The standard composition theorem says that if \(f\) is completely monotone and \(g\) is Bernstein, then \(f\circ g\) is completely monotone. Hence \(\Phi_{\alpha,\beta}'\) is completely monotone, and \(\Phi_{\alpha,\beta}\) is Bernstein for \(0<\beta\le1\). The endpoint \(\beta=0\) is \(\Phi_{\alpha,0}(s)=s^\alpha\), also Bernstein.

Combining this local positive range with Wakrim's source proof that \(\Phi_{\alpha,\beta}\) is not Bernstein for \(\beta>1\), the exact range is
\[
\Phi_{\alpha,\beta}\in BF
\quad\Longleftrightarrow\quad
0\le\beta\le1.
\]
\end{proof}

\begin{lemma}[negative derivative of completely monotone function is completely monotone]
\label{res:chiu-yin-cm-derivative-reflection-lemma}
If \(g\) is completely monotone on \((0,\infty)\), then \(-g'\) is completely monotone wherever the derivative exists.
\end{lemma}
\begin{proof}
Chiu--Yin Remark 3.5 conjectures the converse of Lemma 3.1: in the Sparre Andersen model, complete monotonicity of the ascending ladder-height density should imply complete monotonicity of the claim-size density. This pass proves only the classical Cramer--Lundberg/equilibrium-ladder slice.

\emph{Setup.}
In the classical Cramer--Lundberg/equilibrium formulation, let \(H\) be the claim-size distribution, let
\[
\overline H(x)=1-H(x),
\]
and assume \(H\) is absolutely continuous on \((0,\infty)\). Let
\[
m=\int_0^\infty \overline H(u)\,du
\]
be the finite mean, with \(0<m<\infty\). The equilibrium excess cdf and density are
\[
F_e(x)=\frac1m\int_0^x\overline H(u)\,du,
\qquad
f_e(x)=F_e'(x)=\frac{\overline H(x)}{m}.
\]
In the classical ladder normalization under discussion, the ascending ladder-height density is this equilibrium density:
\[
f_+(x)=f_e(x)=\frac{\overline H(x)}{m}.
\]

If \(g\) is completely monotone on \((0,\infty)\), then for every \(n\ge0\),
\[
(-1)^n g^{(n)}(x)\ge0.
\]
Therefore
\[
(-1)^n(-g')^{(n)}(x)=(-1)^{n+1}g^{(n+1)}(x)\ge0,
\]
so \(-g'\) is completely monotone.

Assume \(f_+\) is completely monotone. Then \(\overline H=m f_+\) is completely monotone. Define the smooth density representative
\[
h_*(x)=-\overline H'(x)=-m f_+'(x).
\]
By the reflection lemma, \(-f_+'\) is completely monotone; multiplying by \(m>0\) preserves complete monotonicity. Hence \(h_*\) is completely monotone.

Since densities are a.e. representatives, the conclusion is that \(H\) admits a completely monotone density version \(h_*\). Thus the converse of Chiu--Yin Lemma 3.1 is true in the classical equilibrium-ladder slice in this density-version sense.

This does not solve the full Sparre Andersen converse in Remark 3.5. The general case includes the weak descending-ladder renewal factor in the representation of \(F^+\), so the original claim density is not recovered by a single derivative of the ladder-height density.
\end{proof}

The next result applies the preceding method to the fixed source problem \textbf{APP-0021}.
\begin{corollary}[APP-0021: Szabo's cutoff problem has \(\alpha\_0=1\)]
\label{res:dagum-c-endpoints-and-upper-gate}
For the Dagum threshold \(c(\beta)=\inf\{\alpha\ge0:x^{-\alpha}(1+x^\beta)^{-1}\in CM\}\), one has \(c(1)=0\), \(c(2)=1\), and \(c(\beta)\le\beta/2\) for \(1<\beta\le2\).

This resolves the fixed application label \textbf{APP-0021}: Determine the exact cutoff \(\alpha\_0\) for \(y^\alpha H\_d(y)\).
\emph{Source reference.} \cite{app-app-0021-source}.
\end{corollary}
\begin{proof}
The source theorem proves complete monotonicity for \(\alpha\le1\).  Conversely,
as \(y\downarrow0\),
\[
  H_d(y)={1-d\over y}+O(1).
\]
Hence, for \(\alpha>1\),
\[
  y^\alpha H_d(y)=(1-d)y^{\alpha-1}+O(y^\alpha),
\]
whose derivative is positive near zero.  A completely monotone function must be
non-increasing, so no \(\alpha>1\) is admissible.
\end{proof}

\begin{lemma}[Szabo psi difference exact alpha0 equals one]
\label{res:szabo-psi-difference-alpha0-exact-one}
For \(0<d<1\), the sharp cutoff for complete monotonicity of \(y^\alpha[\psi(y+d)-\psi(y)-d/y]\) is \(\alpha_0=1\).
\end{lemma}
\begin{proof}
Let \(0<d<1\), \(y>0\), and
\[
H_d(y)=\psi(y+d)-\psi(y)-\frac{d}{y}.
\]
Then
\[
y^\alpha H_d(y)
\]
is strictly completely monotone on \((0,\infty)\) if \(\alpha\le 1\), and is not completely monotone if \(\alpha>1\). Hence the sharp cutoff in Szabo Open Problem 1.5 is
\[
\alpha_0=1.
\]

As \(y\downarrow0\), the digamma Laurent expansion gives
\[
\psi(y)=-\frac1y-\gamma+O(y),
\qquad
\psi(y+d)=\psi(d)+O(y).
\]
Therefore
\[
H_d(y)=\psi(y+d)-\psi(y)-\frac{d}{y}
=\frac{1-d}{y}+O(1).
\]
Since \(0<d<1\), the leading coefficient is positive. Hence
\[
y^\alpha H_d(y)=(1-d)y^{\alpha-1}+O(y^\alpha).
\]
If \(\alpha>1\), then
\[
\frac{d}{dy}\left(y^\alpha H_d(y)\right)
=(1-d)(\alpha-1)y^{\alpha-2}+O(y^{\alpha-1})>0
\]
for all sufficiently small \(y>0\). A completely monotone function must be nonincreasing, so no \(\alpha>1\) can be completely monotone.

Combining this with the source-imported strict complete monotonicity for \(\alpha\le1\) proves \(\alpha_0=1\).

Szabo's original function is recovered by \(y=x+a\) and \(d=b-a\), where \(a\ge0\), \(b>0\), and \(0<b-a<1\). Translation from \((-a,\infty)\) to \((0,\infty)\) preserves complete-monotonicity sign tests.
\end{proof}

\begin{lemma}[Bondesson Steutel c=1/2 Catalan beta Hausdorff representation]
\label{res:bs-c-half-catalan-hausdorff-representation}
For the Bondesson--Steutel distribution at \(c=1/2\), \(P_n(1/2)=C_n/(2\cdot4^n)=\frac{1}{\pi}\int_0^1 x^{n-1/2}(1-x)^{1/2}\,dx\), where \(C_n\) is the \(n\)-th Catalan number.
\end{lemma}
\begin{proof}
For \(c=1/2\), the source pgf is
\[
P(z)=\frac{1-\sqrt{1-z}}{z}.
\]
Using the Catalan generating function
\[
\sum_{n\ge0}C_n x^n=\frac{1-\sqrt{1-4x}}{2x},
\]
with \(x=z/4\), we get
\[
P(z)=\frac12\sum_{n\ge0}C_n\left(\frac z4\right)^n.
\]
Hence
\[
P_n(1/2)=\frac{C_n}{2\cdot4^n}
=\frac{1}{2\cdot4^n(n+1)}\binom{2n}{n}.
\]

Now
\[
\frac1\pi\int_0^1 x^{n-1/2}(1-x)^{1/2}\,dx
=\frac1\pi B(n+1/2,3/2).
\]
Since \(\Gamma(3/2)=\sqrt\pi/2\) and
\[
\Gamma(n+1/2)=\frac{(2n)!\sqrt\pi}{4^n n!},
\]
this beta integral equals
\[
\frac{(2n)!}{2\cdot4^n n!(n+1)!}
=\frac{C_n}{2\cdot4^n}
=P_n(1/2).
\]
Therefore
\[
P_n(1/2)=\frac1\pi\int_0^1 x^{n-1/2}(1-x)^{1/2}\,dx.
\]
This is a Hausdorff moment representation on \([0,1]\), so \((P_n(1/2))\) is completely monotone as a sequence.

For \(c=1/2\), the canonical sequence is
\[
r_n=(n+1/2)P_n(1/2).
\]
Using the Hausdorff representation,
\[
r_n=\frac{n+1/2}{\pi}\int_0^1 x^{n-1/2}(1-x)^{1/2}\,dx.
\]
Because
\[
\frac{d}{dx}x^{n+1/2}=(n+1/2)x^{n-1/2},
\]
integration by parts gives
\[
r_n
=\frac1\pi\int_0^1 (1-x)^{1/2}\,d(x^{n+1/2})
=\frac{1}{2\pi}\int_0^1 x^{n+1/2}(1-x)^{-1/2}\,dx.
\]
The boundary terms vanish at \(0\) and \(1\). Thus
\[
r_n=\int_0^1 x^n\,d\nu(x),
\qquad
d\nu(x)=\frac{1}{2\pi}x^{1/2}(1-x)^{-1/2}\,dx.
\]
This is a positive finite measure on \([0,1]\). Hence \((r_n)\) is completely monotone.
\end{proof}

\begin{lemma}[Simon moment property gives entire moment generating normal form]
\label{res:simon-leroy-moment-property-normal-form}
For \(\alpha,\beta,\gamma>0\), if Simon's moment property \(M_{\alpha,\beta,\gamma}\) holds for a positive random variable \(X\), then \((\Gamma(\beta))^\gamma F^{(\gamma)}_{\alpha,\beta}(z)={\bf E}[e^{zX}]\) for every \(z\in\mathbb C\).
\end{lemma}
\begin{proof}
\emph{Setup.}
Let
\[
F^{(\gamma)}_{\alpha,\beta}(z)
=\sum_{n\ge0}\frac{z^n}{\Gamma(\beta+\alpha n)^\gamma},
\qquad \alpha,\beta,\gamma>0.
\]

Assume Simon's moment property \(M_{\alpha,\beta,\gamma}\): there is a positive random variable \(X\) with
\[
{\bf E}[X^s]
=\Gamma(1+s)\left(\frac{\Gamma(\beta)}{\Gamma(\beta+\alpha s)}\right)^\gamma,
\qquad s>0.
\]

For \(n\ge0\), this gives
\[
{\bf E}[X^n]
=n!\left(\frac{\Gamma(\beta)}{\Gamma(\beta+\alpha n)}\right)^\gamma.
\]

Set
\[
a_n=\left(\frac{\Gamma(\beta)}{\Gamma(\beta+\alpha n)}\right)^\gamma.
\]
By the standard Gamma-ratio asymptotic,
\[
\frac{a_{n+1}}{a_n}
=\left(\frac{\Gamma(\beta+\alpha n)}{\Gamma(\beta+\alpha(n+1))}\right)^\gamma
\sim (\alpha n)^{-\alpha\gamma}\to0.
\]
Hence \(\sum_{n\ge0}a_n z^n\) has infinite radius of convergence.

For \(t\ge0\), monotone convergence gives
\[
{\bf E}[e^{tX}]
=\sum_{n\ge0}\frac{t^n{\bf E}[X^n]}{n!}
=\sum_{n\ge0}a_n t^n
=(\Gamma(\beta))^\gamma F^{(\gamma)}_{\alpha,\beta}(t).
\]
The right-hand side is finite for every \(t\ge0\). Therefore \(e^{|z|X}\) is integrable for every \(z\in\mathbb C\), and dominated convergence gives
\[
{\bf E}[e^{zX}]
=\sum_{n\ge0}\frac{z^n{\bf E}[X^n]}{n!}
=(\Gamma(\beta))^\gamma F^{(\gamma)}_{\alpha,\beta}(z),
\qquad z\in\mathbb C.
\]

This proves the claim.

For \(x>0\), substituting \(z=-x\) yields
\[
F^{(\gamma)}_{\alpha,\beta}(-x)
=\Gamma(\beta)^{-\gamma}{\bf E}[e^{-xX}].
\]
This is a positive Laplace transform. More explicitly, for every \(k\ge0\),
\[
(-1)^k\frac{d^k}{dx^k}F^{(\gamma)}_{\alpha,\beta}(-x)
=\Gamma(\beta)^{-\gamma}{\bf E}[X^k e^{-xX}]\ge0,
\]
where differentiating under the expectation is justified on compact subintervals of \((0,\infty)\) by the finite moment \({\bf E}[X^k]\).

This proves the claim.
\end{proof}

\begin{lemma}[Alzer Berg reciprocal Gini Gamma quotient complete monotonicity forces a plus b at most...]
\label{res:ab-reciprocal-sum-gate-a-plus-b-le-one-third}
If \(x\mapsto1/P_{a,b}(u,v;x)\) is completely monotone on \((0,\infty)\) for every \(v>u>0\), then \(a+b\le1/3\).
\end{lemma}
\begin{proof}
For \(a\ne b\), the Gini mean is
\[
G_{a,b}(u,v)
=\left(\frac{u^a+v^a}{u^b+v^b}\right)^{1/(a-b)}.
\]
The case \(a=b\) is obtained by continuity. Put \(v=u+h=u(1+t)\). Then
\[
\log \frac{G_{a,b}(u,u+h)}{u}
=\frac{1}{a-b}
\left[
\log(1+(1+t)^a)-\log(1+(1+t)^b)
\right].
\]
For a real parameter \(c\),
\[
\log(1+(1+t)^c)
=\log2+\frac c2t+\frac{c(c-2)}8t^2+O(t^3).
\]
Thus
\[
\log \frac{G_{a,b}(u,u+h)}{u}
=\frac12t+\frac{a+b-2}{8}t^2+O(t^3).
\]
Exponentiating gives
\[
G_{a,b}(u,u+h)
=u+\frac h2+\frac{a+b-1}{8u}h^2+O(h^3).
\]

Let \(Q_{a,b}=1/P_{a,b}\), \(r=a+b\), and \(y=x+u\). With \(v=u+h\),
\[
\log Q_{a,b}(u,u+h;x)
=\log\Gamma(y+h)-\log\Gamma(y)
-h\,\psi(x+G_{a,b}(u,u+h)).
\]
Using the Gini expansion,
\[
x+G_{a,b}(u,u+h)
=y+\frac h2+\frac{r-1}{8u}h^2+O(h^3).
\]
Taylor expansion gives
\[
\log\Gamma(y+h)-\log\Gamma(y)
=h\psi(y)+\frac{h^2}{2}\psi'(y)+\frac{h^3}{6}\psi''(y)+O(h^4),
\]
and
\[
h\psi\left(y+\frac h2+\frac{r-1}{8u}h^2+O(h^3)\right)
=h\psi(y)+\frac{h^2}{2}\psi'(y)
h^3\left(\frac{r-1}{8u}\psi'(y)+\frac18\psi''(y)\right)+O(h^4).
\]
Therefore
\[
\log Q_{a,b}(u,u+h;x)
=h^3\left(
\frac{\psi''(y)}{24}
-\frac{r-1}{8u}\psi'(y)
\right)+O(h^4).
\]

Assume \(x\mapsto Q_{a,b}(u,v;x)\) is completely monotone for every \(v>u>0\). Then it is decreasing and nonnegative. Standard Gamma and digamma asymptotics give
\[
\lim_{x\to\infty}Q_{a,b}(u,v;x)=1.
\]
Hence \(Q_{a,b}(u,v;x)\ge1\), so \(\log Q_{a,b}(u,v;x)\ge0\).

For fixed \(u,x>0\), let \(h\downarrow0\). The leading coefficient in the expansion must be nonnegative:
\[
\frac{\psi''(y)}{24}
-\frac{r-1}{8u}\psi'(y)
\ge0.
\]
Since \(\psi'(y)>0\),
\[
r\le 1+\frac{u\psi''(y)}{3\psi'(y)}.
\]
Here \(y=x+u\), so for each \(y>0\) and every \(0<u<y\),
\[
r\le 1+\frac{u\psi''(y)}{3\psi'(y)}.
\]
Letting \(u\uparrow y\) gives
\[
r\le 1+\frac{y\psi''(y)}{3\psi'(y)}
\qquad (y>0).
\]
Finally, as \(y\downarrow0\),
\[
\psi'(y)\sim y^{-2},
\qquad
\psi''(y)\sim -2y^{-3}.
\]
Therefore
\[
1+\frac{y\psi''(y)}{3\psi'(y)}\to \frac13,
\]
and hence
\[
a+b=r\le \frac13.
\]
\end{proof}

\begin{lemma}[Bulboaca-Zayed gamma quotient derivative sign normal form]
\label{res:bz-gamma-quotient-derivative-normal-form}
For \(x>-1\), away from removable points, the sign of \(\widetilde F'(x)\) is the sign of \(N(x)=(x+6)(x^2+6)\psi(x+1)(\log(x^2+6)-\log(x+6))-(x^2+12x-6)\log\Gamma(x+1)\).
\end{lemma}
\begin{proof}
Set
\[
P(x)=(x+6)(x^2+6),\qquad Q(x)=x^2+12x-6.
\]
For \(x\ge8\), \(Q(x)>0\). The standard inequalities
\[
\psi(x+1)>\log x
\]
and the Robbins/Stirling upper bound imply
\[
\log\Gamma(x+1)\le x\log x
\qquad (x\ge8).
\]
Indeed the Robbins upper bound gives
\[
\log\Gamma(x+1)
<
\left(x+\frac12\right)\log x-x+\frac12\log(2\pi)+\frac1{12x},
\]
and the remaining excess
\[
\frac12\log x-x+\frac12\log(2\pi)+\frac1{12x}
\]
is already negative at \(x=8\) and decreasing thereafter.

It remains to lower-bound \(D(x)\). Let
\[
u(x)=\frac{x^2+6}{x+6}.
\]
For \(x\ge8\), \(u(x)>1\), so
\[
\log u(x)\ge \frac{2(u(x)-1)}{u(x)+1}.
\]
A direct rational simplification gives
\[
\frac{2(u(x)-1)}{u(x)+1}
-\frac{xQ(x)}{P(x)}
=
\frac{x^2(x^3-3x^2-18x-78)}
{(x^2+x+12)(x+6)(x^2+6)}.
\]
The denominator is positive, and \(x^3-3x^2-18x-78\) is increasing for \(x\ge8\) and has value \(98>0\) at \(x=8\). Hence
\[
D(x)=\log u(x)\ge \frac{xQ(x)}{P(x)}
\qquad (x\ge8).
\]

Combining the three inequalities,
\[
N(x)
=P(x)\psi(x+1)D(x)-Q(x)\log\Gamma(x+1)
>
P(x)(\log x)D(x)-Q(x)x\log x
\ge0.
\]
Therefore \(N(x)>0\), and consequently
\[
\widetilde F'(x)>0
\qquad (x\ge8).
\]
\end{proof}

The next result applies the preceding method to the fixed source problem \textbf{APP-0010}.
\begin{corollary}[APP-0010: Nielsen \(k\)-beta derivative-ratio monotonicity]
\label{res:nielsen-k-beta-derivative-ratio-monotonicity}
For \(k>0\), \(n\ge0\), and \(f_k(x)=x\beta_k(x)\), the ratio \(f_k^{(n+1)}(x)/(f_k^{(n)}(x)f_k^{(n+2)}(x))\) is strictly increasing on \((0,\infty)\) when \(n\) is odd and strictly decreasing on \((0,\infty)\) when \(n\) is even.

This resolves the fixed application label \textbf{APP-0010}: Nielsen \(k\)-beta derivative-ratio monotonicity
\emph{Source reference.} \cite{app-app-0010-source}.
\end{corollary}
\begin{proof}
Let \(\mu\) be a positive Borel measure on \([0,\infty)\) with finite moments after the tilt \(e^{-xt}\). Define
\[
a_j(x)=\int_0^\infty t^j e^{-xt}\,d\mu(t)>0
\]
for the needed indices. For fixed \(n\ge0\), set
\[
B_n(x)=\frac{a_{n+1}(x)}{a_n(x)a_{n+2}(x)}.
\]
Then \(B_n\) is strictly increasing whenever \(a_{n+1}(x)>0\).

Indeed \(a_j'(x)=-a_{j+1}(x)\). Put
\[
r_j(x)=\frac{a_{j+1}(x)}{a_j(x)}.
\]
Moment log-convexity, equivalently Cauchy--Schwarz applied to
\[
\int t^j e^{-xt}\,d\mu(t),
\]
gives \(r_{j+1}(x)\ge r_j(x)\). Differentiating,
\[
\frac{d}{dx}\log B_n(x)
=-\frac{a_{n+2}}{a_{n+1}}
+\frac{a_{n+1}}{a_n}
+\frac{a_{n+3}}{a_{n+2}}
=r_n(x)-r_{n+1}(x)+r_{n+2}(x).
\]
Since \(r_{n+2}(x)\ge r_{n+1}(x)\) and \(r_n(x)>0\), the logarithmic derivative is positive. Hence \(B_n\) is strictly increasing.

Let
\[
f_k(x)=x\beta_k(x).
\]
From the integral representation of \(\beta_k\), integration by parts gives
\[
f_k(x)
=x\int_0^\infty e^{-xt}\frac{dt}{1+e^{-kt}}
=\frac12+\int_0^\infty e^{-xt}\frac{k e^{-kt}}{(1+e^{-kt})^2}\,dt.
\]
Thus \(f_k\) is the Laplace transform of a positive measure consisting of an atom \(1/2\) at \(0\) plus the positive density
\[
w_k(t)=\frac{k e^{-kt}}{(1+e^{-kt})^2}
\]
on \((0,\infty)\). All tilted moments are finite, and \(a_{j}(x)>0\) for \(j\ge1\).

For \(a_j(x)=(-1)^j f_k^{(j)}(x)\), the bridge lemma gives that
\[
B_{k,n}(x)=\frac{a_{n+1}(x)}{a_n(x)a_{n+2}(x)}
\]
is strictly increasing for every \(k>0\), \(n\ge0\), and \(x>0\). Since
\[
R_{k,n}(x)=(-1)^{n+1}B_{k,n}(x),
\]
we get the precise monotonicity law:

if \(n\) is odd, \(R_{k,n}\) is strictly increasing on \((0,\infty)\);
if \(n\) is even, \(R_{k,n}\) is strictly decreasing on \((0,\infty)\).

This solves the Yin--Zhang Nielsen \(k\)-beta derivative-ratio open problem in the parity-refined form.
\end{proof}

\begin{proposition}[Baricz V\_q completely monotone in parameter q plus one]
\label{res:baricz-vq-q-complete-monotonicity}
For every fixed \(x>0\), the function \(a\mapsto V_{a-1}(x)\) is completely monotone on \((0,\infty)\).
\end{proposition}
\begin{proof}
Put \(u=t^2-x^2\). Then \(t=(u+x^2)^{1/2}\) and
\[
dt=\frac{du}{2(u+x^2)^{1/2}}.
\]
Since \(e^{x^2}e^{-t^2}=e^{-u}\), the factor \(2\) in the definition cancels the \(1/2\) from \(dt\), giving
\[
V_q(x)=\frac1{\Gamma(q+1)}
\int_0^\infty e^{-u}u^q(u+x^2)^{-1/2}\,du.
\]
This integral is finite for \(q>-1\): near \(0\) the integrand is \(O(u^q)\), and at infinity it is exponentially decaying.

For \(r>0\),
\[
r^{-1/2}=\frac1{\sqrt\pi}\int_0^\infty s^{-1/2}e^{-rs}\,ds.
\]
With \(r=u+x^2\), Tonelli's theorem applies because the integrand is nonnegative. Thus
\[
\begin{aligned}
V_q(x)
&=\frac1{\Gamma(q+1)\sqrt\pi}
\int_0^\infty\int_0^\infty
e^{-u}u^q s^{-1/2}e^{-(u+x^2)s}\,ds\,du\\
&=\frac1{\Gamma(q+1)\sqrt\pi}
\int_0^\infty s^{-1/2}e^{-x^2s}
\left(\int_0^\infty u^q e^{-(1+s)u}\,du\right)\,ds\\
&=\frac1{\sqrt\pi}\int_0^\infty
s^{-1/2}e^{-x^2s}(1+s)^{-(q+1)}\,ds.
\end{aligned}
\]
The last step uses
\[
\int_0^\infty u^q e^{-(1+s)u}\,du
=\frac{\Gamma(q+1)}{(1+s)^{q+1}}.
\]

This proves the claim.

Let \(a=q+1>0\). The normal form becomes
\[
V_{a-1}(x)=\frac1{\sqrt\pi}\int_0^\infty
s^{-1/2}e^{-x^2s}(1+s)^{-a}\,ds.
\]
For each integer \(n\ge0\), differentiation under the integral gives
\[
(-1)^n\frac{d^n}{da^n}V_{a-1}(x)
=\frac1{\sqrt\pi}\int_0^\infty
(\log(1+s))^n s^{-1/2}e^{-x^2s}(1+s)^{-a}\,ds\ge0.
\]
The differentiation is justified locally uniformly in \(a>0\): near \(s=0\), \((\log(1+s))^n s^{-1/2}=O(s^{n-1/2})\), and as \(s\to\infty\), the factor \(e^{-x^2s}\) dominates every logarithmic and polynomial factor.

Thus \(a\mapsto V_{a-1}(x)\) is completely monotone on \((0,\infty)\). Equivalently, after the change of variables \(y=\log(1+s)\), it is a positive Laplace transform in \(a\).

This proves the claim.

For \(a>0\), write
\[
F_x(a)=V_{a-1}(x)
=\int_{(0,\infty)} e^{-ay}\,d\nu_x(y),
\]
where \(\nu_x\) is the positive pushforward of
\[
\frac1{\sqrt\pi}s^{-1/2}e^{-x^2s}\,ds
\]
under \(y=\log(1+s)\). This measure is not concentrated at one point because it has positive density on \(s>0\), hence on \(y>0\).

For \(a_1\ne a_2\) and \(0<\lambda<1\), strict Holder inequality gives
\[
F_x(\lambda a_1+(1-\lambda)a_2)
<
F_x(a_1)^\lambda F_x(a_2)^{1-\lambda}.
\]
Strictness holds because \(e^{-a_1y}\) and \(e^{-a_2y}\) are not proportional \(\nu_x\)-almost everywhere unless \(\nu_x\) is supported at a single \(y\).

Thus \(a\mapsto F_x(a)\) is strictly log-convex on \((0,\infty)\). Since \(a=q+1\), the shift preserves strict log-convexity, so \(q\mapsto V_q(x)\) is strictly log-convex on \((-1,\infty)\) for every fixed \(x>0\).

This proves the claim.
\end{proof}

The next result applies the preceding method to the fixed source problem \textbf{APP-0022}.
\begin{corollary}[APP-0022: Chen-Choi Gurland gamma-ratio conjecture is logarithmically completely monotone]
\label{res:chen-choi-gurland-logf-strict-cm}
For \(F(x)=\Gamma(1/x)\Gamma(3/x)/\Gamma(2/x)^2\), the function \(x\mapsto\log F(x)\) is strictly completely monotone on \((0,\infty)\).

This resolves the fixed application label \textbf{APP-0022}: Prove logarithmic complete monotonicity of the Gurland gamma ratio conjectured in the source.
\emph{Source reference.} \cite{app-app-0022-source}.
\end{corollary}
\begin{proof}
Set
\[
  A(u)=2\log(u+2)-\log(u+1)-\log(u+3).
\]
The identity
\[
 A(u)=\int_0^\infty e^{-ut}{e^{-t}(1-e^{-t})^2\over t}\,dt
\]
shows that \(A\) is strictly completely monotone.  The Weierstrass product for
\(\Gamma\) gives
\[
  \log F(x)=\log {4\over3}+\sum_{m\ge1}A(mx).
\]
Termwise differentiation is justified on compact subintervals of
\((0,\infty)\), and each differentiated summand has the required strict
alternating sign.  This proves logarithmic complete monotonicity.
\end{proof}

\begin{proposition}[Sibisi Prabhakar Q measure is stable subordination of transform-normalized Pollard meas...]
\label{res:prabhakar-q-stable-subordination-normal-form}
For \(0<\alpha<1\), \(\gamma>0\), \(\beta>\alpha\gamma\), and \(\lambda>0\), let \(P^\gamma_{\alpha,\beta}\) be Sibisi's transform-normalized Pollard measure satisfying \(E^\gamma_{\alpha,\beta}(-u)=\int_0^\infty e^{-ur}\,dP^\gamma_{\alpha,\beta}(r)\). Then the measure \(Q^\gamma_{\alpha,\beta}(A\mid\lambda)=\int\Pr((\lambda r)^{1/\alpha}S_\alpha\in A)\,dP^\gamma_{\alpha,\beta}(r)\) satisfies \(E^\gamma_{\alpha,\beta}(-\lambda x^\alpha)=\int_0^\infty e^{-xt}\,dQ^\gamma_{\alpha,\beta}(t\mid\lambda)\).
\end{proposition}
\begin{proof}
Let \(P\) be a finite positive measure on \([0,\infty)\), let
\[
F(u)=\int_0^\infty e^{-ur}\,dP(r),
\]
let \(0<\alpha<1\), \(\lambda>0\), and let \(S_\alpha\) be positive \(\alpha\)-stable:
\[
\mathbb E e^{-xS_\alpha}=e^{-x^\alpha}.
\]
Define a finite positive measure \(Q\) by
\[
Q(A)=\int_{[0,\infty)}
\Pr\{(\lambda r)^{1/\alpha}S_\alpha\in A\}\,dP(r).
\]
Then Tonelli's theorem gives, for \(x>0\),
\[
\int_0^\infty e^{-xt}\,dQ(t)
=\int_{[0,\infty)}
\mathbb E e^{-x(\lambda r)^{1/\alpha}S_\alpha}\,dP(r)
=\int_{[0,\infty)}e^{-\lambda r x^\alpha}\,dP(r)
=F(\lambda x^\alpha).
\]

Apply the lemma to Sibisi's transform-normalized \(P=P^\gamma_{\alpha,\beta}\) and
\[
F(u)=E^\gamma_{\alpha,\beta}(-u).
\]
For \(0<\alpha<1\), \(\gamma>0\), \(\beta>\alpha\gamma\), and \(\lambda>0\), set
\[
Q^\gamma_{\alpha,\beta}(A\mid\lambda)
=\int_{[0,\infty)}
\Pr\{(\lambda r)^{1/\alpha}S_\alpha\in A\}\,
dP^\gamma_{\alpha,\beta}(r).
\]
Then
\[
\int_0^\infty e^{-xt}\,dQ^\gamma_{\alpha,\beta}(t\mid\lambda)
=E^\gamma_{\alpha,\beta}(-\lambda x^\alpha).
\]
By uniqueness of finite Laplace transforms, this is Sibisi's \(Q^\gamma_{\alpha,\beta}(\cdot\mid\lambda)\).

If \(S_\alpha\) has density \(f_\alpha\), then the non-atomic density part is
\[
q^\gamma_{\alpha,\beta}(t\mid\lambda)
=\int_{(0,\infty)}
(\lambda r)^{-1/\alpha}
f_\alpha\!\left(\frac{t}{(\lambda r)^{1/\alpha}}\right)
dP^\gamma_{\alpha,\beta}(r),
\]
with an atom \(P^\gamma_{\alpha,\beta}(\{0\})\delta_0\) if \(P^\gamma_{\alpha,\beta}\) has an atom at zero. When \(P^\gamma_{\alpha,\beta}\) is represented by Sibisi's Pollard density \(p^\gamma_{\alpha,\beta}(r)\), this becomes
\[
q^\gamma_{\alpha,\beta}(t\mid\lambda)
=\int_0^\infty
(\lambda r)^{-1/\alpha}
f_\alpha\!\left(\frac{t}{(\lambda r)^{1/\alpha}}\right)
p^\gamma_{\alpha,\beta}(r)\,dr.
\]

Because
\[
E^\gamma_{\alpha,\beta}(0)=\frac1{\Gamma(\beta)},
\]
the representing measures \(P^\gamma_{\alpha,\beta}\) and \(Q^\gamma_{\alpha,\beta}(\cdot\mid\lambda)\) have total mass \(1/\Gamma(\beta)\). They are finite positive measures, not generally probability measures. The normalized probability version is
\[
\widehat Q^\gamma_{\alpha,\beta}=\Gamma(\beta)Q^\gamma_{\alpha,\beta},
\qquad
\mathcal L\widehat Q^\gamma_{\alpha,\beta}(x)
=\Gamma(\beta)E^\gamma_{\alpha,\beta}(-\lambda x^\alpha).
\]

The boundary \(\beta=\alpha\gamma\) is not included here because Sibisi's ordinary convolution density uses \(1/\Gamma(\beta-\alpha\gamma)\). A limiting treatment is a separate open frontier.
\end{proof}

The next result applies the preceding method to the fixed source problem \textbf{APP-0012}.
\begin{corollary}[APP-0012: Baricz \(V\_q\) strict \(q\)-log-convexity]
\label{res:baricz-vq-strict-q-logconvexity}
For every fixed \(x>0\), the function \(q\mapsto V_q(x)\) is strictly log-convex on \((-1,\infty)\).

This resolves the fixed application label \textbf{APP-0012}: Baricz \(V\_q\) strict \(q\)-log-convexity
\emph{Source reference.} \cite{app-app-0012-source}.
\end{corollary}
\begin{proof}
Put \(u=t^2-x^2\).  Then
\begin{equation}
\label{eq:baricz-u-normal-form}
V_q(x)=\frac1{\Gamma(q+1)}
\int_0^\infty e^{-u}u^q(u+x^2)^{-1/2}\,du.
\end{equation}
Using
\[
r^{-1/2}=\frac1{\sqrt\pi}\int_0^\infty s^{-1/2}e^{-rs}\,ds
\]
and Tonelli's theorem, we obtain
\begin{equation}
\label{eq:baricz-laplace-normal-form}
V_q(x)=\frac1{\sqrt\pi}\int_0^\infty
s^{-1/2}e^{-x^2s}(1+s)^{-(q+1)}\,ds.
\end{equation}
Let \(a=q+1>0\).  Differentiation under the integral gives
\begin{equation}
\label{eq:baricz-complete-monotonicity}
(-1)^n\frac{d^n}{da^n}V_{a-1}(x)
=\frac1{\sqrt\pi}\int_0^\infty
(\log(1+s))^n s^{-1/2}e^{-x^2s}(1+s)^{-a}\,ds\ge0.
\end{equation}
Equivalently, after the change of variables \(y=\log(1+s)\), \(V_{a-1}(x)\) is a positive Laplace transform in \(a\).  The representing measure is not concentrated at a single point, because the density in \eqref{eq:baricz-laplace-normal-form} is positive on \(s>0\).  Strict Holder inequality for two different exponents therefore gives strict log-convexity in \(a\), and the shift \(a=q+1\) gives the stated strict log-convexity in \(q\).
\end{proof}

\begin{lemma}[Baricz V\_q parameter q log-convexity open problem]
\label{res:baricz-vq-q-logconvexity-open-problem}
For Baricz's function \(V_q(x)=\frac{2e^{x^2}}{\Gamma(q+1)}\int_x^\infty e^{-t^2}(t^2-x^2)^q\,dt\), prove that for every fixed \(x>0\), the map \(q\mapsto V_q(x)\) is log-convex on \((-1,\infty)\).
\end{lemma}
\begin{proof}
Put \(u=t^2-x^2\). Then \(t=(u+x^2)^{1/2}\) and
\[
dt=\frac{du}{2(u+x^2)^{1/2}}.
\]
Since \(e^{x^2}e^{-t^2}=e^{-u}\), the factor \(2\) in the definition cancels the \(1/2\) from \(dt\), giving
\[
V_q(x)=\frac1{\Gamma(q+1)}
\int_0^\infty e^{-u}u^q(u+x^2)^{-1/2}\,du.
\]
This integral is finite for \(q>-1\): near \(0\) the integrand is \(O(u^q)\), and at infinity it is exponentially decaying.

For \(r>0\),
\[
r^{-1/2}=\frac1{\sqrt\pi}\int_0^\infty s^{-1/2}e^{-rs}\,ds.
\]
With \(r=u+x^2\), Tonelli's theorem applies because the integrand is nonnegative. Thus
\[
\begin{aligned}
V_q(x)
&=\frac1{\Gamma(q+1)\sqrt\pi}
\int_0^\infty\int_0^\infty
e^{-u}u^q s^{-1/2}e^{-(u+x^2)s}\,ds\,du\\
&=\frac1{\Gamma(q+1)\sqrt\pi}
\int_0^\infty s^{-1/2}e^{-x^2s}
\left(\int_0^\infty u^q e^{-(1+s)u}\,du\right)\,ds\\
&=\frac1{\sqrt\pi}\int_0^\infty
s^{-1/2}e^{-x^2s}(1+s)^{-(q+1)}\,ds.
\end{aligned}
\]
The last step uses
\[
\int_0^\infty u^q e^{-(1+s)u}\,du
=\frac{\Gamma(q+1)}{(1+s)^{q+1}}.
\]

This proves the claim.

Let \(a=q+1>0\). The normal form becomes
\[
V_{a-1}(x)=\frac1{\sqrt\pi}\int_0^\infty
s^{-1/2}e^{-x^2s}(1+s)^{-a}\,ds.
\]
For each integer \(n\ge0\), differentiation under the integral gives
\[
(-1)^n\frac{d^n}{da^n}V_{a-1}(x)
=\frac1{\sqrt\pi}\int_0^\infty
(\log(1+s))^n s^{-1/2}e^{-x^2s}(1+s)^{-a}\,ds\ge0.
\]
The differentiation is justified locally uniformly in \(a>0\): near \(s=0\), \((\log(1+s))^n s^{-1/2}=O(s^{n-1/2})\), and as \(s\to\infty\), the factor \(e^{-x^2s}\) dominates every logarithmic and polynomial factor.

Thus \(a\mapsto V_{a-1}(x)\) is completely monotone on \((0,\infty)\). Equivalently, after the change of variables \(y=\log(1+s)\), it is a positive Laplace transform in \(a\).

This proves the claim.

For \(a>0\), write
\[
F_x(a)=V_{a-1}(x)
=\int_{(0,\infty)} e^{-ay}\,d\nu_x(y),
\]
where \(\nu_x\) is the positive pushforward of
\[
\frac1{\sqrt\pi}s^{-1/2}e^{-x^2s}\,ds
\]
under \(y=\log(1+s)\). This measure is not concentrated at one point because it has positive density on \(s>0\), hence on \(y>0\).

For \(a_1\ne a_2\) and \(0<\lambda<1\), strict Holder inequality gives
\[
F_x(\lambda a_1+(1-\lambda)a_2)
<
F_x(a_1)^\lambda F_x(a_2)^{1-\lambda}.
\]
Strictness holds because \(e^{-a_1y}\) and \(e^{-a_2y}\) are not proportional \(\nu_x\)-almost everywhere unless \(\nu_x\) is supported at a single \(y\).

Thus \(a\mapsto F_x(a)\) is strictly log-convex on \((0,\infty)\). Since \(a=q+1\), the shift preserves strict log-convexity, so \(q\mapsto V_q(x)\) is strictly log-convex on \((-1,\infty)\) for every fixed \(x>0\).

This proves the claim.
\end{proof}

\begin{theorem}[Simon gamma quotient derivative is CM with positive Laplace density]
\label{res:simon-gamma-quotient-derivative-positive-laplace-density}
For every \(0<\alpha<1\), the derivative of \(F_\alpha(x)=\Gamma(x+\alpha)/(\Gamma(x)x^\alpha)\) is completely monotone, with a positive Laplace-density representation obtained from the decreasing kernel \(J_\alpha\).
\end{theorem}
\begin{proof}
For \(0<\alpha<1\), put
\[
F_\alpha(x)=\frac{\Gamma(x+\alpha)}{\Gamma(x)x^\alpha},\qquad x>0.
\]
Then \(F_\alpha\) is a Bernstein function. In fact, \(1-F_\alpha\) is completely monotone.

The stronger route asserting that \(F_\alpha\) is a complete Bernstein function is false.

Let \(\beta=1-\alpha\in(0,1)\). Since \(\Gamma(x+1)=x\Gamma(x)\), the beta integral gives
\[
F_\alpha(x)
=x^\beta\frac{\Gamma(x+\alpha)}{\Gamma(x+1)}
=\frac{x^\beta}{\Gamma(\beta)}
\int_0^\infty e^{-xu}(e^u-1)^{-\alpha}\,du.
\]
Also
\[
1=\frac{x^\beta}{\Gamma(\beta)}
\int_0^\infty e^{-xu}u^{-\alpha}\,du.
\]
Thus
\[
1-F_\alpha(x)
=\frac{x^\beta}{\Gamma(\beta)}
\int_0^\infty e^{-xu}\left(u^{-\alpha}-(e^u-1)^{-\alpha}\right)\,du.
\]

Define
\[
J_\alpha(t)=\int_0^t (t-u)^{\alpha-1}(e^u-1)^{-\alpha}\,du.
\]
By Laplace convolution,
\[
\mathcal L[J_\alpha](x)
=\Gamma(\alpha)x^{-\alpha}
\int_0^\infty e^{-xu}(e^u-1)^{-\alpha}\,du.
\]
Therefore
\[
F_\alpha(x)
=\frac{x}{\Gamma(\alpha)\Gamma(1-\alpha)}\mathcal L[J_\alpha](x).
\]

After the substitution \(u=tv\),
\[
J_\alpha(t)=
\int_0^1
(1-v)^{\alpha-1}
\left(\frac{t}{e^{tv}-1}\right)^\alpha\,dv.
\]
For fixed \(v\in(0,1]\), the function
\[
t\mapsto \frac{t}{e^{tv}-1}
\]
is decreasing because
\[
\frac{d}{dt}\frac{t}{e^{tv}-1}
=
\frac{e^{tv}-1-tv e^{tv}}{(e^{tv}-1)^2}\le0.
\]
The integrand is dominated by \((1-v)^{\alpha-1}v^{-\alpha}\), which is integrable on \((0,1)\). Hence \(J_\alpha\) is decreasing. Moreover,
\[
J_\alpha(0+)
=\int_0^1(1-v)^{\alpha-1}v^{-\alpha}\,dv
=\Gamma(\alpha)\Gamma(1-\alpha),
\]
and \(J_\alpha(t)\to0\) as \(t\to\infty\).

Let \(\mu_\alpha=-dJ_\alpha\). This is a positive finite measure with total mass
\[
\mu_\alpha([0,\infty))=\Gamma(\alpha)\Gamma(1-\alpha).
\]
Stieltjes integration by parts gives, for \(x>0\),
\[
\int_0^\infty e^{-xt}\,\mu_\alpha(dt)
=
\Gamma(\alpha)\Gamma(1-\alpha)-x\mathcal L[J_\alpha](x).
\]
Using the formula for \(F_\alpha\),
\[
1-F_\alpha(x)
=
\frac{1}{\Gamma(\alpha)\Gamma(1-\alpha)}
\int_0^\infty e^{-xt}\,\mu_\alpha(dt).
\]
Thus \(1-F_\alpha\) is completely monotone.

Consequently,
\[
F_\alpha'(x)
=
\frac{1}{\Gamma(\alpha)\Gamma(1-\alpha)}
\int_0^\infty t e^{-xt}\,\mu_\alpha(dt),
\]
so \(F_\alpha'\) is completely monotone. Since \(F_\alpha(x)>0\), \(F_\alpha\) is a Bernstein function.

For \(t>0\), differentiating the \(v\)-integral under the same endpoint domination gives \(\mu_\alpha(dt)=m_\alpha(t)\,dt\), where
\[
m_\alpha(t)
=
\alpha\int_0^1
(1-v)^{\alpha-1}
\left(\frac{t}{e^{tv}-1}\right)^\alpha
\left(
\frac{v e^{tv}}{e^{tv}-1}-\frac1t
\right)\,dv.
\]
The bracket is positive, because with \(a=tv>0\),
\[
\frac{a e^a}{e^a-1}>1.
\]
Therefore the Laplace-density route in the argument is also closed.

The case \(\alpha=1/2\) is the specialization
\[
F_{1/2}(x)=\frac{\Gamma(x+1/2)}{\Gamma(x)\sqrt{x}},
\]
and the same complement-CM representation proves that \(F_{1/2}\) is Bernstein.

For \(z=-r+i0\) with \(r\in(\alpha,1)\), take the upper-lip branch of \(z^{-\alpha}\). Then
\[
F_\alpha(-r+i0)
=
\frac{\Gamma(\alpha-r)}{\Gamma(-r)}r^{-\alpha}e^{-i\pi\alpha}.
\]
Here \(\Gamma(\alpha-r)<0\) and \(\Gamma(-r)<0\), so the real gamma ratio is positive. Hence
\[
\operatorname{Im} F_\alpha(-r+i0)<0.
\]
A complete Bernstein function has the Pick property and maps the upper half-plane into itself. Therefore \(F_\alpha\) is not a complete Bernstein function for \(0<\alpha<1\).
\end{proof}

\begin{lemma}[Simon gamma quotient is Bernstein for alpha in (0,1)]
\label{res:simon-gamma-quotient-bf-alpha-window-open-problem}
For every \(0<\alpha<1\), the function \(F_\alpha(x)=\Gamma(x+\alpha)/(\Gamma(x)x^\alpha)\) is a Bernstein function on \((0,\infty)\).
\end{lemma}
\begin{proof}
Source artifacts:

Gomilko--Tomilov source: arXiv:2210.03474, "To the theory of generalized Stieltjes transforms"
Simon source: Thomas Simon, "Moment problems related to Bernstein functions", Annales de la Faculte des sciences de Toulouse 29 (2020), 577--594, DOI 10.5802/afst.1640

For \(\alpha,\beta>0\), \(s,t\ge0\), and \(x>0\), the beta identity gives
\[
(x+s)^{-\alpha}(x+t)^{-\beta}
=
\frac{\Gamma(\alpha+\beta)}{\Gamma(\alpha)\Gamma(\beta)}
\int_0^1
\frac{u^{\alpha-1}(1-u)^{\beta-1}}
{\bigl(x+us+(1-u)t\bigr)^{\alpha+\beta}}
\,du.
\]

If
\[
f(x)=\int_{[0,\infty)}(x+s)^{-\alpha}\,d\mu(s),
\qquad
g(x)=\int_{[0,\infty)}(x+t)^{-\beta}\,d\nu(t),
\]
then Tonelli gives
\[
f(x)g(x)=\int_{[0,\infty)}(x+r)^{-(\alpha+\beta)}\,d\kappa(r),
\]
where \(\kappa\) is the positive pushforward measure defined by
\[
\int \varphi(r)\,d\kappa(r)
=
\frac{\Gamma(\alpha+\beta)}{\Gamma(\alpha)\Gamma(\beta)}
\int\!\!\int\!\!\int_0^1
\varphi\bigl(us+(1-u)t\bigr)
u^{\alpha-1}(1-u)^{\beta-1}
\,du\,d\mu(s)\,d\nu(t).
\]

Thus pure generalized Stieltjes transforms satisfy
\[
\mathcal S_\alpha^{0}\mathcal S_\beta^{0}\subseteq \mathcal S_{\alpha+\beta}^{0}.
\]
With the usual nonnegative constant term convention, constants and lower-order terms require the standard order-lift convention already recorded in the source vocabulary; no application should rely on constants being absorbed unless that convention is explicitly cited.

For \(0<\alpha<1\), set
\[
F_\alpha(x)=\frac{\Gamma(x+\alpha)}{\Gamma(x)x^\alpha},
\qquad
L_\alpha(x)=\log F_\alpha(x).
\]
Then
\[
L_\alpha'(x)
=
\psi(x+\alpha)-\psi(x)-\frac{\alpha}{x}
=
\int_0^\infty e^{-xt}
\left(
\frac{1-e^{-\alpha t}}{1-e^{-t}}-\alpha
\right)\,dt.
\]
The kernel is positive for \(t>0\), because \(u\mapsto 1-e^{-ut}\) is strictly concave on \([0,1]\), hence
\[
1-e^{-\alpha t}>\alpha(1-e^{-t}).
\]
Therefore \(L_\alpha'\) is completely monotone, and \(1/F_\alpha\) is logarithmically completely monotone.

This does not prove that \(F_\alpha\) is a Bernstein function. Simon's source explicitly states that the Bernstein character of the corresponding \(\Phi(\lambda)=\Gamma(1-t+\lambda)/(\lambda^{1-t}\Gamma(\lambda))\) is a puzzling question, and also says the reciprocal LCM property is necessary but not sufficient.
\end{proof}

\section{Laplace, Stieltjes, and Bernstein transport}
A recurring mechanism is permanence: positive kernels remain positive after mixtures, Stieltjes transforms, Bernstein integration, stable subordination, and complete-monotone closure operations.  This section collects the reusable transport tools and their direct consequences.
\begin{theorem}[complete monotone functions closed under product positive sums and positive mixtures]
\label{res:cm-closure-product-positive-mixture}
Completely monotone functions on \((0,\infty)\) are closed under pointwise products, nonnegative finite sums, and nonnegative parameter integrals whenever the integral is finite.
\end{theorem}
\begin{proof}
The density
\[
\phi_0(t)=\frac{1}{t(\pi^2+\log^2 t)}
\]
is completely monotone on \((0,\infty)\). Consequently \(I_R\) is a Stieltjes function. Moreover, the antiderivative
\[
\widetilde I_R(x)=a+\int_0^\infty (1-e^{-xt})\phi_0(t)\,dt
\]
is a complete Bernstein function.

For \(u=\log t\),
\[
\int_0^1 e^{-au}\sin(\pi a)\,da
=\frac{\pi(1+e^{-u})}{u^2+\pi^2}.
\]
Substituting \(u=\log t\) gives
\[
\int_0^1 t^{-a}\sin(\pi a)\,da
=\frac{\pi(1+t^{-1})}{\pi^2+\log^2 t}.
\]
Therefore
\[
\phi_0(t)
=\frac{1}{t(\pi^2+\log^2 t)}
=\frac{1}{\pi(1+t)}\int_0^1 t^{-a}\sin(\pi a)\,da.
\]

For \(0<a<1\), \(t^{-a}\) is completely monotone because
\[
t^{-a}=\frac1{\Gamma(a)}\int_0^\infty e^{-ts}s^{a-1}\,ds.
\]
Also \(t\mapsto(1+t)^{-1}\) is completely monotone. Complete monotonicity is closed under products and positive mixtures, and \(\sin(\pi a)>0\) on \((0,1)\). Hence \(\phi_0\) is completely monotone.

If \(\phi\) is completely monotone and
\[
F(x)=\int_0^\infty e^{-xt}\phi(t)\,dt
\]
is finite for \(x>0\), then \(F\) is Stieltjes: write \(\phi(t)=\int_0^\infty e^{-ts}\,d\mu(s)\) and use Tonelli to obtain
\[
F(x)=\int_0^\infty\frac{d\mu(s)}{x+s}.
\]
Applying this to \(\phi_0\) proves that \(I_R\) is Stieltjes.

Finally, the source already proves that
\[
\widetilde I_R(x)=a+\int_0^\infty (1-e^{-xt})\phi_0(t)\,dt
\]
is a Bernstein function. The Levy density \(\phi_0\) is completely monotone, and it satisfies the Bernstein integrability condition
\[
\int_0^\infty (1\wedge t)\phi_0(t)\,dt<\infty.
\]
The standard complete-Bernstein criterion for Levy densities therefore gives that \(\widetilde I_R\) is a complete Bernstein function.
\end{proof}

\begin{proposition}[Prabhakar Q density is stable density mixture over Pollard density]
\label{res:prabhakar-q-density-mixture-formula}
If the transform-normalized Pollard measure \(P^\gamma_{\alpha,\beta}\) has density \(p^\gamma_{\alpha,\beta}\) and \(S_\alpha\) has density \(f_\alpha\), then the non-atomic part of Sibisi's \(Q^\gamma_{\alpha,\beta}(\cdot\mid\lambda)\) has density \(q^\gamma_{\alpha,\beta}(t\mid\lambda)=\int_0^\infty (\lambda r)^{-1/\alpha}f_\alpha(t/(\lambda r)^{1/\alpha})p^\gamma_{\alpha,\beta}(r)\,dr\), with a matching pushed-forward atom at zero if \(P^\gamma_{\alpha,\beta}\) has one.
\end{proposition}
\begin{proof}
Let \(P\) be a finite positive measure on \([0,\infty)\), let
\[
F(u)=\int_0^\infty e^{-ur}\,dP(r),
\]
let \(0<\alpha<1\), \(\lambda>0\), and let \(S_\alpha\) be positive \(\alpha\)-stable:
\[
\mathbb E e^{-xS_\alpha}=e^{-x^\alpha}.
\]
Define a finite positive measure \(Q\) by
\[
Q(A)=\int_{[0,\infty)}
\Pr\{(\lambda r)^{1/\alpha}S_\alpha\in A\}\,dP(r).
\]
Then Tonelli's theorem gives, for \(x>0\),
\[
\int_0^\infty e^{-xt}\,dQ(t)
=\int_{[0,\infty)}
\mathbb E e^{-x(\lambda r)^{1/\alpha}S_\alpha}\,dP(r)
=\int_{[0,\infty)}e^{-\lambda r x^\alpha}\,dP(r)
=F(\lambda x^\alpha).
\]

Apply the lemma to Sibisi's transform-normalized \(P=P^\gamma_{\alpha,\beta}\) and
\[
F(u)=E^\gamma_{\alpha,\beta}(-u).
\]
For \(0<\alpha<1\), \(\gamma>0\), \(\beta>\alpha\gamma\), and \(\lambda>0\), set
\[
Q^\gamma_{\alpha,\beta}(A\mid\lambda)
=\int_{[0,\infty)}
\Pr\{(\lambda r)^{1/\alpha}S_\alpha\in A\}\,
dP^\gamma_{\alpha,\beta}(r).
\]
Then
\[
\int_0^\infty e^{-xt}\,dQ^\gamma_{\alpha,\beta}(t\mid\lambda)
=E^\gamma_{\alpha,\beta}(-\lambda x^\alpha).
\]
By uniqueness of finite Laplace transforms, this is Sibisi's \(Q^\gamma_{\alpha,\beta}(\cdot\mid\lambda)\).

If \(S_\alpha\) has density \(f_\alpha\), then the non-atomic density part is
\[
q^\gamma_{\alpha,\beta}(t\mid\lambda)
=\int_{(0,\infty)}
(\lambda r)^{-1/\alpha}
f_\alpha\!\left(\frac{t}{(\lambda r)^{1/\alpha}}\right)
dP^\gamma_{\alpha,\beta}(r),
\]
with an atom \(P^\gamma_{\alpha,\beta}(\{0\})\delta_0\) if \(P^\gamma_{\alpha,\beta}\) has an atom at zero. When \(P^\gamma_{\alpha,\beta}\) is represented by Sibisi's Pollard density \(p^\gamma_{\alpha,\beta}(r)\), this becomes
\[
q^\gamma_{\alpha,\beta}(t\mid\lambda)
=\int_0^\infty
(\lambda r)^{-1/\alpha}
f_\alpha\!\left(\frac{t}{(\lambda r)^{1/\alpha}}\right)
p^\gamma_{\alpha,\beta}(r)\,dr.
\]

Because
\[
E^\gamma_{\alpha,\beta}(0)=\frac1{\Gamma(\beta)},
\]
the representing measures \(P^\gamma_{\alpha,\beta}\) and \(Q^\gamma_{\alpha,\beta}(\cdot\mid\lambda)\) have total mass \(1/\Gamma(\beta)\). They are finite positive measures, not generally probability measures. The normalized probability version is
\[
\widehat Q^\gamma_{\alpha,\beta}=\Gamma(\beta)Q^\gamma_{\alpha,\beta},
\qquad
\mathcal L\widehat Q^\gamma_{\alpha,\beta}(x)
=\Gamma(\beta)E^\gamma_{\alpha,\beta}(-\lambda x^\alpha).
\]

The boundary \(\beta=\alpha\gamma\) is not included here because Sibisi's ordinary convolution density uses \(1/\Gamma(\beta-\alpha\gamma)\). A limiting treatment is a separate open frontier.
\end{proof}

\begin{lemma}[Karp Sitnik Lemma 1 gives Stieltjes normal form for sigma one generalized hypergeometri...]
\label{res:ks-lemma1-stieltjes-normal-form}
Under the Karp--Sitnik positive-density hypotheses \(b_i>a_i>0\), the function \(F(x)={}_{q+1}F_q(1,a_1,\ldots,a_q;b_1,\ldots,b_q;-x)\) is a normalized Stieltjes function on \((0,\infty)\).
\end{lemma}
\begin{proof}
Let \(q\ge1\) and \(b_i>a_i>0\). Define
\[
F(x)={}_{q+1}F_q(1,a_1,\ldots,a_q;b_1,\ldots,b_q;-x)
\]
and
\[
R(x)=
\frac{{}_{q+1}F_q(1,a_1+1,\ldots,a_q+1;b_1+1,\ldots,b_q+1;-x)}
{{}_{q+1}F_q(1,a_1,\ldots,a_q;b_1,\ldots,b_q;-x)}.
\]
Then \(R\) is a Stieltjes function on \((0,\infty)\), hence completely monotone. Moreover
\[
xR(x)
\]
is a complete Bernstein function up to the positive scalar \(\prod_i b_i/a_i\).

For \(q=1\), this gives the Gauss/Beta slice:
\[
\frac{{}_2F_1(1,a+1;b+1;-x)}{{}_2F_1(1,a;b;-x)}
\]
is Stieltjes and completely monotone for \(b>a>0\).

Karp--Sitnik's positive-density representation gives a compact-support Stieltjes form for \(F\):
\[
F(x)=\int_0^1 \frac{d\mu(s)}{1+sx},
\qquad \mu\ge0,
\qquad \mu([0,1])=F(0)=1.
\]
Thus \(F\) is a nonzero Stieltjes function.

By the standard Stieltjes/complete-Bernstein duality, if \(F\) is nonzero Stieltjes, then
\[
G(x)=\frac1{F(x)}
\]
is a complete Bernstein function. Since \(G(0)=1\), the complete Bernstein representation gives
\[
G(x)=1+\alpha x+\int_0^\infty \frac{x}{x+t}\,d\nu(t),
\qquad \alpha\ge0,\quad \nu\ge0.
\]
Therefore
\[
\frac{G(x)-1}{x}
=\alpha+\int_0^\infty \frac{1}{x+t}\,d\nu(t)
\]
is Stieltjes.

Karp--Sitnik's contiguous defect identity for \(\sigma=1,\delta=1\) is
\[
1-F(x)
=x\,{}_{q+1}F_q(1,a_1+1,\ldots,a_q+1;b_1+1,\ldots,b_q+1;-x)
\prod_{i=1}^q\frac{a_i}{b_i}.
\]
Hence
\[
R(x)=
\prod_{i=1}^q\frac{b_i}{a_i}\,
\frac{1-F(x)}{xF(x)}
=
\prod_{i=1}^q\frac{b_i}{a_i}\,
\frac{G(x)-1}{x}.
\]
A positive scalar multiple of a Stieltjes function is Stieltjes. Thus \(R\) is Stieltjes and completely monotone.

Finally,
\[
xR(x)=\prod_{i=1}^q\frac{b_i}{a_i}\,(G(x)-1).
\]
Since subtracting the finite value \(G(0)=1\) from a complete Bernstein function leaves the same nonnegative drift and Levy measure part, \(G(x)-1\) is complete Bernstein. Hence \(xR(x)\) is complete Bernstein up to the stated positive scalar.

This is a strong local theorem and a reusable Stieltjes/CBF bridge. The subsequent primary-source wording audit did not find an explicit Karp--Sitnik open problem asking for the CM/Stieltjes/CBF upgrade, so this result is bridge-only unless a separate source asks for exactly this strengthening.
\end{proof}

\begin{lemma}[reciprocal arctan logarithmic derivative has source imported complete monotonicity]
\label{res:jovanovic-treml-arctan-logderivative-cm}
The function \(x\mapsto ((1+x^2)\arctan x)^{-1}\) is completely monotone on \((0,\infty)\), by the Jovanovic--Treml positive-kernel theorem.
\end{lemma}
\begin{proof}
If \(1/\arctan x\) were a Stieltjes function, then its reciprocal \(\arctan x\) would be a Bernstein function. A Bernstein function has completely monotone derivative.

But
\[
(\arctan x)'=\frac1{1+x^2},
\]
and
\[
\left(\frac1{1+x^2}\right)''
=\frac{2(3x^2-1)}{(1+x^2)^3},
\]
which is negative for \(0<x<1/\sqrt3\). Hence \((\arctan x)'\) is not completely monotone, \(\arctan x\) is not Bernstein, and \(1/\arctan x\) is not Stieltjes.
\end{proof}

\begin{lemma}[heat flow L2 energy is completely monotone in heat time]
\label{res:heat-flow-l2-energy-completely-monotone}
For every probability measure \(\mu\) on \(\mathbb R^d\), the heat-flow \(L^2\) energy \(N_2(t)=\int_{\mathbb R^d}(G_t*\mu)(x)^2\,dx\) is completely monotone on \((0,\infty)\).
\end{lemma}
\begin{proof}
Let \(\mu\) be a probability measure on \(\mathbb R^d\), and let \(G_t\) be the heat kernel with Fourier transform
\[
\widehat G_t(\xi)=e^{-t|\xi|^2}.
\]
For \(p_t=G_t*\mu\), define
\[
N_2(t)=\int_{\mathbb R^d}p_t(x)^2\,dx,\qquad t>0.
\]
Then \(N_2\) is completely monotone on \((0,\infty)\). If \(S_2(t)=1-N_2(t)\), then \(S_2'\) is completely monotone on \((0,\infty)\). Hence, for each \(t_0>0\), the normalized increment \(S_2(t)-S_2(t_0)\) is Bernstein as a function of \(t\ge t_0\).

Since \(|\widehat\mu(\xi)|\le1\), the function \(e^{-t|\xi|^2}\widehat\mu(\xi)\) is in \(L^2(\mathbb R^d)\) for every \(t>0\). Thus \(p_t\in L^2\), and Plancherel gives
\[
N_2(t)
=(2\pi)^{-d}\int_{\mathbb R^d}e^{-2t|\xi|^2}|\widehat\mu(\xi)|^2\,d\xi.
\]
Define \(\nu_\mu\) to be the locally finite positive pushforward of
\[
(2\pi)^{-d}|\widehat\mu(\xi)|^2\,d\xi
\]
under the map \(\xi\mapsto|\xi|^2\). Then
\[
N_2(t)=\int_0^\infty e^{-2tr}\,d\nu_\mu(r).
\]
The integral is finite for every \(t>0\), because \(|\widehat\mu|\le1\) and \(\int_{\mathbb R^d}e^{-2t|\xi|^2}\,d\xi<\infty\).

For \(k\ge0\), differentiation under the integral is justified at \(t>0\) by the Gaussian factor. It yields
\[
(-1)^kN_2^{(k)}(t)
=\int_0^\infty (2r)^k e^{-2tr}\,d\nu_\mu(r)\ge0.
\]
Therefore \(N_2\) is completely monotone.

Now \(S_2'(t)=-N_2'(t)\), and
\[
S_2'(t)=\int_0^\infty 2r e^{-2tr}\,d\nu_\mu(r),
\]
another positive Laplace transform. Hence \(S_2'\) is completely monotone. This implies that every normalized increment \(S_2(t)-S_2(t_0)\), \(t\ge t_0>0\), is a Bernstein function on that restricted half-line.
\end{proof}

\begin{lemma}[reciprocal of nonzero Bernstein function is completely monotone]
\label{res:bf-reciprocal-is-cm}
If \(g\) is a nonzero Bernstein function on \((0,\infty)\), then \(1/g\) is completely monotone.
\end{lemma}
\begin{proof}
For
\[
f_{\alpha,\beta}(x)=x^{-\alpha}(1+x^\beta)^{-1},\qquad x>0,
\]
the following hold.

1. If \(\alpha\ge0\) and \(0\le\beta\le1\), then \(f_{\alpha,\beta}\) is completely monotone.
2. If \(1<\beta\le2\) and \(\alpha\ge\beta/2\), then \(f_{\alpha,\beta}\) is completely monotone.
3. If \(\beta=2\), then \(f_{\alpha,2}\) is completely monotone if and only if \(\alpha\ge1\).

Consequently, for the Dagum threshold
\[
c(\beta)=\inf\{\alpha\ge0:f_{\alpha,\beta}\in CM\},
\]
one has \(c(1)=0\), \(c(2)=1\), and \(c(\beta)\le\beta/2\) for \(1<\beta\le2\).

For \(0<\beta\le1\), \(x^\beta\) is a Bernstein function, so \(1+x^\beta\) is Bernstein. A nonzero Bernstein function has completely monotone reciprocal. Hence \((1+x^\beta)^{-1}\) is completely monotone. The factor \(x^{-\alpha}\) is completely monotone for \(\alpha\ge0\), and products of completely monotone functions are completely monotone. The case \(\beta=0\) is immediate because \(f_{\alpha,0}(x)=2^{-1}x^{-\alpha}\).

For \(1<\beta\le2\), put \(\gamma=\beta/2\in(1/2,1]\). The elementary Laplace identity
\[
\frac1{x(1+x^2)}=\int_0^\infty e^{-xt}(1-\cos t)\,dt
\]
shows that \(h(x)=1/(x(1+x^2))\) is completely monotone. Since \(x^\gamma\) is Bernstein, \(h(x^\gamma)\) is completely monotone. But
\[
h(x^\gamma)=\frac1{x^\gamma(1+x^{2\gamma})}
=\frac1{x^{\beta/2}(1+x^\beta)}
=f_{\beta/2,\beta}(x).
\]
Multiplication by \(x^{-(\alpha-\beta/2)}\) proves complete monotonicity for \(\alpha\ge\beta/2\).

For \(\beta=2\), sufficiency for \(\alpha\ge1\) follows from the preceding paragraph. For necessity, first \(f_{0,2}(x)=1/(1+x^2)\) is not completely monotone, for example because its inverse Laplace kernel is \(\sin t\).

Let \(0<\alpha<1\). The inverse Laplace kernel of \(x^{-\alpha}\) is \(t^{\alpha-1}/\Gamma(\alpha)\), and the inverse Laplace kernel of \(1/(1+x^2)\) is \(\sin t\). Hence the inverse kernel of \(f_{\alpha,2}\) is
\[
\eta_\alpha(t)=\frac1{\Gamma(\alpha)}\int_0^t(t-s)^{\alpha-1}\sin s\,ds.
\]
At \(t=2\pi\),
\[
\Gamma(\alpha)\eta_\alpha(2\pi)
=\int_0^\pi\left((2\pi-s)^{\alpha-1}-(\pi-s)^{\alpha-1}\right)\sin s\,ds.
\]
Since \(\alpha-1<0\), the bracket is negative for \(0<s<\pi\). Therefore \(\eta_\alpha(2\pi)<0\). By uniqueness of the Laplace transform for locally finite measures, \(f_{\alpha,2}\) cannot be completely monotone. Thus \(f_{\alpha,2}\in CM\) exactly when \(\alpha\ge1\).
\end{proof}

\begin{lemma}[Baskakov alpha one r six positive inverse Laplace density factorization]
\label{res:baskakov-alpha1-r6-positive-density-factorization}
For \(m=3\), equivalently \(r=6\), the Baskakov \(\alpha=1\) function \(f^{[6]}_1(x)=1/((1+x)^6-x^6)\) has inverse-Laplace density \(\rho_3(t)=\frac{4}{3}e^{-t/2}(1-\cos(t/(2\sqrt3)))^2(2+\cos(t/(2\sqrt3)))\), hence is completely monotone.
\end{lemma}
\begin{proof}
Set
\[
D_m(x)=(1+x)^{2m}-x^{2m}.
\]
The roots are obtained from
\[
\frac{1+x}{x}=\omega_k,\qquad
\omega_k=e^{\pi i k/m},\qquad k=1,\ldots,2m-1.
\]
Thus
\[
\lambda_{k,m}
=\frac{1}{\omega_k-1}
=-\frac12-\frac{i}{2}\cot\frac{\pi k}{2m}.
\]
At this root,
\[
D_m'(\lambda_{k,m})
=2m\lambda_{k,m}^{2m-1}(\omega_k^{-1}-1),
\]
and direct simplification gives the real residue
\[
R_{k,m}
=\frac{1}{D_m'(\lambda_{k,m})}
=\frac{(-1)^{m+k}}{2m}
\left(2\sin\frac{\pi k}{2m}\right)^{2m-2}.
\]
Consequently the inverse-Laplace density is
\[
\rho_m(t)
=\sum_{k=1}^{2m-1}R_{k,m}e^{\lambda_{k,m}t}.
\]
Pairing conjugate roots gives
\[
\rho_m(t)
=\frac{e^{-t/2}}{2m}
\left[
2^{2m-2}
+2\sum_{k=1}^{m-1}
(-1)^{m+k}
\left(2\sin\frac{\pi k}{2m}\right)^{2m-2}
\cos\!\left(
\frac{t}{2}\cot\frac{\pi k}{2m}
\right)
\right].
\]

This proves the finite trigonometric density formula used by the diagnostic
candidate.

For \(m=2\), this gives
\[
\rho_2(t)=e^{-t/2}\left(1-\cos\frac t2\right)
=2e^{-t/2}\sin^2\frac t4\ge0,
\]
recovering the previously admitted \(r=4,\alpha=1\) seed.

For \(m=3\),
\[
\rho_3(t)
=\frac{e^{-t/2}}6
\left[
16+2\cos\left(\frac{\sqrt3\,t}{2}\right)
-18\cos\left(\frac{t}{2\sqrt3}\right)
\right].
\]
Set \(u=t/(2\sqrt3)\).  Since
\(\cos(3u)=4\cos^3u-3\cos u\),
\[
\rho_3(t)
=\frac{4}{3}e^{-t/2}(1-\cos u)^2(2+\cos u)\ge0.
\]
Thus \(r=6,\alpha=1\) is completely monotone by an explicit positive density.

For \(m=4\), write \(h_4(t)=e^{t/2}\rho_4(t)\).  The density formula gives
\[
h_4(t)
=8+2\cos\frac t2
-\frac{(2-\sqrt2)^3}{4}
\cos\left(\frac{1+\sqrt2}{2}t\right)
-\frac{(2+\sqrt2)^3}{4}
\cos\left(\frac{\sqrt2-1}{2}t\right).
\]
At \(t=10\pi\),
\[
\cos(t/2)=\cos(5\pi)=-1,
\]
and both other cosine terms equal \(-\cos(5\pi\sqrt2)\).  Since
\[
(2-\sqrt2)^3+(2+\sqrt2)^3=40,
\]
we obtain
\[
h_4(10\pi)=6+10\cos(5\pi\sqrt2).
\]
Now
\[
0<5\sqrt2-7<\frac14
\]
because \(7/5<\sqrt2<29/20\).  Hence, with
\(\delta=5\sqrt2-7\),
\[
\cos(5\pi\sqrt2)=\cos(7\pi+\pi\delta)=-\cos(\pi\delta)
<-\cos\frac\pi4=-\frac{\sqrt2}{2}<-\frac35.
\]
Therefore
\[
h_4(10\pi)<6+10\left(-\frac35\right)=0,
\]
and
\[
\rho_4(10\pi)=e^{-5\pi}h_4(10\pi)<0.
\]
Since the inverse-Laplace transform is unique for measures on
\([0,\infty)\), \(f^{[8]}_1(x)=1/((1+x)^8-x^8)\) cannot be completely monotone.
\end{proof}

\begin{lemma}[Bazhlekova two term square root derivative polynomial normal form]
\label{res:bazhlekova-odd-derivative-polynomial-normal-form}
For \(h(s)=\sqrt{c s^a+d s^b}\), with \(c,d>0\), \(\Delta=a-b\), \(B=b/2\), and \(y=(c/d)s^\Delta\), the derivatives have the form \(h^{(n)}(s)=\sqrt d\,s^{B-n}(1+y)^{1/2-n}Q_n(y)\), where \(Q_0=1\) and \(Q_{n+1}=(B-n)(1+y)Q_n+\Delta y((1+y)Q_n'+(1/2-n)Q_n)\).
\end{lemma}
\begin{proof}
If for an odd \(n=2q+1\) there is \(y_0>0\) such that
\[
Q_n(y_0)<0,
\]
then for every \(c,d>0\) the point
\[
s_0=\left(\frac{d y_0}{c}\right)^{1/(a-b)}
\]
satisfies \(h^{(n)}(s_0)<0\).

A Bernstein function \(h\) has completely monotone derivative, hence for odd \(n\ge3\) it must satisfy \(h^{(n)}\ge0\). Thus the displayed sign proves that \(h\) is not Bernstein.

More directly for the Bazhlekova transform,
\[
F_x(s)=e^{-x h(s)}=\mathcal L\{w_t(x,\cdot)\}(s),
\]
we have
\[
F_x^{(n)}(s_0)=-x h^{(n)}(s_0)+O(x^2).
\]
For odd \(n\), complete monotonicity would require \(F_x^{(n)}(s_0)\le0\). If \(h^{(n)}(s_0)<0\), then \(F_x^{(n)}(s_0)>0\) for all sufficiently small \(x>0\), so \(w_t(x,\cdot)\) cannot be nonnegative.

The previous exact inner-gap example is recovered by the recurrence. For
\[
(a,b)=\left(\frac{28}{25},\frac1{50}\right),
\]
one has
\[
a-b=\frac{11}{10},\qquad
(a-1)^2+(b-1)^2=\frac{2437}{2500}<1,
\]
and
\[
Q_5(1)=-\frac{5570045943}{10000000000}<0.
\]
This matches the earlier value
\[
h^{(5)}(1)
=-\frac{5570045943\sqrt2}{320000000000}<0
\]
when \(c=d=1\).

The seventh derivative adds a new exact residual seed. For
\[
(a,b)=\left(\frac{107}{100},\frac1{100}\right),
\]
the point lies in the residual inner gap and exact Sturm counting gives no positive-root obstruction for \(Q_5\). But
\[
Q_7\left(\frac32\right)
=-\frac{56394696198326721417}{327680000000000000}<0.
\]

The ninth derivative adds another seed. For
\[
(a,b)=\left(\frac{53}{50},\frac1{100}\right),
\]
exact Sturm counting gives no positive-root obstruction for \(Q_5\) or \(Q_7\), but
\[
Q_9\left(\frac74\right)
=-\frac{3262653919060921278409379787}
{262144000000000000000000}<0.
\]

The finite \(5,7,9\) diagnostic does not cover the residual inner gap. At
\[
(a,b)=\left(\frac32,\frac25\right)
\]
we have
\[
0<b<a-1,\qquad a-b=\frac{11}{10}>1,\qquad
(a-1)^2+(b-1)^2=\frac{61}{100}<1.
\]
Exact Sturm counts show that \(Q_5,Q_7,Q_9\) have no positive roots. Since each has positive value at \(0\), each is positive on \(y>0\). Thus the fifth, seventh, and ninth derivative tests cannot decide this exponent pair.

The same noncoverage phenomenon also occurs at the near-boundary residual point
\[
(a,b)=\left(\frac{11}{10},\frac1{20}\right),
\qquad
(a-1)^2+(b-1)^2=\frac{73}{80}<1.
\]
Again \(Q_5,Q_7,Q_9\) have no positive roots and are positive at \(0\).

The derivative-polynomial normal form for \(h^{(n)}\).
The odd-derivative small-\(x\) obstruction criterion.
The exact fifth/seventh/ninth negative seeds above.
The exact finite \(5,7,9\) no-cover seeds above.
\end{proof}

The next result applies the preceding method to the fixed source problem \textbf{APP-0014}.
\begin{theorem}[APP-0014: Ramanujan integral is Stieltjes and has a complete Bernstein primitive]
\label{res:ramanujan-density-logsquare-cm}
The function \(\phi_0(t)=1/[t(\pi^2+\log^2 t)]\) is completely monotone on \((0,\infty)\).

This resolves the fixed application label \textbf{APP-0014}: Classify the Ramanujan integral by a standard positive-kernel transform rather than only by inequalities.
\emph{Source reference.} \cite{app-app-0014-source}.
\end{theorem}
\begin{proof}
For \(t>0\),
\[
 {1\over t(\pi^2+\log^2t)}
   ={1\over \pi(1+t)}\int_0^1 t^{-u}\sin(\pi u)\,du .
\]
The right side is a positive superposition of completely monotone powers and is
therefore completely monotone as a function of \(t\).  The Laplace transform of
a completely monotone density is Stieltjes.  Integrating the same density
against \(1-e^{-xt}\) gives a complete Bernstein function by the standard
Stieltjes-density characterization of complete Bernstein functions.
\end{proof}

\begin{lemma}[Karp Sitnik sigma one delta one shifted generalized hypergeometric ratio is Stieltjes a...]
\label{res:ks-general-q-delta1-sigma1-ratio-stieltjes-cm}
For \(q\ge1\) and \(b_i>a_i>0\), the shifted ratio \({}_{q+1}F_q(1,a_1+1,\ldots,a_q+1;b_1+1,\ldots,b_q+1;-x)/{}_{q+1}F_q(1,a_1,\ldots,a_q;b_1,\ldots,b_q;-x)\) is a Stieltjes function, hence completely monotone, on \((0,\infty)\).
\end{lemma}
\begin{proof}
Let \(q\ge1\) and \(b_i>a_i>0\). Define
\[
F(x)={}_{q+1}F_q(1,a_1,\ldots,a_q;b_1,\ldots,b_q;-x)
\]
and
\[
R(x)=
\frac{{}_{q+1}F_q(1,a_1+1,\ldots,a_q+1;b_1+1,\ldots,b_q+1;-x)}
{{}_{q+1}F_q(1,a_1,\ldots,a_q;b_1,\ldots,b_q;-x)}.
\]
Then \(R\) is a Stieltjes function on \((0,\infty)\), hence completely monotone. Moreover
\[
xR(x)
\]
is a complete Bernstein function up to the positive scalar \(\prod_i b_i/a_i\).

For \(q=1\), this gives the Gauss/Beta slice:
\[
\frac{{}_2F_1(1,a+1;b+1;-x)}{{}_2F_1(1,a;b;-x)}
\]
is Stieltjes and completely monotone for \(b>a>0\).

Karp--Sitnik's positive-density representation gives a compact-support Stieltjes form for \(F\):
\[
F(x)=\int_0^1 \frac{d\mu(s)}{1+sx},
\qquad \mu\ge0,
\qquad \mu([0,1])=F(0)=1.
\]
Thus \(F\) is a nonzero Stieltjes function.

By the standard Stieltjes/complete-Bernstein duality, if \(F\) is nonzero Stieltjes, then
\[
G(x)=\frac1{F(x)}
\]
is a complete Bernstein function. Since \(G(0)=1\), the complete Bernstein representation gives
\[
G(x)=1+\alpha x+\int_0^\infty \frac{x}{x+t}\,d\nu(t),
\qquad \alpha\ge0,\quad \nu\ge0.
\]
Therefore
\[
\frac{G(x)-1}{x}
=\alpha+\int_0^\infty \frac{1}{x+t}\,d\nu(t)
\]
is Stieltjes.

Karp--Sitnik's contiguous defect identity for \(\sigma=1,\delta=1\) is
\[
1-F(x)
=x\,{}_{q+1}F_q(1,a_1+1,\ldots,a_q+1;b_1+1,\ldots,b_q+1;-x)
\prod_{i=1}^q\frac{a_i}{b_i}.
\]
Hence
\[
R(x)=
\prod_{i=1}^q\frac{b_i}{a_i}\,
\frac{1-F(x)}{xF(x)}
=
\prod_{i=1}^q\frac{b_i}{a_i}\,
\frac{G(x)-1}{x}.
\]
A positive scalar multiple of a Stieltjes function is Stieltjes. Thus \(R\) is Stieltjes and completely monotone.

Finally,
\[
xR(x)=\prod_{i=1}^q\frac{b_i}{a_i}\,(G(x)-1).
\]
Since subtracting the finite value \(G(0)=1\) from a complete Bernstein function leaves the same nonnegative drift and Levy measure part, \(G(x)-1\) is complete Bernstein. Hence \(xR(x)\) is complete Bernstein up to the stated positive scalar.

This is a strong local theorem and a reusable Stieltjes/CBF bridge. The subsequent primary-source wording audit did not find an explicit Karp--Sitnik open problem asking for the CM/Stieltjes/CBF upgrade, so this result is bridge-only unless a separate source asks for exactly this strengthening.
\end{proof}

\begin{lemma}[reciprocal arctan is logarithmically completely monotone]
\label{res:reciprocal-arctan-lcm}
The function \(x\mapsto1/\arctan x\) is logarithmically completely monotone on \((0,\infty)\).
\end{lemma}
\begin{proof}
If \(1/\arctan x\) were a Stieltjes function, then its reciprocal \(\arctan x\) would be a Bernstein function. A Bernstein function has completely monotone derivative.

But
\[
(\arctan x)'=\frac1{1+x^2},
\]
and
\[
\left(\frac1{1+x^2}\right)''
=\frac{2(3x^2-1)}{(1+x^2)^3},
\]
which is negative for \(0<x<1/\sqrt3\). Hence \((\arctan x)'\) is not completely monotone, \(\arctan x\) is not Bernstein, and \(1/\arctan x\) is not Stieltjes.
\end{proof}

\begin{lemma}[Bondesson Steutel c=1/2 canonical r sequence completely monotone]
\label{res:bs-c-half-canonical-r-sequence-cm}
For the Bondesson--Steutel distribution at \(c=1/2\), the canonical sequence \(r_n=(n+1/2)P_n(1/2)\) is completely monotone, with \(r_n=\frac{1}{2\pi}\int_0^1 x^{n+1/2}(1-x)^{-1/2}\,dx\).
\end{lemma}
\begin{proof}
For \(c=1/2\), the source pgf is
\[
P(z)=\frac{1-\sqrt{1-z}}{z}.
\]
Using the Catalan generating function
\[
\sum_{n\ge0}C_n x^n=\frac{1-\sqrt{1-4x}}{2x},
\]
with \(x=z/4\), we get
\[
P(z)=\frac12\sum_{n\ge0}C_n\left(\frac z4\right)^n.
\]
Hence
\[
P_n(1/2)=\frac{C_n}{2\cdot4^n}
=\frac{1}{2\cdot4^n(n+1)}\binom{2n}{n}.
\]

Now
\[
\frac1\pi\int_0^1 x^{n-1/2}(1-x)^{1/2}\,dx
=\frac1\pi B(n+1/2,3/2).
\]
Since \(\Gamma(3/2)=\sqrt\pi/2\) and
\[
\Gamma(n+1/2)=\frac{(2n)!\sqrt\pi}{4^n n!},
\]
this beta integral equals
\[
\frac{(2n)!}{2\cdot4^n n!(n+1)!}
=\frac{C_n}{2\cdot4^n}
=P_n(1/2).
\]
Therefore
\[
P_n(1/2)=\frac1\pi\int_0^1 x^{n-1/2}(1-x)^{1/2}\,dx.
\]
This is a Hausdorff moment representation on \([0,1]\), so \((P_n(1/2))\) is completely monotone as a sequence.

For \(c=1/2\), the canonical sequence is
\[
r_n=(n+1/2)P_n(1/2).
\]
Using the Hausdorff representation,
\[
r_n=\frac{n+1/2}{\pi}\int_0^1 x^{n-1/2}(1-x)^{1/2}\,dx.
\]
Because
\[
\frac{d}{dx}x^{n+1/2}=(n+1/2)x^{n-1/2},
\]
integration by parts gives
\[
r_n
=\frac1\pi\int_0^1 (1-x)^{1/2}\,d(x^{n+1/2})
=\frac{1}{2\pi}\int_0^1 x^{n+1/2}(1-x)^{-1/2}\,dx.
\]
The boundary terms vanish at \(0\) and \(1\). Thus
\[
r_n=\int_0^1 x^n\,d\nu(x),
\qquad
d\nu(x)=\frac{1}{2\pi}x^{1/2}(1-x)^{-1/2}\,dx.
\]
This is a positive finite measure on \([0,1]\). Hence \((r_n)\) is completely monotone.
\end{proof}

\begin{lemma}[Ramanujan integral I\_R is Stieltjes]
\label{res:ramanujan-integral-stieltjes}
The Ramanujan integral \(I_R(x)=\int_0^\infty e^{-xt}\,dt/[t(\pi^2+\log^2 t)]\) is a Stieltjes function on \((0,\infty)\).
\end{lemma}
\begin{proof}
The density
\[
\phi_0(t)=\frac{1}{t(\pi^2+\log^2 t)}
\]
is completely monotone on \((0,\infty)\). Consequently \(I_R\) is a Stieltjes function. Moreover, the antiderivative
\[
\widetilde I_R(x)=a+\int_0^\infty (1-e^{-xt})\phi_0(t)\,dt
\]
is a complete Bernstein function.

For \(u=\log t\),
\[
\int_0^1 e^{-au}\sin(\pi a)\,da
=\frac{\pi(1+e^{-u})}{u^2+\pi^2}.
\]
Substituting \(u=\log t\) gives
\[
\int_0^1 t^{-a}\sin(\pi a)\,da
=\frac{\pi(1+t^{-1})}{\pi^2+\log^2 t}.
\]
Therefore
\[
\phi_0(t)
=\frac{1}{t(\pi^2+\log^2 t)}
=\frac{1}{\pi(1+t)}\int_0^1 t^{-a}\sin(\pi a)\,da.
\]

For \(0<a<1\), \(t^{-a}\) is completely monotone because
\[
t^{-a}=\frac1{\Gamma(a)}\int_0^\infty e^{-ts}s^{a-1}\,ds.
\]
Also \(t\mapsto(1+t)^{-1}\) is completely monotone. Complete monotonicity is closed under products and positive mixtures, and \(\sin(\pi a)>0\) on \((0,1)\). Hence \(\phi_0\) is completely monotone.

If \(\phi\) is completely monotone and
\[
F(x)=\int_0^\infty e^{-xt}\phi(t)\,dt
\]
is finite for \(x>0\), then \(F\) is Stieltjes: write \(\phi(t)=\int_0^\infty e^{-ts}\,d\mu(s)\) and use Tonelli to obtain
\[
F(x)=\int_0^\infty\frac{d\mu(s)}{x+s}.
\]
Applying this to \(\phi_0\) proves that \(I_R\) is Stieltjes.

Finally, the source already proves that
\[
\widetilde I_R(x)=a+\int_0^\infty (1-e^{-xt})\phi_0(t)\,dt
\]
is a Bernstein function. The Levy density \(\phi_0\) is completely monotone, and it satisfies the Bernstein integrability condition
\[
\int_0^\infty (1\wedge t)\phi_0(t)\,dt<\infty.
\]
The standard complete-Bernstein criterion for Levy densities therefore gives that \(\widetilde I_R\) is a complete Bernstein function.
\end{proof}

\section{Determinants and moment-ratio tools}
Moment-ratio and determinant inequalities provide a compressed way to propagate positivity through derivative hierarchies.  These tools are particularly useful for Mills-ratio and Bessel-type estimates where the same normal form controls many levels at once.
\begin{theorem}[Mills ratio positive Laplace representation and complete monotonicity]
\label{res:from-mills-laplace-cm-normal-form}
For \(t>0\), the standard normal Mills ratio satisfies \(r(t)=\int_0^\infty e^{-tu-u^2/2}\,du\). Hence \(r\) is completely monotone on \((0,\infty)\).
\end{theorem}
\begin{proof}
For the standard normal Mills ratio, the normalizing constants cancel:
\[
r(t)=\frac{\int_t^\infty e^{-z^2/2}\,dz}{e^{-t^2/2}}
=e^{t^2/2}\int_t^\infty e^{-z^2/2}\,dz.
\]
Put \(z=t+u\). Then
\[
r(t)=\int_0^\infty e^{-tu-u^2/2}\,du.
\]
For \(t>0\), differentiating under the integral gives
\[
(-1)^n r^{(n)}(t)=\int_0^\infty u^n e^{-tu-u^2/2}\,du>0.
\]
Thus \(r\) is completely monotone on \((0,\infty)\).

From the same integral,
\[
r'(t)=-\int_0^\infty u e^{-tu-u^2/2}\,du.
\]
Since
\[
\frac{d}{du}e^{-tu-u^2/2}=-(t+u)e^{-tu-u^2/2},
\]
integration over \([0,\infty)\) gives
\[
\int_0^\infty (t+u)e^{-tu-u^2/2}\,du=1.
\]
Therefore \(\int_0^\infty u e^{-tu-u^2/2}\,du=1-tr(t)\), and
\[
r'(t)=tr(t)-1.
\]

Set \(P_0=1\), \(Q_0=0\), and suppose
\[
r^{(n)}(t)=P_n(t)r(t)+Q_n(t).
\]
Differentiating and using \(r'=tr-1\),
\[
r^{(n+1)}(t)=(P_n'(t)+tP_n(t))r(t)+(Q_n'(t)-P_n(t)).
\]
Thus
\[
P_{n+1}=P_n'+tP_n,\qquad Q_{n+1}=Q_n'-P_n.
\]
This inductively gives the claimed polynomial recurrence.

For each integer \(L\ge0\), define
\[
f_L(t)=(-1)^L r^{(L)}(t)=\int_0^\infty u^L e^{-tu-u^2/2}\,du.
\]
This is a positive Laplace transform. Let
\[
a_j(t)=(-1)^j f_L^{(j)}(t)=\int_0^\infty u^{L+j}e^{-tu-u^2/2}\,du.
\]
Normalize the positive measure \(u^L e^{-tu-u^2/2}du\) to a probability law \(X\). Up to the positive factor \(a_0(t)^2\), the determinant expression is
\[
\mathbb E[X^4]-4\mathbb E[X^3]\mathbb E[X]+3\mathbb E[X^2]^2.
\]
With \(\mu=\mathbb E[X]\), this equals
\[
\mathbb E[(X-\mu)^4]+3\operatorname{Var}(X)^2\ge0.
\]
Hence
\[
f_L^{(4)}f_L-4f_L^{(3)}f_L'+3(f_L'')^2\ge0.
\]
Since \(f_L^{(j)}=(-1)^L r^{(L+j)}\), this is exactly
\[
r^{(L+4)}r^{(L)}-4r^{(L+3)}r^{(L+1)}+3(r^{(L+2)})^2\ge0.
\]
\end{proof}

\begin{theorem}[Yang Tian Bessel W general zero partial fraction proves W\_nu sqrt is Bernstein]
\label{res:yt-bessel-w-general-zero-partial-fraction}
For every \(\nu>-1\), if \(j_{\nu+1,n}\) are the positive zeros of \(J_{\nu+1}\), then \(G_\nu(s)=W_\nu(\sqrt{s})\) satisfies \(G_\nu(s)=2(\nu+1)+2\sum_{n\ge1}s/(s+j_{\nu+1,n}^2)\) and \(G_\nu'(s)=2\sum_{n\ge1}j_{\nu+1,n}^2/(s+j_{\nu+1,n}^2)^2\). Consequently \(G_\nu\) is a Bernstein function on \((0,\infty)\).
\end{theorem}
\begin{proof}
Put \(\mu=\nu+1>0\). The canonical product for the modified Bessel function gives, locally uniformly in \(z\),
\[
I_\mu(z)=\frac{(z/2)^\mu}{\Gamma(\mu+1)}
\prod_{n=1}^{\infty}\left(1+\frac{z^2}{j_{\mu,n}^2}\right),
\]
where \(j_{\mu,n}\) are the positive zeros of \(J_\mu\). Therefore
\[
z\frac{I_\mu'(z)}{I_\mu(z)}
=\mu+2\sum_{n=1}^{\infty}\frac{z^2}{z^2+j_{\mu,n}^2}.
\]
The recurrence
\[
I_\mu'(z)=I_{\mu-1}(z)-\frac{\mu}{z}I_\mu(z)
\]
then yields
\[
W_\nu(z)
=z\frac{I_{\mu-1}(z)}{I_\mu(z)}
=2\mu+2\sum_{n=1}^{\infty}\frac{z^2}{z^2+j_{\mu,n}^2}.
\]

Set \(s=z^2\). Then
\[
G_\nu(s):=W_\nu(\sqrt s)
=2(\nu+1)+2\sum_{n=1}^{\infty}\frac{s}{s+j_{\nu+1,n}^2}.
\]

Termwise differentiation is justified because \(j_{\nu+1,n}\sim \pi n\), so the differentiated sums converge locally uniformly on \((0,\infty)\). Thus
\[
G_\nu'(s)
=2\sum_{n=1}^{\infty}
\frac{j_{\nu+1,n}^2}{(s+j_{\nu+1,n}^2)^2}.
\]
For every \(k\ge0\),
\[
(-1)^k\frac{d^k}{ds^k}
\frac{j_{\nu+1,n}^2}{(s+j_{\nu+1,n}^2)^2}
=
(k+1)!\frac{j_{\nu+1,n}^2}{(s+j_{\nu+1,n}^2)^{k+2}}
\ge0.
\]
The locally uniformly convergent sum preserves these inequalities, hence \(G_\nu'\) is completely monotone and \(G_\nu\) is a Bernstein function.

If \(0<\tau\le1/2\), then \(x^{2\tau}\) is a Bernstein function. By Bernstein-function composition closure,
\[
x\mapsto W_\nu(x^\tau)=G_\nu(x^{2\tau})
\]
is a Bernstein function on \((0,\infty)\).
\end{proof}

\begin{proposition}[positive Laplace tilted moments imply strict monotonicity of a\_{n+1}/(a\_n a\_{n+2})]
\label{res:cm-laplace-moment-ratio-monotonicity}
Let \(a_j(x)=\int_0^\infty t^j e^{-xt}\,d\mu(t)>0\) be tilted moments of a positive Laplace measure for the needed indices. Then, for every \(n\ge0\), \(B_n(x)=a_{n+1}(x)/(a_n(x)a_{n+2}(x))\) is strictly increasing on its domain.
\end{proposition}
\begin{proof}
Let \(\mu\) be a positive Borel measure on \([0,\infty)\) with finite moments after the tilt \(e^{-xt}\). Define
\[
a_j(x)=\int_0^\infty t^j e^{-xt}\,d\mu(t)>0
\]
for the needed indices. For fixed \(n\ge0\), set
\[
B_n(x)=\frac{a_{n+1}(x)}{a_n(x)a_{n+2}(x)}.
\]
Then \(B_n\) is strictly increasing whenever \(a_{n+1}(x)>0\).

Indeed \(a_j'(x)=-a_{j+1}(x)\). Put
\[
r_j(x)=\frac{a_{j+1}(x)}{a_j(x)}.
\]
Moment log-convexity, equivalently Cauchy--Schwarz applied to
\[
\int t^j e^{-xt}\,d\mu(t),
\]
gives \(r_{j+1}(x)\ge r_j(x)\). Differentiating,
\[
\frac{d}{dx}\log B_n(x)
=-\frac{a_{n+2}}{a_{n+1}}
+\frac{a_{n+1}}{a_n}
+\frac{a_{n+3}}{a_{n+2}}
=r_n(x)-r_{n+1}(x)+r_{n+2}(x).
\]
Since \(r_{n+2}(x)\ge r_{n+1}(x)\) and \(r_n(x)>0\), the logarithmic derivative is positive. Hence \(B_n\) is strictly increasing.

Let
\[
f_k(x)=x\beta_k(x).
\]
From the integral representation of \(\beta_k\), integration by parts gives
\[
f_k(x)
=x\int_0^\infty e^{-xt}\frac{dt}{1+e^{-kt}}
=\frac12+\int_0^\infty e^{-xt}\frac{k e^{-kt}}{(1+e^{-kt})^2}\,dt.
\]
Thus \(f_k\) is the Laplace transform of a positive measure consisting of an atom \(1/2\) at \(0\) plus the positive density
\[
w_k(t)=\frac{k e^{-kt}}{(1+e^{-kt})^2}
\]
on \((0,\infty)\). All tilted moments are finite, and \(a_{j}(x)>0\) for \(j\ge1\).

For \(a_j(x)=(-1)^j f_k^{(j)}(x)\), the bridge lemma gives that
\[
B_{k,n}(x)=\frac{a_{n+1}(x)}{a_n(x)a_{n+2}(x)}
\]
is strictly increasing for every \(k>0\), \(n\ge0\), and \(x>0\). Since
\[
R_{k,n}(x)=(-1)^{n+1}B_{k,n}(x),
\]
we get the precise monotonicity law:

if \(n\) is odd, \(R_{k,n}\) is strictly increasing on \((0,\infty)\);
if \(n\) is even, \(R_{k,n}\) is strictly decreasing on \((0,\infty)\).

This solves the Yin--Zhang Nielsen \(k\)-beta derivative-ratio open problem in the parity-refined form.
\end{proof}

\begin{theorem}[Mills ratio all-L fourth-order determinant inequality]
\label{res:from-mills-derivative-determinant-all-l}
For every integer \(L\ge0\) and every \(t>0\), \(r^{(L+4)}(t)r^{(L)}(t)-4r^{(L+3)}(t)r^{(L+1)}(t)+3(r^{(L+2)}(t))^2\ge0\), where \(r\) is the standard normal Mills ratio.
\end{theorem}
\begin{proof}
For the standard normal Mills ratio, the normalizing constants cancel:
\[
r(t)=\frac{\int_t^\infty e^{-z^2/2}\,dz}{e^{-t^2/2}}
=e^{t^2/2}\int_t^\infty e^{-z^2/2}\,dz.
\]
Put \(z=t+u\). Then
\[
r(t)=\int_0^\infty e^{-tu-u^2/2}\,du.
\]
For \(t>0\), differentiating under the integral gives
\[
(-1)^n r^{(n)}(t)=\int_0^\infty u^n e^{-tu-u^2/2}\,du>0.
\]
Thus \(r\) is completely monotone on \((0,\infty)\).

From the same integral,
\[
r'(t)=-\int_0^\infty u e^{-tu-u^2/2}\,du.
\]
Since
\[
\frac{d}{du}e^{-tu-u^2/2}=-(t+u)e^{-tu-u^2/2},
\]
integration over \([0,\infty)\) gives
\[
\int_0^\infty (t+u)e^{-tu-u^2/2}\,du=1.
\]
Therefore \(\int_0^\infty u e^{-tu-u^2/2}\,du=1-tr(t)\), and
\[
r'(t)=tr(t)-1.
\]

Set \(P_0=1\), \(Q_0=0\), and suppose
\[
r^{(n)}(t)=P_n(t)r(t)+Q_n(t).
\]
Differentiating and using \(r'=tr-1\),
\[
r^{(n+1)}(t)=(P_n'(t)+tP_n(t))r(t)+(Q_n'(t)-P_n(t)).
\]
Thus
\[
P_{n+1}=P_n'+tP_n,\qquad Q_{n+1}=Q_n'-P_n.
\]
This inductively gives the claimed polynomial recurrence.

For each integer \(L\ge0\), define
\[
f_L(t)=(-1)^L r^{(L)}(t)=\int_0^\infty u^L e^{-tu-u^2/2}\,du.
\]
This is a positive Laplace transform. Let
\[
a_j(t)=(-1)^j f_L^{(j)}(t)=\int_0^\infty u^{L+j}e^{-tu-u^2/2}\,du.
\]
Normalize the positive measure \(u^L e^{-tu-u^2/2}du\) to a probability law \(X\). Up to the positive factor \(a_0(t)^2\), the determinant expression is
\[
\mathbb E[X^4]-4\mathbb E[X^3]\mathbb E[X]+3\mathbb E[X^2]^2.
\]
With \(\mu=\mathbb E[X]\), this equals
\[
\mathbb E[(X-\mu)^4]+3\operatorname{Var}(X)^2\ge0.
\]
Hence
\[
f_L^{(4)}f_L-4f_L^{(3)}f_L'+3(f_L'')^2\ge0.
\]
Since \(f_L^{(j)}=(-1)^L r^{(L+j)}\), this is exactly
\[
r^{(L+4)}r^{(L)}-4r^{(L+3)}r^{(L+1)}+3(r^{(L+2)})^2\ge0.
\]
\end{proof}

The next result applies the preceding method to the fixed source problem \textbf{APP-0040}.
\begin{corollary}[APP-0040: Yang-Tian Bessel-\(W\) power-Bernstein conjecture is true]
\label{res:yt-bessel-w-power-bernstein-conjecture}
For \(\tau\in(0,1/2]\) and \(\nu>-1\), prove that \(x\mapsto W_\nu(x^\tau)\), where \(W_\nu(x)=xI_\nu(x)/I_{\nu+1}(x)\), is a Bernstein function on \((0,\infty)\).

This resolves the fixed application label \textbf{APP-0040}: Solve Yang and Tian's Bessel-\(W\) power-Bernstein conjecture on the source interval.
\emph{Source reference.} \cite{app-app-0040-source}.
\end{corollary}
\begin{proof}
Put \(\mu=\nu+1>0\). The canonical product for the modified Bessel function gives, locally uniformly in \(z\),
\[
I_\mu(z)=\frac{(z/2)^\mu}{\Gamma(\mu+1)}
\prod_{n=1}^{\infty}\left(1+\frac{z^2}{j_{\mu,n}^2}\right),
\]
where \(j_{\mu,n}\) are the positive zeros of \(J_\mu\). Therefore
\[
z\frac{I_\mu'(z)}{I_\mu(z)}
=\mu+2\sum_{n=1}^{\infty}\frac{z^2}{z^2+j_{\mu,n}^2}.
\]
The recurrence
\[
I_\mu'(z)=I_{\mu-1}(z)-\frac{\mu}{z}I_\mu(z)
\]
then yields
\[
W_\nu(z)
=z\frac{I_{\mu-1}(z)}{I_\mu(z)}
=2\mu+2\sum_{n=1}^{\infty}\frac{z^2}{z^2+j_{\mu,n}^2}.
\]

Set \(s=z^2\). Then
\[
G_\nu(s):=W_\nu(\sqrt s)
=2(\nu+1)+2\sum_{n=1}^{\infty}\frac{s}{s+j_{\nu+1,n}^2}.
\]

Termwise differentiation is justified because \(j_{\nu+1,n}\sim \pi n\), so the differentiated sums converge locally uniformly on \((0,\infty)\). Thus
\[
G_\nu'(s)
=2\sum_{n=1}^{\infty}
\frac{j_{\nu+1,n}^2}{(s+j_{\nu+1,n}^2)^2}.
\]
For every \(k\ge0\),
\[
(-1)^k\frac{d^k}{ds^k}
\frac{j_{\nu+1,n}^2}{(s+j_{\nu+1,n}^2)^2}
=
(k+1)!\frac{j_{\nu+1,n}^2}{(s+j_{\nu+1,n}^2)^{k+2}}
\ge0.
\]
The locally uniformly convergent sum preserves these inequalities, hence \(G_\nu'\) is completely monotone and \(G_\nu\) is a Bernstein function.

If \(0<\tau\le1/2\), then \(x^{2\tau}\) is a Bernstein function. By Bernstein-function composition closure,
\[
x\mapsto W_\nu(x^\tau)=G_\nu(x^{2\tau})
\]
is a Bernstein function on \((0,\infty)\).
\end{proof}

\begin{lemma}[Yang Tian Bessel W\_nu x power tau Bernstein conjecture for tau <= 1/2 and nu > -1]
\label{res:yang-tian-bessel-w-bernstein-conjecture}
For \(\tau\in(0,1/2]\) and \(\nu>-1\), the function \(x\mapsto W_\nu(x^\tau)\), where \(W_\nu(x)=xI_\nu(x)/I_{\nu+1}(x)\), is a Bernstein function on \((0,\infty)\).
\end{lemma}
\begin{proof}
Put \(\mu=\nu+1>0\). The canonical product for the modified Bessel function gives, locally uniformly in \(z\),
\[
I_\mu(z)=\frac{(z/2)^\mu}{\Gamma(\mu+1)}
\prod_{n=1}^{\infty}\left(1+\frac{z^2}{j_{\mu,n}^2}\right),
\]
where \(j_{\mu,n}\) are the positive zeros of \(J_\mu\). Therefore
\[
z\frac{I_\mu'(z)}{I_\mu(z)}
=\mu+2\sum_{n=1}^{\infty}\frac{z^2}{z^2+j_{\mu,n}^2}.
\]
The recurrence
\[
I_\mu'(z)=I_{\mu-1}(z)-\frac{\mu}{z}I_\mu(z)
\]
then yields
\[
W_\nu(z)
=z\frac{I_{\mu-1}(z)}{I_\mu(z)}
=2\mu+2\sum_{n=1}^{\infty}\frac{z^2}{z^2+j_{\mu,n}^2}.
\]

Set \(s=z^2\). Then
\[
G_\nu(s):=W_\nu(\sqrt s)
=2(\nu+1)+2\sum_{n=1}^{\infty}\frac{s}{s+j_{\nu+1,n}^2}.
\]

Termwise differentiation is justified because \(j_{\nu+1,n}\sim \pi n\), so the differentiated sums converge locally uniformly on \((0,\infty)\). Thus
\[
G_\nu'(s)
=2\sum_{n=1}^{\infty}
\frac{j_{\nu+1,n}^2}{(s+j_{\nu+1,n}^2)^2}.
\]
For every \(k\ge0\),
\[
(-1)^k\frac{d^k}{ds^k}
\frac{j_{\nu+1,n}^2}{(s+j_{\nu+1,n}^2)^2}
=
(k+1)!\frac{j_{\nu+1,n}^2}{(s+j_{\nu+1,n}^2)^{k+2}}
\ge0.
\]
The locally uniformly convergent sum preserves these inequalities, hence \(G_\nu'\) is completely monotone and \(G_\nu\) is a Bernstein function.

If \(0<\tau\le1/2\), then \(x^{2\tau}\) is a Bernstein function. By Bernstein-function composition closure,
\[
x\mapsto W_\nu(x^\tau)=G_\nu(x^{2\tau})
\]
is a Bernstein function on \((0,\infty)\).
\end{proof}

The next result applies the preceding method to the fixed source problem \textbf{APP-0032}.
\begin{corollary}[APP-0032: From's all-\(L\) Mills-ratio bound family is explicit]
\label{res:from-mills-all-l-alternating-r-bound-family}
Let \(P_0=1\), \(P_1=t\), \(P_{n+1}=tP_n+nP_{n-1}\), and \(B_0=0\), \(B_1=1\), \(B_{n+1}=nB_{n-1}-tB_n\). With \(U_L(t)\) as in the all-\(L\) moment-ratio bound, the standard normal Mills ratio satisfies, for even \(L\ge0\), \(r(t)>[B_{L+1}-U_LB_L]/[P_{L+1}+U_LP_L]\), and for odd \(L\ge1\), \(r(t)<[U_LB_L-B_{L+1}]/[P_{L+1}+U_LP_L]\), for all \(t\ge0\).

This resolves the fixed application label \textbf{APP-0032}: Solve From's open all-\(L\) Mills-ratio bound problem from Remark 6.5.
\emph{Source reference.} \cite{app-app-0032-source}.
\end{corollary}
\begin{proof}
Set
\[
m_L(t)=\frac{M_{L+1}(t)}{M_L(t)}.
\]

The recurrence gives
\[
\frac{M_{L+2}}{M_L}=L+1-tm_L,
\]
\[
\frac{M_{L+3}}{M_L}=(t^2+L+2)m_L-t(L+1),
\]
and
\[
\frac{M_{L+4}}{M_L}=(L+1)(t^2+L+3)-t(t^2+2L+5)m_L.
\]

Substitution into the determinant inequality and expansion give
\[
(L+1)(t^2+4L+6)-t(t^2+4L+7)m_L-(t^2+4L+8)m_L^2>0.
\]

Equivalently,
\[
(t^2+4L+8)m_L^2+t(t^2+4L+7)m_L-(L+1)(t^2+4L+6)<0.
\]

Since \(m_L>0\), the positive root yields the uniform strict bound
\[
m_L(t)<U_L(t),
\]
where
\[
U_L(t)=
\frac{
\sqrt{t^2(t^2+4L+7)^2+4(L+1)(t^2+4L+8)(t^2+4L+6)}
-t(t^2+4L+7)
}
{2(t^2+4L+8)}.
\]

This is the determinant-to-bound extractor. It is a single theorem for every \(L\ge0\).

Define polynomials \(P_n(t)\) by
\[
P_0=1,\qquad P_1=t,\qquad P_{n+1}=tP_n+nP_{n-1}.
\]
Define \(B_n(t)\) by
\[
B_0=0,\qquad B_1=1,\qquad B_{n+1}=nB_{n-1}-tB_n.
\]

Then the same recurrence gives
\[
M_n(t)=(-1)^nP_n(t)r(t)+B_n(t).
\]

The ratio bound \(M_{L+1}<U_LM_L\) becomes
\[
\left((-1)^{L+1}P_{L+1}-(-1)^L U_LP_L\right)r
<U_LB_L-B_{L+1}.
\]

Since \(P_n(t)>0\) for \(t>0\) and \(P_{2j}(0)>0\), the coefficient has sign \((-1)^{L+1}\). Thus the bounds alternate:

For even \(L\ge0\),
\[
r(t)>
\frac{B_{L+1}(t)-U_L(t)B_L(t)}
{P_{L+1}(t)+U_L(t)P_L(t)}.
\]

For odd \(L\ge1\),
\[
r(t)<
\frac{U_L(t)B_L(t)-B_{L+1}(t)}
{P_{L+1}(t)+U_L(t)P_L(t)}.
\]

These formulas are valid at \(t=0\) by direct evaluation of the recurrences and for \(t>0\) by the positivity of \(M_L\).
\end{proof}

The next result applies the preceding method to the fixed source problem \textbf{APP-0016}.
\begin{corollary}[APP-0016: From's Mills-ratio bound extends to every derivative level]
\label{res:from-mills-explicit-bound-family-frontier}
After the Mills-ratio Laplace/complete-monotone normal form, derivative recurrence, and all-\(L\) determinant bridge, the remaining From frontier is solved by a closed, explicit alternating upper/lower bound family valid uniformly for all \(L\ge1\).

This resolves the fixed application label \textbf{APP-0016}: Lift the published Mills-ratio bound to the full derivative hierarchy.
\emph{Source reference.} \cite{app-app-0016-source}.
\end{corollary}
\begin{proof}
The density defining \(M_L\) is a positive moment density.  Its order-four
Hankel-minor consequence gives
\[
M_{L+4}M_L-4M_{L+3}M_{L+1}+3M_{L+2}^2>0.
\]
Using the recurrence \(M_{n+1}=nM_{n-1}-tM_n\), this determinant inequality
reduces to a quadratic inequality for \(m_L\); solving the admissible positive
root gives \(m_L<U_L\).  Iterating the same recurrence expresses \(M_L\) in the
form
\[
M_L=(-1)^L(P_L r-B_L),
\]
and substituting the level-\(L\) quotient bound gives the displayed alternating
upper and lower bounds for \(r\).
\end{proof}

\section{Counterexamples and endpoint obstructions}
The theory also records negative information.  Endpoint expansions, rational witnesses, and explicit sign failures identify the boundary of the positive-kernel method and prevent false conjectures from being absorbed as theorems.
The next result applies the preceding method to the fixed source problem \textbf{APP-0042}.
\begin{counterexample}[APP-0042: Modified Bessel square-root log-concavity fails]
\label{res:not-bessel-i-sqrt-log-concavity-nu-ge-0}
not(For every \(\nu\ge0\), the function \(u\mapsto \sqrt{u}\,I_\nu(u)\) is strictly log-concave on \((0,\infty)\).)

This resolves the fixed application label \textbf{APP-0042}: Is \(u\mapsto \sqrt{u}I\_\nu(u)\) strictly log-concave for all \(\nu\ge0\)?
\emph{Source reference.} \cite{app-app-0042-source}.
\end{counterexample}
\begin{proof}
Let
\[
g_\nu(u)=\log(\sqrt u\,I_\nu(u)),
\qquad
r_\nu(u)=\frac{I_\nu'(u)}{I_\nu(u)}.
\]
Then
\[
g_\nu''(u)=r_\nu'(u)-\frac{1}{2u^2}.
\]
The modified Bessel equation
\[
u^2I_\nu''(u)+uI_\nu'(u)-(u^2+\nu^2)I_\nu(u)=0
\]
gives
\[
\frac{I_\nu''(u)}{I_\nu(u)}
=1+\frac{\nu^2}{u^2}-\frac{r_\nu(u)}{u}.
\]
Since
\[
r_\nu'(u)=\frac{I_\nu''(u)}{I_\nu(u)}-r_\nu(u)^2,
\]
we get
\[
g_\nu''(u)
=
1+\frac{\nu^2-\frac12}{u^2}
-\frac{r_\nu(u)}{u}
-r_\nu(u)^2.
\]
This proves the claim.

For \(\nu=0\),
\[
r_0(u)=\frac{I_0'(u)}{I_0(u)}=\frac{I_1(u)}{I_0(u)}.
\]
At \(u=10\), the Riccati expression is
\[
g_0''(10)=\frac{199}{200}-\frac{r_0(10)}{10}-r_0(10)^2.
\]
It is enough to prove
\[
r_0(10)<\frac{9487}{10000},
\]
because
\[
\frac{199}{200}
-\frac{1}{10}\frac{9487}{10000}
-\left(\frac{9487}{10000}\right)^2
=\frac{9831}{100000000}>0.
\]

\[
I_0(10)=\sum_{k=0}^{\infty}\frac{25^k}{(k!)^2},
\qquad
I_1(10)=\sum_{k=0}^{\infty}\frac{5\cdot25^k}{k!(k+1)!}.
\]
With \(N=20\), set
\[
L_0=\sum_{k=0}^{20}\frac{25^k}{(k!)^2}<I_0(10).
\]
For the \(I_1\) tail after \(N=20\), the term ratio is bounded by
\[
q=\frac{25}{22\cdot23}<1,
\]
so the script computes an exact rational upper bound \(U_1>I_1(10)\) and verifies
\[
10000\,U_1<9487\,L_0.
\]
Therefore \(I_1(10)/I_0(10)<9487/10000\), and hence
\[
g_0''(10)>0.
\]

This proves the claim.
\end{proof}

The next result applies the preceding method to the fixed source problem \textbf{APP-0017}.
\begin{counterexample}[APP-0017: Baricz gamma-quotient Bernstein-for-all problem has a rational counterexample]
\label{res:baricz-gamma-quotient-a2b3-not-bf}
For \(a=2\) and \(b=3\), the Baricz Gamma quotient \(x\mapsto\Gamma(x)\Gamma(x-a+b)/(\Gamma(x-a)\Gamma(x+b))\) is not a Bernstein function on \((2,\infty)\).

This resolves the fixed application label \textbf{APP-0017}: Decide whether the gamma quotient is Bernstein throughout the full parameter range.
\emph{Source reference.} \cite{app-app-0017-source}.
\end{counterexample}
\begin{proof}
Take \(a=2\), \(b=3\), and set \(y=x-2\).  The quotient becomes
\[
  g(y)={y(y+1)\over (y+3)(y+4)} .
\]
A direct differentiation gives
\[
  g'''(y)=
 {36(y^4+8y^3+12y^2-40y-94)\over (y+3)^4(y+4)^4}.
\]
Thus \(g'''(1)<0\).  If \(g\) were Bernstein, then \(g'\) would be completely
monotone and hence \(g'''\ge0\).  The displayed value contradicts that
necessary condition.
\end{proof}

The next result applies the preceding method to the fixed source problem \textbf{APP-0041}.
\begin{corollary}[APP-0041: Qi \(h\_\lambda\) degree-four complete-monotonicity conjecture is false]
\label{res:ma-weigert-log-function-dk-chain-conjecture}
In the Ma--Weigert log-function class \(\mathcal F_{1,n}\), the derivative-sign regions \(D_k=\{f:L^k f\ge0\}\), with \(L=-d/dx\), form a descending chain: \(D_{k+1}\subseteq D_k\) for every \(k,n\).

This resolves the fixed application label \textbf{APP-0041}: Settle Qi's degree-4 complete-monotonicity conjecture for \(h\_\lambda\).
\emph{Source reference.} \cite{app-app-0041-source}.
\end{corollary}
\begin{proof}
We first prove that for each \(k\ge0\),
\[
L^k f(x)=x^{-k-1}p_k(\log x)
\]
for some polynomial \(p_k\).

For \(k=0\), this is the definition with \(p_0=p\).  Suppose
\[
L^k f(x)=x^{-k-1}p_k(\log x).
\]
Then
\[
\begin{aligned}
L^{k+1}f(x)
&=-\frac{d}{dx}\left(x^{-k-1}p_k(\log x)\right)\\
&=x^{-k-2}\left((k+1)p_k(\log x)-p_k'(\log x)\right).
\end{aligned}
\]
Thus \(p_{k+1}(y)=(k+1)p_k(y)-p_k'(y)\), which is again a polynomial.

Since every polynomial in \(\log x\) grows slower than every positive power of \(x\),
\[
x^{-k-1}p_k(\log x)\to0
\qquad (x\to\infty).
\]
Therefore
\[
\lim_{x\to\infty}L^k f(x)=0
\]
for every \(k\ge0\).

Assume \(f\in D_{k+1}\).  Then
\[
L^{k+1}f(x)\ge0
\]
on \((0,\infty)\).  Put
\[
g(x)=L^k f(x).
\]
By definition of \(L\),
\[
-g'(x)=L^{k+1}f(x)\ge0.
\]
Using \(g(x)\to0\) as \(x\to\infty\), integrate from \(x\) to \(R\):
\[
g(x)-g(R)=\int_x^R L^{k+1}f(t)\,dt.
\]
Letting \(R\to\infty\) gives
\[
L^k f(x)=g(x)=\int_x^\infty L^{k+1}f(t)\,dt\ge0.
\]
Hence \(f\in D_k\).  Therefore
\[
D_{k+1}\subseteq D_k
\]
for all \(k\ge0\) and all \(n\).
\end{proof}

\begin{lemma}[Girjoaba Rasa c(n,k) full sequence not Hausdorff for any n>=1]
\label{res:gr-cnk-full-sequence-not-hausdorff-any-n}
For every \(n\ge1\), the full Girjoaba--Rasa sequence \(k\mapsto c(n,k)\), \(0\le k\le2n\), is not a Hausdorff moment sequence in the natural order.
\end{lemma}
\begin{proof}
The source sequence is
\[
c(n,k)=\binom{2n}{k}^{-2}
\sum_{j=0}^k \binom{n}{j}^2\binom{n}{k-j}^2,
\]
with binomial coefficients outside their natural range interpreted as \(0\).

For \(n=1\):
\[
c(1,0)=\binom20^{-2}\binom10^2\binom10^2=1.
\]
For \(k=1\),
\[
c(1,1)=\binom21^{-2}
\left(\binom10^2\binom11^2+\binom11^2\binom10^2\right)
=\frac14(1+1)=\frac12.
\]
For \(k=2\),
\[
c(1,2)=\binom22^{-2}
\left(\binom11^2\binom11^2\right)=1.
\]
Thus the full \(n=1\) sequence begins
\[
1,\quad \frac12,\quad 1.
\]

If \((h_k)\) is a Hausdorff moment sequence on \([0,1]\), then
\[
h_k=\int_0^1 t^k\,d\mu(t)
\]
for some positive measure \(\mu\).  Since \(0\le t\le1\), we have \(t^{k+1}\le t^k\), and therefore
\[
h_{k+1}=\int_0^1 t^{k+1}\,d\mu(t)
\le
\int_0^1 t^k\,d\mu(t)=h_k.
\]
Every Hausdorff moment sequence is nonincreasing.

The \(n=1\) Girjoaba--Rasa sequence violates this necessary condition because
\[
c(1,2)=1>\frac12=c(1,1).
\]

The obstruction is actually uniform in \(n\ge1\).  At the right endpoint,
\[
c(n,2n)=1,
\]
because only \(j=n\) contributes and \(\binom{2n}{2n}=1\).  At the adjacent point \(k=2n-1\), only \(j=n-1\) and \(j=n\) contribute, so
\[
\sum_{j=0}^{2n-1}\binom{n}{j}^2\binom{n}{2n-1-j}^2
=2n^2.
\]
Since \(\binom{2n}{2n-1}=2n\), this gives
\[
c(n,2n-1)=\frac{2n^2}{(2n)^2}=\frac12.
\]
Thus \(c(n,2n)>c(n,2n-1)\) for every \(n\ge1\), so no full sequence \(0\le k\le2n\) can be Hausdorff moment in the natural order.
\end{proof}

\begin{counterexample}[Girjoaba Rasa c(n,k) n=1 Hausdorff moment counterexample]
\label{res:gr-cnk-n1-hausdorff-counterexample}
The \(n=1\) Girjoaba--Rasa sequence \(c(1,0)=1,c(1,1)=1/2,c(1,2)=1\) violates the nonincreasing condition for Hausdorff moment sequences.
\end{counterexample}
\begin{proof}
The source sequence is
\[
c(n,k)=\binom{2n}{k}^{-2}
\sum_{j=0}^k \binom{n}{j}^2\binom{n}{k-j}^2,
\]
with binomial coefficients outside their natural range interpreted as \(0\).

For \(n=1\):
\[
c(1,0)=\binom20^{-2}\binom10^2\binom10^2=1.
\]
For \(k=1\),
\[
c(1,1)=\binom21^{-2}
\left(\binom10^2\binom11^2+\binom11^2\binom10^2\right)
=\frac14(1+1)=\frac12.
\]
For \(k=2\),
\[
c(1,2)=\binom22^{-2}
\left(\binom11^2\binom11^2\right)=1.
\]
Thus the full \(n=1\) sequence begins
\[
1,\quad \frac12,\quad 1.
\]

If \((h_k)\) is a Hausdorff moment sequence on \([0,1]\), then
\[
h_k=\int_0^1 t^k\,d\mu(t)
\]
for some positive measure \(\mu\).  Since \(0\le t\le1\), we have \(t^{k+1}\le t^k\), and therefore
\[
h_{k+1}=\int_0^1 t^{k+1}\,d\mu(t)
\le
\int_0^1 t^k\,d\mu(t)=h_k.
\]
Every Hausdorff moment sequence is nonincreasing.

The \(n=1\) Girjoaba--Rasa sequence violates this necessary condition because
\[
c(1,2)=1>\frac12=c(1,1).
\]

The obstruction is actually uniform in \(n\ge1\).  At the right endpoint,
\[
c(n,2n)=1,
\]
because only \(j=n\) contributes and \(\binom{2n}{2n}=1\).  At the adjacent point \(k=2n-1\), only \(j=n-1\) and \(j=n\) contribute, so
\[
\sum_{j=0}^{2n-1}\binom{n}{j}^2\binom{n}{2n-1-j}^2
=2n^2.
\]
Since \(\binom{2n}{2n-1}=2n\), this gives
\[
c(n,2n-1)=\frac{2n^2}{(2n)^2}=\frac12.
\]
Thus \(c(n,2n)>c(n,2n-1)\) for every \(n\ge1\), so no full sequence \(0\le k\le2n\) can be Hausdorff moment in the natural order.
\end{proof}

\begin{lemma}[exact concavity criterion for sqrt(c s power a + d s power b) two-term power symbol]
\label{res:sqrt-two-term-symbol-exact-concavity-criterion}
For \(h(s)=\sqrt{c s^a+d s^b}\), \(c,d>0\), the exact condition for \(h\) to be concave on \((0,\infty)\) is: if \(a=b\), then \(0\le a\le2\); if \(a\ne b\), then \(a(a-2)\le0\), \(b(b-2)\le0\), and either \(a^2+b^2-ab-a-b\le0\) or \((a^2+b^2-ab-a-b)^2\le a(a-2)b(b-2)\).
\end{lemma}
\begin{proof}
Write \(f(s)=c s^a+d s^b\). Then
\[
h''(s)=\frac{2f(s)f''(s)-f'(s)^2}{4f(s)^{3/2}}.
\]
A direct expansion gives
\[
4h(s)^3h''(s)
=c^2 a(a-2)s^{2a-2}
+2cd\bigl(a^2+b^2-ab-a-b\bigr)s^{a+b-2}
+d^2b(b-2)s^{2b-2}.
\]
If \(a=b\), then \(h(s)=\sqrt{c+d}\,s^{a/2}\), so \(h\) is concave on \((0,\infty)\) exactly when \(0\le a\le2\).

Assume now \(a\ne b\), and set
\[
t=\frac cd\,s^{a-b}.
\]
As \(s\) ranges over \((0,\infty)\), so does \(t\). The sign of \(h''\) is therefore the sign of
\[
Q(t)=A t^2+2D t+C,\qquad t>0,
\]
where
\[
A=a(a-2),\qquad C=b(b-2),\qquad D=a^2+b^2-ab-a-b.
\]

Thus \(h\) is concave exactly when \(Q(t)\le0\) for every \(t>0\).

The condition \(Q(t)\le0\) on \((0,\infty)\) first forces
\[
A\le0,\qquad C\le0,
\]
by taking \(t\to\infty\) and \(t\to0^+\). These are equivalent to \(a,b\in[0,2]\).

If \(D\le0\), then all three coefficients of \(Q\) are nonpositive, hence \(Q(t)\le0\) for every \(t>0\).

If \(D>0\), the concave quadratic has its maximum at
\[
t_*=-\frac DA>0
\]
when \(A<0\); endpoint cases \(A=0\) are included by continuity and by the same final inequality. The maximum condition is
\[
Q(t_*)=C-\frac{D^2}{A}\le0.
\]
Since \(A<0\), this is equivalent to
\[
D^2\le AC.
\]
When \(A=0\) and \(D>0\), the function \(Q(t)=2Dt+C\) eventually becomes positive, and \(D^2\le AC=0\) fails. The \(C=0\) endpoint is symmetric. Hence no endpoint exception is missing.

Therefore, for \(a\ne b\), \(h\) is concave on \((0,\infty)\) if and only if
\[
A\le0,\qquad C\le0,\qquad
\left(D\le0\quad\text{or}\quad D^2\le AC\right).
\]

Equivalently,
\[
a(a-2)\le0,\qquad b(b-2)\le0,
\]
and
\[
a^2+b^2-ab-a-b\le0
\]
or
\[
\bigl(a^2+b^2-ab-a-b\bigr)^2
\le a(a-2)b(b-2).
\]

In the Bazhlekova two-term gap regime
\[
1<a<2,\qquad 0<b<a-1,
\]
the two exponents straddle \(1\). In this regime,
\[
D^2-AC=(a-b)^2\left((a-1)^2+(b-1)^2-1\right).
\]
Thus the inner disk condition
\[
(a-1)^2+(b-1)^2\le1
\]
is exactly the condition under which the second-derivative concavity obstruction is absent. The remaining question there is not a concavity question; it requires testing whether \(\sqrt g\) is Bernstein by higher derivatives or proving sign change in the inverse Laplace transform.
\end{proof}

\begin{counterexample}[Keady self-bijection completely monotone inverse not completely monotone example]
\label{res:keady-self-bijection-inverse-cm-negative-example}
The function \(f(x)=x^{-1}+100e^{-x}\) is a completely monotone decreasing bijection \((0,\infty)\to(0,\infty)\), but \(f^{-1}\) is not completely monotone.
\end{counterexample}
\begin{proof}
There exists a completely monotone decreasing bijection
\[
f:(0,\infty)\to(0,\infty)
\]
whose inverse is not completely monotone. One explicit example is
\[
f(x)=\frac1x+100e^{-x}.
\]

For every \(n\ge0\),
\[
(-1)^n f^{(n)}(x)=\frac{n!}{x^{n+1}}+100e^{-x}>0.
\]
Thus \(f\) is strictly completely monotone. Also
\[
f'(x)=-x^{-2}-100e^{-x}<0,
\]
so \(f\) is strictly decreasing. Finally,
\[
\lim_{x\downarrow0}f(x)=\infty,
\qquad
\lim_{x\to\infty}f(x)=0.
\]
Therefore \(f\) is a decreasing bijection from \((0,\infty)\) onto \((0,\infty)\).

Let \(g=f^{-1}\) and write \(y=f(x)\). Inverse differentiation gives
\[
g'''(f(x))=\frac{3(f''(x))^2-f'''(x)f'(x)}{(f'(x))^5}.
\]
For \(f_a(x)=x^{-1}+ae^{-x}\), define
\[
N_a(x)=3(f_a''(x))^2-f_a'''(x)f_a'(x).
\]
A direct calculation gives
\[
N_a(x)=\frac6{x^6}+ae^{-x}\left(-\frac6{x^4}+\frac{12}{x^3}-\frac1{x^2}\right)+2a^2e^{-2x}.
\]
At \(a=100\) and \(x_0=1/8\),
\[
N_{100}(1/8)=1{,}572{,}864-1{,}849{,}600e^{-1/8}+20{,}000e^{-1/4}.
\]
Using \(e^{-1/8}>7/8\) and \(e^{-1/4}<1\),
\[
N_{100}(1/8)<1{,}572{,}864-1{,}849{,}600\cdot\frac78+20{,}000=-25{,}536<0.
\]
Since \(f'(1/8)<0\), the denominator \((f'(1/8))^5\) is negative. Hence
\[
g'''(f(1/8))=\frac{N_{100}(1/8)}{(f'(1/8))^5}>0.
\]
A completely monotone function must satisfy \((-1)^3g'''\ge0\), i.e. \(g'''\le0\). Therefore \(g=f^{-1}\) is not completely monotone.
\end{proof}

The next result applies the preceding method to the fixed source problem \textbf{APP-0028}.
\begin{counterexample}[APP-0028: Qi-Agarwal/Yin divisor-polygamma parity is corrected]
\label{res:qa-divisor-polygamma-even-claim-refuted}
The Qi--Agarwal/Yin claim that \(f_{2\ell}\) is not completely monotone for every \(\ell\in\mathbb N\) is false, because \(f_2\) is completely monotone.

This resolves the fixed application label \textbf{APP-0028}: Decide the Qi-Agarwal/Yin divisor-polygamma parity problem.
\emph{Source reference.} \cite{app-app-0028-source}.
\end{counterexample}
\begin{proof}
The standard polygamma representation gives, for \(k\ge1\),
\[
 (-1)^{k+1}\psi^{(k)}(x)
 =\int_0^\infty {t^ke^{-xt}\over1-e^{-t}}\,dt .
\]
Thus \(\psi^{(k)}\) is a positive completely monotone function whenever
\(k\) is odd.  If \(n\) is odd and \(km=n\), both \(k\) and \(m\) are odd, so
each summand \([\psi^{(k)}]^m\) is a finite product of positive completely
monotone functions.  Finite sums preserve complete monotonicity.

For \(n=2\), the divisor pairs are \((1,2)\) and \((2,1)\), hence
\[
 f_2(x)=[\psi'(x)]^2+\psi''(x).
\]
Qi--Agarwal record the sharp theorem that
\[
 [\psi'(x)]^2+\lambda\psi''(x)
\]
is completely monotone on \((0,\infty)\) if and only if \(\lambda\le1\).
Taking \(\lambda=1\) proves that \(f_2\) is completely monotone.  This refutes
the stated even-parity clause at \(\ell=1\).
\end{proof}

The next result applies the preceding method to the fixed source problem \textbf{APP-0044}.
\begin{theorem}[APP-0044: Bessel ratio quadratic lower bound is false]
\label{res:from-mills-all-l-moment-ratio-quadratic-bound}
Let \(M_L(t)=(-1)^L r^{(L)}(t)=\int_0^\infty u^L e^{-tu-u^2/2}\,du\) for the standard normal Mills ratio. For every \(L\ge0\) and \(t\ge0\), with \(m_L=M_{L+1}/M_L\), one has \(m_L(t)<U_L(t)\), where \(U_L(t)\) is the positive root \(\frac{\sqrt{t^2(t^2+4L+7)^2+4(L+1)(t^2+4L+8)(t^2+4L+6)}-t(t^2+4L+7)}{2(t^2+4L+8)}\).

This resolves the fixed application label \textbf{APP-0044}: Prove \(q\_\nu(u)^2+(2\nu+1)q\_\nu(u)/u+(\nu+1/2)/u^2>1\) for all \(\nu\ge0,u>0\), with \(q\_\nu(u)=I\_{\nu+1}(u)/I\_\nu(u)\).
\emph{Source reference.} \cite{app-app-0044-source}.
\end{theorem}
\begin{proof}
Set
\[
m_L(t)=\frac{M_{L+1}(t)}{M_L(t)}.
\]

The recurrence gives
\[
\frac{M_{L+2}}{M_L}=L+1-tm_L,
\]
\[
\frac{M_{L+3}}{M_L}=(t^2+L+2)m_L-t(L+1),
\]
and
\[
\frac{M_{L+4}}{M_L}=(L+1)(t^2+L+3)-t(t^2+2L+5)m_L.
\]

Substitution into the determinant inequality and expansion give
\[
(L+1)(t^2+4L+6)-t(t^2+4L+7)m_L-(t^2+4L+8)m_L^2>0.
\]

Equivalently,
\[
(t^2+4L+8)m_L^2+t(t^2+4L+7)m_L-(L+1)(t^2+4L+6)<0.
\]

Since \(m_L>0\), the positive root yields the uniform strict bound
\[
m_L(t)<U_L(t),
\]
where
\[
U_L(t)=
\frac{
\sqrt{t^2(t^2+4L+7)^2+4(L+1)(t^2+4L+8)(t^2+4L+6)}
-t(t^2+4L+7)
}
{2(t^2+4L+8)}.
\]

This is the determinant-to-bound extractor. It is a single theorem for every \(L\ge0\).

Define polynomials \(P_n(t)\) by
\[
P_0=1,\qquad P_1=t,\qquad P_{n+1}=tP_n+nP_{n-1}.
\]
Define \(B_n(t)\) by
\[
B_0=0,\qquad B_1=1,\qquad B_{n+1}=nB_{n-1}-tB_n.
\]

Then the same recurrence gives
\[
M_n(t)=(-1)^nP_n(t)r(t)+B_n(t).
\]

The ratio bound \(M_{L+1}<U_LM_L\) becomes
\[
\left((-1)^{L+1}P_{L+1}-(-1)^L U_LP_L\right)r
<U_LB_L-B_{L+1}.
\]

Since \(P_n(t)>0\) for \(t>0\) and \(P_{2j}(0)>0\), the coefficient has sign \((-1)^{L+1}\). Thus the bounds alternate:

For even \(L\ge0\),
\[
r(t)>
\frac{B_{L+1}(t)-U_L(t)B_L(t)}
{P_{L+1}(t)+U_L(t)P_L(t)}.
\]

For odd \(L\ge1\),
\[
r(t)<
\frac{U_L(t)B_L(t)-B_{L+1}(t)}
{P_{L+1}(t)+U_L(t)P_L(t)}.
\]

These formulas are valid at \(t=0\) by direct evaluation of the recurrences and for \(t>0\) by the positivity of \(M_L\).
\end{proof}

The next result applies the preceding method to the fixed source problem \textbf{APP-0035}.
\begin{counterexample}[APP-0035: Keady self-bijection inverse-CM question is negative]
\label{res:keady-q3-self-bijection-inverse-cm-negative-answer}
Keady Question 3 has a negative answer for the self-range case: there is a completely monotone decreasing bijection \((0,\infty)\to(0,\infty)\) whose inverse is not completely monotone.

This resolves the fixed application label \textbf{APP-0035}: Ask whether a CM self-bijection \((0,\infty)\to(0,\infty)\) must have a CM inverse (Keady Q3).
\emph{Source reference.} \cite{app-app-0035-source}.
\end{counterexample}
\begin{proof}
There exists a completely monotone decreasing bijection
\[
f:(0,\infty)\to(0,\infty)
\]
whose inverse is not completely monotone. One explicit example is
\[
f(x)=\frac1x+100e^{-x}.
\]

For every \(n\ge0\),
\[
(-1)^n f^{(n)}(x)=\frac{n!}{x^{n+1}}+100e^{-x}>0.
\]
Thus \(f\) is strictly completely monotone. Also
\[
f'(x)=-x^{-2}-100e^{-x}<0,
\]
so \(f\) is strictly decreasing. Finally,
\[
\lim_{x\downarrow0}f(x)=\infty,
\qquad
\lim_{x\to\infty}f(x)=0.
\]
Therefore \(f\) is a decreasing bijection from \((0,\infty)\) onto \((0,\infty)\).

Let \(g=f^{-1}\) and write \(y=f(x)\). Inverse differentiation gives
\[
g'''(f(x))=\frac{3(f''(x))^2-f'''(x)f'(x)}{(f'(x))^5}.
\]
For \(f_a(x)=x^{-1}+ae^{-x}\), define
\[
N_a(x)=3(f_a''(x))^2-f_a'''(x)f_a'(x).
\]
A direct calculation gives
\[
N_a(x)=\frac6{x^6}+ae^{-x}\left(-\frac6{x^4}+\frac{12}{x^3}-\frac1{x^2}\right)+2a^2e^{-2x}.
\]
At \(a=100\) and \(x_0=1/8\),
\[
N_{100}(1/8)=1{,}572{,}864-1{,}849{,}600e^{-1/8}+20{,}000e^{-1/4}.
\]
Using \(e^{-1/8}>7/8\) and \(e^{-1/4}<1\),
\[
N_{100}(1/8)<1{,}572{,}864-1{,}849{,}600\cdot\frac78+20{,}000=-25{,}536<0.
\]
Since \(f'(1/8)<0\), the denominator \((f'(1/8))^5\) is negative. Hence
\[
g'''(f(1/8))=\frac{N_{100}(1/8)}{(f'(1/8))^5}>0.
\]
A completely monotone function must satisfy \((-1)^3g'''\ge0\), i.e. \(g'''\le0\). Therefore \(g=f^{-1}\) is not completely monotone.
\end{proof}

\appendix

\section{Technical examples and source-specific certificates}
This appendix collects finite certificates, source-specific endpoint checks, numerical witnesses, and applications whose details are useful for verification but would interrupt the main narrative.
\begin{counterexample}[Baskakov alpha one r eight rational function is not completely monotone]
\label{res:baskakov-alpha1-r8-not-cm}
The Baskakov \(\alpha=1\), \(r=8\) function \(f^{[8]}_1(x)=1/((1+x)^8-x^8)\) is not completely monotone on \((0,\infty)\).
\end{counterexample}
\begin{proof}
Set
\[
D_m(x)=(1+x)^{2m}-x^{2m}.
\]
The roots are obtained from
\[
\frac{1+x}{x}=\omega_k,\qquad
\omega_k=e^{\pi i k/m},\qquad k=1,\ldots,2m-1.
\]
Thus
\[
\lambda_{k,m}
=\frac{1}{\omega_k-1}
=-\frac12-\frac{i}{2}\cot\frac{\pi k}{2m}.
\]
At this root,
\[
D_m'(\lambda_{k,m})
=2m\lambda_{k,m}^{2m-1}(\omega_k^{-1}-1),
\]
and direct simplification gives the real residue
\[
R_{k,m}
=\frac{1}{D_m'(\lambda_{k,m})}
=\frac{(-1)^{m+k}}{2m}
\left(2\sin\frac{\pi k}{2m}\right)^{2m-2}.
\]
Consequently the inverse-Laplace density is
\[
\rho_m(t)
=\sum_{k=1}^{2m-1}R_{k,m}e^{\lambda_{k,m}t}.
\]
Pairing conjugate roots gives
\[
\rho_m(t)
=\frac{e^{-t/2}}{2m}
\left[
2^{2m-2}
+2\sum_{k=1}^{m-1}
(-1)^{m+k}
\left(2\sin\frac{\pi k}{2m}\right)^{2m-2}
\cos\!\left(
\frac{t}{2}\cot\frac{\pi k}{2m}
\right)
\right].
\]

This proves the finite trigonometric density formula used by the diagnostic
candidate.

For \(m=2\), this gives
\[
\rho_2(t)=e^{-t/2}\left(1-\cos\frac t2\right)
=2e^{-t/2}\sin^2\frac t4\ge0,
\]
recovering the previously admitted \(r=4,\alpha=1\) seed.

For \(m=3\),
\[
\rho_3(t)
=\frac{e^{-t/2}}6
\left[
16+2\cos\left(\frac{\sqrt3\,t}{2}\right)
-18\cos\left(\frac{t}{2\sqrt3}\right)
\right].
\]
Set \(u=t/(2\sqrt3)\).  Since
\(\cos(3u)=4\cos^3u-3\cos u\),
\[
\rho_3(t)
=\frac{4}{3}e^{-t/2}(1-\cos u)^2(2+\cos u)\ge0.
\]
Thus \(r=6,\alpha=1\) is completely monotone by an explicit positive density.

For \(m=4\), write \(h_4(t)=e^{t/2}\rho_4(t)\).  The density formula gives
\[
h_4(t)
=8+2\cos\frac t2
-\frac{(2-\sqrt2)^3}{4}
\cos\left(\frac{1+\sqrt2}{2}t\right)
-\frac{(2+\sqrt2)^3}{4}
\cos\left(\frac{\sqrt2-1}{2}t\right).
\]
At \(t=10\pi\),
\[
\cos(t/2)=\cos(5\pi)=-1,
\]
and both other cosine terms equal \(-\cos(5\pi\sqrt2)\).  Since
\[
(2-\sqrt2)^3+(2+\sqrt2)^3=40,
\]
we obtain
\[
h_4(10\pi)=6+10\cos(5\pi\sqrt2).
\]
Now
\[
0<5\sqrt2-7<\frac14
\]
because \(7/5<\sqrt2<29/20\).  Hence, with
\(\delta=5\sqrt2-7\),
\[
\cos(5\pi\sqrt2)=\cos(7\pi+\pi\delta)=-\cos(\pi\delta)
<-\cos\frac\pi4=-\frac{\sqrt2}{2}<-\frac35.
\]
Therefore
\[
h_4(10\pi)<6+10\left(-\frac35\right)=0,
\]
and
\[
\rho_4(10\pi)=e^{-5\pi}h_4(10\pi)<0.
\]
Since the inverse-Laplace transform is unique for measures on
\([0,\infty)\), \(f^{[8]}_1(x)=1/((1+x)^8-x^8)\) cannot be completely monotone.
\end{proof}

\begin{counterexample}[Qi h\_lambda degree four transforms are not completely monotone on source ranges]
\label{res:qi-hlambda-x4-source-range-not-cm}
For \(h_\lambda(x)=\Psi(x)-(x^2+\lambda x+12)/(12x^4(x+1)^2)\), \(x^4h_\lambda(x)\) is not completely monotone for every \(\lambda\le0\), and \(x^4(-h_\mu(x))\) is not completely monotone for every \(\mu\ge4\).
\end{counterexample}
\begin{proof}
A completely monotone function must be nonincreasing, so its first derivative cannot be positive at any point.

Use the recurrence expansions at \(x=0^+\):
\[
\psi'(x)=\frac1{x^2}+\zeta(2)-2\zeta(3)x+3\zeta(4)x^2+O(x^3),
\]
and
\[
\psi''(x)=-\frac2{x^3}-2\zeta(3)+6\zeta(4)x+O(x^2).
\]
Therefore
\[
\Psi(x)
=\frac1{x^4}-\frac2{x^3}+\frac{\pi^2}{3x^2}
-\frac{4\zeta(3)}x+O(1).
\]
Also
\[
\frac{x^2+\lambda x+12}{12x^4(1+x)^2}
=\frac1{x^4}+\left(\frac{\lambda}{12}-2\right)\frac1{x^3}
+\left(\frac{37}{12}-\frac{\lambda}{6}\right)\frac1{x^2}
+O(x^{-1}).
\]
Subtracting gives
\[
h_\lambda(x)
=-\frac{\lambda}{12x^3}
+\left(\frac{\lambda}{6}+\frac{\pi^2}{3}-\frac{37}{12}\right)\frac1{x^2}
+O(x^{-1}).
\]

If \(\lambda<0\), then
\[
x^4h_\lambda(x)=-\frac{\lambda}{12}x+O(x^2),
\]
so
\[
\frac{d}{dx}\bigl[x^4h_\lambda(x)\bigr]
=-\frac{\lambda}{12}+O(x)>0
\]
for all sufficiently small \(x>0\). Hence \(x^4h_\lambda\) is not completely monotone.

If \(\lambda=0\), then
\[
x^4h_0(x)
=\left(\frac{\pi^2}{3}-\frac{37}{12}\right)x^2+O(x^3).
\]
Since \(4\pi^2>37\), the leading coefficient is positive, and
\[
\frac{d}{dx}\bigl[x^4h_0(x)\bigr]
=2\left(\frac{\pi^2}{3}-\frac{37}{12}\right)x+O(x^2)>0
\]
for all sufficiently small \(x>0\). Hence \(x^4h_0\) is not completely monotone.

Similarly, for \(\mu\ge4\),
\[
x^4(-h_\mu(x))=\frac{\mu}{12}x+O(x^2),
\]
so
\[
\frac{d}{dx}\bigl[x^4(-h_\mu(x))\bigr]
=\frac{\mu}{12}+O(x)>0
\]
for all sufficiently small \(x>0\). Hence \(x^4(-h_\mu)\) is not completely monotone for every \(\mu\ge4\).

Thus the degree-four transforms required by Qi's conjecture are not completely monotone on the conjectured source ranges. The conjecture in Remark 7.4 is false.
\end{proof}

\begin{proposition}[Keady self-bijection inverse third derivative sign certificate at f(1/8)]
\label{res:keady-inverse-third-derivative-sign-certificate}
For \(f(x)=x^{-1}+100e^{-x}\) and \(g=f^{-1}\), the inverse-derivative certificate gives \(g'''(f(1/8))>0\).
\end{proposition}
\begin{proof}
There exists a completely monotone decreasing bijection
\[
f:(0,\infty)\to(0,\infty)
\]
whose inverse is not completely monotone. One explicit example is
\[
f(x)=\frac1x+100e^{-x}.
\]

For every \(n\ge0\),
\[
(-1)^n f^{(n)}(x)=\frac{n!}{x^{n+1}}+100e^{-x}>0.
\]
Thus \(f\) is strictly completely monotone. Also
\[
f'(x)=-x^{-2}-100e^{-x}<0,
\]
so \(f\) is strictly decreasing. Finally,
\[
\lim_{x\downarrow0}f(x)=\infty,
\qquad
\lim_{x\to\infty}f(x)=0.
\]
Therefore \(f\) is a decreasing bijection from \((0,\infty)\) onto \((0,\infty)\).

Let \(g=f^{-1}\) and write \(y=f(x)\). Inverse differentiation gives
\[
g'''(f(x))=\frac{3(f''(x))^2-f'''(x)f'(x)}{(f'(x))^5}.
\]
For \(f_a(x)=x^{-1}+ae^{-x}\), define
\[
N_a(x)=3(f_a''(x))^2-f_a'''(x)f_a'(x).
\]
A direct calculation gives
\[
N_a(x)=\frac6{x^6}+ae^{-x}\left(-\frac6{x^4}+\frac{12}{x^3}-\frac1{x^2}\right)+2a^2e^{-2x}.
\]
At \(a=100\) and \(x_0=1/8\),
\[
N_{100}(1/8)=1{,}572{,}864-1{,}849{,}600e^{-1/8}+20{,}000e^{-1/4}.
\]
Using \(e^{-1/8}>7/8\) and \(e^{-1/4}<1\),
\[
N_{100}(1/8)<1{,}572{,}864-1{,}849{,}600\cdot\frac78+20{,}000=-25{,}536<0.
\]
Since \(f'(1/8)<0\), the denominator \((f'(1/8))^5\) is negative. Hence
\[
g'''(f(1/8))=\frac{N_{100}(1/8)}{(f'(1/8))^5}>0.
\]
A completely monotone function must satisfy \((-1)^3g'''\ge0\), i.e. \(g'''\le0\). Therefore \(g=f^{-1}\) is not completely monotone.
\end{proof}

\begin{proposition}[Yin p,k-digamma bracket near zero asymptotic x inverse]
\label{res:yin-pk-digamma-near-zero-asymptotic}
For fixed \(p\in\mathbb N\) and \(k>0\), the Yin bracket satisfies \(B_{p,k}(x)=x^{-1}+O(|\log x|)\) as \(x\to0^+\). Consequently \(x^\alpha B_{p,k}(x)\sim x^{\alpha-1}\).
\end{proposition}
\begin{proof}
\emph{Setup.}
Fix \(p\in\mathbb N\) and \(k>0\). Define
\[
B_{p,k}(x)=
\frac1k\log\frac{pkx}{x+k(p+1)}-\psi_{p,k}(x),
\qquad
\delta_{p,k,\alpha}(x)=x^\alpha B_{p,k}(x).
\]

The source gives the finite formula
\[
\psi_{p,k}(x)=\frac1k\log(pk)-\sum_{n=0}^p\frac1{x+nk}.
\]
Therefore
\[
\begin{aligned}
B_{p,k}(x)
&=
\frac1k\log\frac{pkx}{x+k(p+1)}
-\frac1k\log(pk)
+\sum_{n=0}^p\frac1{x+nk}  \\
&=
\frac1k\log\frac{x}{x+k(p+1)}
+\sum_{n=0}^p\frac1{x+nk}.
\end{aligned}
\]

This normal form also gives positivity directly. Put \(t=x/k\). Then
\[
B_{p,k}(x)
=\frac1k\left[
\sum_{n=0}^p\frac1{t+n}
+\log\frac{t}{t+p+1}
\right].
\]
Since
\[
\log\frac{t+p+1}{t}=\int_0^{p+1}\frac{du}{t+u},
\]
and \(u\mapsto(t+u)^{-1}\) is strictly decreasing, the left Riemann sum dominates:
\[
\sum_{n=0}^p\frac1{t+n}
>
\int_0^{p+1}\frac{du}{t+u}.
\]
Hence \(B_{p,k}(x)>0\) for \(x>0\).

This proves the claim.

From the finite normal form,
\[
B_{p,k}(x)
=\frac1x
+\frac1k\log\frac{x}{x+k(p+1)}
+\sum_{n=1}^p\frac1{x+nk}.
\]

As \(x\to0^+\), the finite sum over \(n\ge1\) is \(O(1)\), and
\[
\frac1k\log\frac{x}{x+k(p+1)}=O(|\log x|).
\]
Thus
\[
B_{p,k}(x)=\frac1x+O(|\log x|)
\qquad (x\to0^+).
\]
Equivalently,
\[
xB_{p,k}(x)\to1.
\]
Therefore
\[
\delta_{p,k,\alpha}(x)
=x^\alpha B_{p,k}(x)
=x^{\alpha-1}(xB_{p,k}(x))
\sim x^{\alpha-1}.
\]

This proves the claim.

Assume that \(\delta_{p,k,\alpha}\) is completely monotone on \((0,\infty)\). Then
\[
\delta_{p,k,\alpha}(x)\ge0,\qquad
\delta_{p,k,\alpha}'(x)\le0.
\]
Since \(B_{p,k}(x)>0\), the function \(\delta_{p,k,\alpha}\) is actually positive for every \(x>0\).

Suppose, for contradiction, that \(\alpha>1\). The endpoint asymptotic gives
\[
\lim_{x\to0^+}\delta_{p,k,\alpha}(x)=0.
\]
Fix \(x_0>0\). Positivity gives \(\delta_{p,k,\alpha}(x_0)>0\). Because \(\delta_{p,k,\alpha}(x)\to0\) as \(x\to0^+\), choose \(0<y<x_0\) with
\[
\delta_{p,k,\alpha}(y)<\frac12\delta_{p,k,\alpha}(x_0).
\]
But \(\delta_{p,k,\alpha}'\le0\) means the function is nonincreasing as \(x\) increases, hence for \(0<y<x_0\),
\[
\delta_{p,k,\alpha}(y)\ge\delta_{p,k,\alpha}(x_0),
\]
a contradiction.

Therefore \(\alpha\le1\).

This proves the claim.
\end{proof}

\begin{proposition}[Ramanujan Turan n two upper window has small x counterfamily]
\label{res:ramanujan-turan-n2-upper-window-counterfamily}
For the Ramanujan integral Turan gap, there are parameters \(\alpha\in(0,1/2)\) such that \(H_2(x;\alpha)<0\) for some \(x>0\); specifically, with \(x=e^{-L}\) and \(\alpha_L=1/2-1/(4L)\), this holds for all sufficiently large \(L\).
\end{proposition}
\begin{proof}
Put
\[
A_k(x)=(-1)^k I_R^{(k)}(x)
=\int_0^\infty e^{-xt}\frac{t^{k-1}}{\pi^2+\log^2 t}\,dt
\qquad (k\ge1).
\]
Then
\[
H_2(x;\alpha)=A_2(x)^2-\alpha A_1(x)A_3(x).
\]

We study \(x=e^{-L}\) as \(L\to\infty\). With \(y=xt\),
\[
A_k(e^{-L})
=
e^{kL}\int_0^\infty
\frac{e^{-y}y^{k-1}}
{\pi^2+(L+\log y)^2}\,dy.
\]

For fixed \(k=1,2,3\),
\[
A_k(e^{-L})
=
e^{kL}L^{-2}\Gamma(k)
\left(1-\frac{2\psi(k)}{L}+O(L^{-2})\right),
\]
where \(\psi\) is the digamma function. To justify the expansion, split the \(y\)-integral into the central region \(|\log y|\le L^{2/3}\) and its complement. On the central interval, with \(z=\log y\), expand uniformly:
\[
\frac{1}{\pi^2+(L+\log y)^2}
=
L^{-2}\left(1-\frac{2\log y}{L}+O\left(\frac{1+\log^2 y}{L^2}\right)\right),
\]
and integrate against \(e^{-y}y^{k-1}\). The lower complement \(0<y<e^{-L^{2/3}}\) is \(O(e^{-cL^{2/3}})\) for \(k=1,2,3\) after the change \(y=e^{-u}\), and the upper complement \(y>e^{L^{2/3}}\) is super-exponentially small because of \(e^{-y}\). These errors are \(o(L^{-N})\) for every fixed \(N\), which is more than needed for the displayed expansion.

Therefore
\[
\frac{A_2(e^{-L})^2}{A_1(e^{-L})A_3(e^{-L})}
=
\frac{\Gamma(2)^2}{\Gamma(1)\Gamma(3)}
\left(
1+\frac{2(\psi(1)+\psi(3)-2\psi(2))}{L}+O(L^{-2})
\right).
\]
Using
\[
\psi(1)=-\gamma,\qquad
\psi(2)=1-\gamma,\qquad
\psi(3)=\frac32-\gamma,
\]
we get
\[
\psi(1)+\psi(3)-2\psi(2)=-\frac12,
\]
so
\[
\frac{A_2(e^{-L})^2}{A_1(e^{-L})A_3(e^{-L})}
=
\frac12-\frac{1}{2L}+O(L^{-2}).
\]

Choose
\[
\alpha_L=\frac12-\frac{1}{4L}.
\]
For all sufficiently large \(L\), \(\alpha_L\in(0,1/2)\) and
\[
\frac{A_2(e^{-L})^2}{A_1(e^{-L})A_3(e^{-L})}<\alpha_L.
\]
Consequently
\[
H_2(e^{-L};\alpha_L)
=
A_2(e^{-L})^2-\alpha_L A_1(e^{-L})A_3(e^{-L})
<0.
\]

Complete monotonicity would imply nonnegativity at order zero. Hence \(H_2(\cdot;\alpha_L)\) is not completely monotone for these fixed parameters \(\alpha_L\in(0,1/2)\). The source's universal Turan-window statement is false already in the \(n=2\) upper window.

The new mechanism is a logarithmic-density moment-ratio asymptotic obstruction for quadratic Laplace-moment Turan gaps. It is distinct from the public APP-0014 Stieltjes/complete-Bernstein classification of \(I_R\).
\end{proof}

\begin{proposition}[Bulboaca Zayed gamma quotient numerator has one crossing on one eight]
\label{res:bz-gamma-quotient-n-single-crossing-one-eight}
Let \(N(x)\) be the Bulboaca--Zayed derivative-sign numerator from \(T\)-BZ-gamma-quotient-derivative-normal-form. On \([1,8]\), \(N\) has exactly one zero \(\xi\), with \(N(x)<0\) for \(1\le x<\xi\) and \(N(x)>0\) for \(\xi<x\le8\).
\end{proposition}
\begin{proof}
\emph{Setup.}
Use the notation from the previous Bulboaca--Zayed pass:
\[
G(x)=\log\Gamma(x+1),\qquad
D(x)=\log\frac{x^2+6}{x+6},
\]
and
\[
F(x)=\frac{G(x)}{D(x)}.
\]
For \(x>1\), \(D(x)>0\), and
\[
D'(x)=p(x)=\frac{x^2+12x-6}{(x^2+6)(x+6)}.
\]
Let
\[
N_0(x)=\psi(x+1)D(x)-G(x)D'(x).
\]
The derivative-sign numerator recorded in
the corresponding result is
\[
N(x)=((x^2+6)(x+6))\,N_0(x),
\]
so \(N\) and \(N_0\) have the same sign on \(x>1\).

The source itself proves the decreasing part on \((-1,1)\), and the earlier
the preceding theorem proves \(F'(x)>0\) for
\(x\ge8\).  It remains to prove a single sign change of \(N_0\) on \((1,8]\).

Put
\[
P(x)=x^2+12x-6,\qquad Q(x)=(x^2+6)(x+6),
\]
so \(p=P/Q\), and define
\[
r(x)=\frac{\psi(x+1)}{p(x)}=\psi(x+1)\frac{Q(x)}{P(x)}.
\]
Then
\[
P(x)^2r'(x)=S(x),
\]
where
\[
S(x)=\psi'(x+1)Q(x)P(x)+\psi(x+1)K(x),
\]
and
\[
K(x)=Q'(x)P(x)-Q(x)P'(x)
=x^4+24x^3+48x^2-144x-468.
\]
Direct differentiation gives
\[
S'(x)=P(x)\{Q(x)\psi''(x+1)+2Q'(x)\psi'(x+1)\}+K'(x)\psi(x+1).
\]

For \(y=x+1\ge2\), use the standard elementary bounds
\[
\psi(y)>\log y-\frac1y,\qquad
\psi'(y)>\frac1y+\frac{1}{2y^2},\qquad
\psi''(y)>-\frac1{y^2}-\frac2{y^3}.
\]
The last inequality follows from the series
\[
\psi''(y)=-2\sum_{n=0}^{\infty}(n+y)^{-3}
\]
and an integral upper bound for the sum.

Substituting these bounds into \(S'\) yields
\[
S'(x)>
P(x)\frac{5x^4+30x^3+57x^2+12x-90}{(x+1)^3}
K'(x)\left(\log(x+1)-\frac1{x+1}\right).
\]
On \(x\ge1\), \(P(x)>0\).  The quartic numerator has value \(14\) at \(x=1\)
and positive derivative
\[
20x^3+90x^2+114x+12.
\]
Also
\[
K'(x)=4x^3+72x^2+96x-144
\]
has value \(28\) at \(x=1\) and positive derivative for \(x\ge1\), while
\(\log(x+1)-1/(x+1)>0\) for \(x\ge1\).  Hence
\[
S'(x)>0\qquad(1\le x\le8).
\]

At the endpoints,
\[
S(1)=343\left(\frac{\pi^2}{6}-1\right)-539(1-\gamma)<0,
\]
for example using \(\pi^2<987/100\) and \(\gamma<5773/10000\), which gives
\(-66003/10000<0\).  Also \(S(8)>0\), because \(K(8)=17836>0\),
\(\psi'(9)>0\), and \(\psi(9)=H_8-\gamma>0\).  Therefore \(S\) has exactly
one zero on \([1,8]\).  Since \(P^2>0\), \(r'\) has exactly one zero; \(r\)
decreases first and then increases.

Now define
\[
I(x)=D(x)r(x)-G(x).
\]
Then
\[
N_0(x)=p(x)I(x),\qquad I'(x)=D(x)r'(x).
\]
Because \(D(x)>0\) for \(x>1\), \(I\) decreases and then increases with the
same turning point as \(r\).  Moreover \(\lim_{x\downarrow1}I(x)=0\), so
\(I(x)<0\) immediately to the right of \(1\).

Finally,
\[
N_0(8)=(H_8-\gamma)\log5-\frac{11}{70}\log(40320)>0.
\]
For instance, the bounds
\[
\gamma<\frac{5773}{10000},\quad \log5>\frac{1609}{1000},
\quad \log(40320)<\frac{10605}{1000}
\]
give the lower bound
\[
\left(\frac{761}{280}-\frac{5773}{10000}\right)\frac{1609}{1000}
-\frac{11}{70}\frac{10605}{1000}
=\frac{124435951}{70000000}>0.
\]
Thus \(I\), and therefore \(N_0\) and \(N\), has exactly one zero on
\((1,8]\).  Denote it by \(\xi\).  The sign pattern is
\[
N(x)<0\quad(1<x<\xi),\qquad
N(\xi)=0,\qquad
N(x)>0\quad(\xi<x\le8).
\]

A high-precision numerical check places
\[
\xi=1.12620706105164346558\ldots,
\]
matching the source's reported value.  The numerical value is not used in the
proof.

The raw numerator has a removable zero at \(x=1\), since \(G(1)=D(1)=0\).  Its
quadratic leading coefficient is
\[
\lim_{x\downarrow1}\frac{N_0(x)}{(x-1)^2}
=
\frac{7(\pi^2/6-1)-11(1-\gamma)}{98}<0,
\]
again by \(\pi^2<987/100\) and \(\gamma<5773/10000\), which gives
\(-1347/980000<0\).  Therefore any normalization dividing out the removable
\((x-1)^2\) factor is negative at the left endpoint.

The single-crossing candidate is true:
\[
T\text{-BZ-gamma-quotient-N-single-crossing-one-eight}.
\]
Together with the source's theorem on \((-1,1)\) and the previous local
right-tail theorem for \(x\ge8\), this proves the critical-window reduction
\[
T\text{-BZ-gamma-quotient-critical-window-reduction}
\]
and then the source problem
\[
T\text{-Bulboaca-Zayed-gamma-quotient-full-monotonicity}.
\]

The finite-envelope candidate is not separately promoted: the proof above is a
clean analytic ratio-kernel argument rather than the finite rational cover
specified in that candidate.  The diagnostic obstruction-map candidate is also
not needed after the single-crossing proof.
\end{proof}

\begin{proposition}[analytic denominator reduced assembly proves terminal endpoint]
\label{res:q2-analytic-denominator-reduced-assembly-terminal-certificate}
The terminal endpoint follows from a reduced analytic-denominator assembly: the true covers \((0,9/50]\) and \([9/10,1)\), the endpoint signs on \(J\), an analytic-denominator or direct-log finite-middle cover of \([9/50,1409/5000]\cup[293/1000,9/10]\), and a compact one-crossing/dominance certificate on \([1409/5000,293/1000]\setminus J\). This determines \(L_2=Q_2(\xi)\) and \(\mathcal I_2=(Q_2(\xi),3]\).
\end{proposition}
\begin{proof}
This roll proves a true auxiliary finite-middle outside-cover node:

the corresponding result.

The replayable certificate is:

\begin{verbatim}
\end{verbatim}

It proves
\[
Q_2(x)<q_J
\]
on
\[
\left[\frac9{50},\frac{1409}{5000}\right]
\cup
\left[\frac{293}{1000},\frac9{10}\right].
\]

The left bridge \([9/50,221050/1000000]\) uses the already-derived identity
\[
3xZ_4(x)-Z_3(x)
=2x^{-3}+\sum_{k=1}^{\infty}\frac{2x-k}{(x+k)^4}.
\]
For \(x<1/2\), this gives
\[
R(x)>Z_3(1/x).
\]
The certificate proves
\[
Z_3(1/x)>x^{1157/500}
\]
on \([9/50,221050/1000000]\). Since
\[
\frac{1157}{500}<\frac{115727}{50000}<q_J,
\]
this gives \(Q_2(x)<q_J\) on the bridge.

The remaining two subintervals use the direct comparison
\[
\log R(x)>\frac{115727}{50000}\log x.
\]
Since \(\log x<0\), this implies
\[
Q_2(x)<\frac{115727}{50000}<q_J.
\]
The script uses power-of-two range reduction for the atanh logarithm enclosure:
\[
\log r=\log(2^k r)-k\log 2,
\]
so every logarithm is evaluated with an argument in \([1,2)\).

The certificate output was:

\begin{verbatim}
REMAINING_FINITE_MIDDLE_COVER_OK
ZETA_LEFT_SUBCOVER
interval_count 245
domain [180000/1000000,221050/1000000]
min_width 1 max_width 1283
worst_interval [180000/1000000,181282/1000000]
worst_margin_positive True
DIRECT_LEFT_SUBCOVER
interval_count 3739
domain [221050/1000000,281800/1000000]
min_width 1 max_width 238
worst_interval [280389/1000000,280392/1000000]
worst_margin_positive True
DIRECT_RIGHT_SUBCOVER
interval_count 3801
domain [293000/1000000,900000/1000000]
min_width 1 max_width 18969
worst_interval [293124/1000000,293126/1000000]
worst_margin_positive True
\end{verbatim}

Together with the earlier true covers \((0,9/50]\) and \([9/10,1)\), this covers all outside-compact points below \(q_J\).

This roll also proves:

the corresponding result.

The replayable certificate is:

\begin{verbatim}
\end{verbatim}

Let
\[
H(x)=\Lambda(x)+x\Lambda'(x).
\]
Using \(Z_s'(x)=-sZ_{s+1}(x)\) and
\[
\frac{d}{dx}Z_s(1/x)=s x^{-2} Z_{s+1}(1/x),
\]
the script builds a rational interval enclosure for \(H\) using \(Z_3,\ldots,Z_6\) at \(x\) and \(Z_3,\ldots,Z_5\) at \(1/x\). It proves
\[
H(x)>0
\qquad
\left(\frac{1409}{5000}\le x\le\frac{293}{1000}\right).
\]

Since
\[
G'(x)=\log x\,H(x)
\]
and \(\log x<0\) on the compact bracket, this gives
\[
G'(x)<0.
\]
The existing endpoint certificate proves
\[
G(287345/1000000)>0,
\qquad
G(287346/1000000)<0.
\]
Therefore \(G\) has a unique zero
\[
\xi\in J=\left[\frac{287345}{1000000},\frac{287346}{1000000}\right],
\]
with \(G>0\) to the left of \(J\) and \(G<0\) to the right of \(J\) on the compact bracket.

Because
\[
Q_2'(x)=\frac{G(x)}{x(\log x)^2},
\]
\(Q_2\) is increasing before \(\xi\) and decreasing after \(\xi\) on the compact bracket.

The certificate output was:

\begin{verbatim}
COMPACT_MONOTONICITY_CERTIFICATE_OK
endpoint_signs G_left_positive G_right_negative
h_interval_count 128
domain [281800/1000000,293000/1000000]
min_width 87 max_width 88
worst_interval [281887/1000000,281975/1000000]
worst_H_lower_positive True
\end{verbatim}

The true cover package is now:

\[
(0,9/50]\cup
\left[9/50,\frac{1409}{5000}\right]\cup
\left[\frac{293}{1000},9/10\right]\cup
[9/10,1).
\]

All points outside the compact bracket have \(Q_2(x)<q_J\), and \(q_J<Q_2(\xi)\) by the certified inner witness. On the compact bracket, \(Q_2\) has a unique maximum at the unique zero \(\xi\in J\). Thus
\[
L_2=Q_2(\xi).
\]

At \(\beta=Q_2(\xi)\), the defining inequality has equality at \(x=\xi\), so the lower endpoint is excluded. The known upper endpoint remains included. Therefore
\[
\mathcal I_2=(Q_2(\xi),3].
\]

This proves the terminal exact endpoint node.
\end{proof}

\begin{proposition}[compact Taylor monotonicity around J plus outside covers proves terminal endpoint]
\label{res:q2-compact-taylor-monotonicity-j-terminal-certificate}
Using the true endpoint-sign certificate on \(J=[287345/1000000,287346/1000000]\), there is a Taylor-model or derivative-bounded rational certificate proving \(G>0\) on \([1409/5000,287345/1000000]\), \(G<0\) on \([287346/1000000,293/1000]\), and dominance below the witness \(q_J\); combined with finite-middle covers outside the compact bracket and the true endpoint outside covers, this determines \(L_2=Q_2(\xi)\) and \(\mathcal I_2=(Q_2(\xi),3]\).
\end{proposition}
\begin{proof}
This roll proves a true auxiliary finite-middle outside-cover node:

the corresponding result.

The replayable certificate is:

\begin{verbatim}
\end{verbatim}

It proves
\[
Q_2(x)<q_J
\]
on
\[
\left[\frac9{50},\frac{1409}{5000}\right]
\cup
\left[\frac{293}{1000},\frac9{10}\right].
\]

The left bridge \([9/50,221050/1000000]\) uses the already-derived identity
\[
3xZ_4(x)-Z_3(x)
=2x^{-3}+\sum_{k=1}^{\infty}\frac{2x-k}{(x+k)^4}.
\]
For \(x<1/2\), this gives
\[
R(x)>Z_3(1/x).
\]
The certificate proves
\[
Z_3(1/x)>x^{1157/500}
\]
on \([9/50,221050/1000000]\). Since
\[
\frac{1157}{500}<\frac{115727}{50000}<q_J,
\]
this gives \(Q_2(x)<q_J\) on the bridge.

The remaining two subintervals use the direct comparison
\[
\log R(x)>\frac{115727}{50000}\log x.
\]
Since \(\log x<0\), this implies
\[
Q_2(x)<\frac{115727}{50000}<q_J.
\]
The script uses power-of-two range reduction for the atanh logarithm enclosure:
\[
\log r=\log(2^k r)-k\log 2,
\]
so every logarithm is evaluated with an argument in \([1,2)\).

The certificate output was:

\begin{verbatim}
REMAINING_FINITE_MIDDLE_COVER_OK
ZETA_LEFT_SUBCOVER
interval_count 245
domain [180000/1000000,221050/1000000]
min_width 1 max_width 1283
worst_interval [180000/1000000,181282/1000000]
worst_margin_positive True
DIRECT_LEFT_SUBCOVER
interval_count 3739
domain [221050/1000000,281800/1000000]
min_width 1 max_width 238
worst_interval [280389/1000000,280392/1000000]
worst_margin_positive True
DIRECT_RIGHT_SUBCOVER
interval_count 3801
domain [293000/1000000,900000/1000000]
min_width 1 max_width 18969
worst_interval [293124/1000000,293126/1000000]
worst_margin_positive True
\end{verbatim}

Together with the earlier true covers \((0,9/50]\) and \([9/10,1)\), this covers all outside-compact points below \(q_J\).

This roll also proves:

the corresponding result.

The replayable certificate is:

\begin{verbatim}
\end{verbatim}

Let
\[
H(x)=\Lambda(x)+x\Lambda'(x).
\]
Using \(Z_s'(x)=-sZ_{s+1}(x)\) and
\[
\frac{d}{dx}Z_s(1/x)=s x^{-2} Z_{s+1}(1/x),
\]
the script builds a rational interval enclosure for \(H\) using \(Z_3,\ldots,Z_6\) at \(x\) and \(Z_3,\ldots,Z_5\) at \(1/x\). It proves
\[
H(x)>0
\qquad
\left(\frac{1409}{5000}\le x\le\frac{293}{1000}\right).
\]

Since
\[
G'(x)=\log x\,H(x)
\]
and \(\log x<0\) on the compact bracket, this gives
\[
G'(x)<0.
\]
The existing endpoint certificate proves
\[
G(287345/1000000)>0,
\qquad
G(287346/1000000)<0.
\]
Therefore \(G\) has a unique zero
\[
\xi\in J=\left[\frac{287345}{1000000},\frac{287346}{1000000}\right],
\]
with \(G>0\) to the left of \(J\) and \(G<0\) to the right of \(J\) on the compact bracket.

Because
\[
Q_2'(x)=\frac{G(x)}{x(\log x)^2},
\]
\(Q_2\) is increasing before \(\xi\) and decreasing after \(\xi\) on the compact bracket.

The certificate output was:

\begin{verbatim}
COMPACT_MONOTONICITY_CERTIFICATE_OK
endpoint_signs G_left_positive G_right_negative
h_interval_count 128
domain [281800/1000000,293000/1000000]
min_width 87 max_width 88
worst_interval [281887/1000000,281975/1000000]
worst_H_lower_positive True
\end{verbatim}

The true cover package is now:

\[
(0,9/50]\cup
\left[9/50,\frac{1409}{5000}\right]\cup
\left[\frac{293}{1000},9/10\right]\cup
[9/10,1).
\]

All points outside the compact bracket have \(Q_2(x)<q_J\), and \(q_J<Q_2(\xi)\) by the certified inner witness. On the compact bracket, \(Q_2\) has a unique maximum at the unique zero \(\xi\in J\). Thus
\[
L_2=Q_2(\xi).
\]

At \(\beta=Q_2(\xi)\), the defining inequality has equality at \(x=\xi\), so the lower endpoint is excluded. The known upper endpoint remains included. Therefore
\[
\mathcal I_2=(Q_2(\xi),3].
\]

This proves the terminal exact endpoint node.
\end{proof}

\begin{proposition}[remaining finite middle direct log cover plus compact certificate proves terminal endpoint]
\label{res:q2-remaining-finite-middle-log-cover-terminal-certificate}
Using the true outside covers \((0,9/50]\) and \([9/10,1)\), the true endpoint signs on \(J\), and the true enclosure lemmas, there is a finite rational direct-log cover of \([9/50,1409/5000]\cup[293/1000,9/10]\) proving \(Q_2(x)<q_J\) on every interval, together with a compact one-crossing/dominance certificate on \([1409/5000,293/1000]\setminus J\). This determines \(L_2=Q_2(\xi)\) and \(\mathcal I_2=(Q_2(\xi),3]\).
\end{proposition}
\begin{proof}
This roll proves a true auxiliary finite-middle outside-cover node:

the corresponding result.

The replayable certificate is:

\begin{verbatim}
\end{verbatim}

It proves
\[
Q_2(x)<q_J
\]
on
\[
\left[\frac9{50},\frac{1409}{5000}\right]
\cup
\left[\frac{293}{1000},\frac9{10}\right].
\]

The left bridge \([9/50,221050/1000000]\) uses the already-derived identity
\[
3xZ_4(x)-Z_3(x)
=2x^{-3}+\sum_{k=1}^{\infty}\frac{2x-k}{(x+k)^4}.
\]
For \(x<1/2\), this gives
\[
R(x)>Z_3(1/x).
\]
The certificate proves
\[
Z_3(1/x)>x^{1157/500}
\]
on \([9/50,221050/1000000]\). Since
\[
\frac{1157}{500}<\frac{115727}{50000}<q_J,
\]
this gives \(Q_2(x)<q_J\) on the bridge.

The remaining two subintervals use the direct comparison
\[
\log R(x)>\frac{115727}{50000}\log x.
\]
Since \(\log x<0\), this implies
\[
Q_2(x)<\frac{115727}{50000}<q_J.
\]
The script uses power-of-two range reduction for the atanh logarithm enclosure:
\[
\log r=\log(2^k r)-k\log 2,
\]
so every logarithm is evaluated with an argument in \([1,2)\).

The certificate output was:

\begin{verbatim}
REMAINING_FINITE_MIDDLE_COVER_OK
ZETA_LEFT_SUBCOVER
interval_count 245
domain [180000/1000000,221050/1000000]
min_width 1 max_width 1283
worst_interval [180000/1000000,181282/1000000]
worst_margin_positive True
DIRECT_LEFT_SUBCOVER
interval_count 3739
domain [221050/1000000,281800/1000000]
min_width 1 max_width 238
worst_interval [280389/1000000,280392/1000000]
worst_margin_positive True
DIRECT_RIGHT_SUBCOVER
interval_count 3801
domain [293000/1000000,900000/1000000]
min_width 1 max_width 18969
worst_interval [293124/1000000,293126/1000000]
worst_margin_positive True
\end{verbatim}

Together with the earlier true covers \((0,9/50]\) and \([9/10,1)\), this covers all outside-compact points below \(q_J\).

This roll also proves:

the corresponding result.

The replayable certificate is:

\begin{verbatim}
\end{verbatim}

Let
\[
H(x)=\Lambda(x)+x\Lambda'(x).
\]
Using \(Z_s'(x)=-sZ_{s+1}(x)\) and
\[
\frac{d}{dx}Z_s(1/x)=s x^{-2} Z_{s+1}(1/x),
\]
the script builds a rational interval enclosure for \(H\) using \(Z_3,\ldots,Z_6\) at \(x\) and \(Z_3,\ldots,Z_5\) at \(1/x\). It proves
\[
H(x)>0
\qquad
\left(\frac{1409}{5000}\le x\le\frac{293}{1000}\right).
\]

Since
\[
G'(x)=\log x\,H(x)
\]
and \(\log x<0\) on the compact bracket, this gives
\[
G'(x)<0.
\]
The existing endpoint certificate proves
\[
G(287345/1000000)>0,
\qquad
G(287346/1000000)<0.
\]
Therefore \(G\) has a unique zero
\[
\xi\in J=\left[\frac{287345}{1000000},\frac{287346}{1000000}\right],
\]
with \(G>0\) to the left of \(J\) and \(G<0\) to the right of \(J\) on the compact bracket.

Because
\[
Q_2'(x)=\frac{G(x)}{x(\log x)^2},
\]
\(Q_2\) is increasing before \(\xi\) and decreasing after \(\xi\) on the compact bracket.

The certificate output was:

\begin{verbatim}
COMPACT_MONOTONICITY_CERTIFICATE_OK
endpoint_signs G_left_positive G_right_negative
h_interval_count 128
domain [281800/1000000,293000/1000000]
min_width 87 max_width 88
worst_interval [281887/1000000,281975/1000000]
worst_H_lower_positive True
\end{verbatim}

The true cover package is now:

\[
(0,9/50]\cup
\left[9/50,\frac{1409}{5000}\right]\cup
\left[\frac{293}{1000},9/10\right]\cup
[9/10,1).
\]

All points outside the compact bracket have \(Q_2(x)<q_J\), and \(q_J<Q_2(\xi)\) by the certified inner witness. On the compact bracket, \(Q_2\) has a unique maximum at the unique zero \(\xi\in J\). Thus
\[
L_2=Q_2(\xi).
\]

At \(\beta=Q_2(\xi)\), the defining inequality has equality at \(x=\xi\), so the lower endpoint is excluded. The known upper endpoint remains included. Therefore
\[
\mathcal I_2=(Q_2(\xi),3].
\]

This proves the terminal exact endpoint node.
\end{proof}

\begin{proposition}[s=9 inverse zeta tail floor formula with reusable certificate template]
\label{res:zeta-tail-s9-reusable-certificate}
The next case \(s=9\) has an exact formula for \(\lfloor\zeta_n(9)^{-1}\rfloor\) with a reproducible certificate following the same reusable approximant, telescoping-sign, and gap-trap template intended for later exponents.
\end{proposition}
\begin{proof}
For integer \(s>1\), the Hurwitz-tail Euler--Maclaurin expansion has the form
\[
\zeta_n(s)
=n^{1-s}\left[
\frac1{s-1}+\frac{z}{2}
+\sum_{r\ge1}\frac{B_{2r}}{(2r)!}(s)_{2r-1}z^{2r}
\right],
\qquad z=\frac1n.
\]
Formally inverting the bracket gives
\[
\zeta_n(s)^{-1}\sim n^{s-1}G_s(z).
\]

The replay script
computes \(G_s\) symbolically and verifies two things:

1. For \(s=8\), the coefficients through \(z^{15}\) reproduce the staged \(A_8\) and \(B_8-A_8\) coefficients.
2. For \(s=9\), truncating \(n^8G_9(1/n)\) through \(z^{14}\) gives an \(A_9\), and adding the \(z^{15}\) term gives \(B_9\), with exact telescoping residuals of the expected signs.

The generated \(s=9\) lower approximant is
\[
\begin{aligned}
A_9(n)=&8n^8-32n^7+80n^6-128n^5+120n^4-64n^3
+\frac{624}{7}n^2-\frac{512}{7}n-\frac{3324}{7}\\
&+\frac{90304}{21n^2}+\frac{90304}{21n^3}
-\frac{10280944}{245n^4}-\frac{64846304}{735n^5}
+\frac{109940976}{245n^6}.
\end{aligned}
\]
The upper approximant is
\[
B_9(n)=A_9(n)+\frac{384388288}{245n^7}.
\]

Let \(R_A(n)=1/A_9(n)\) and \(R_B(n)=1/B_9(n)\). Exact denominator clearing gives:
\[
R_A(k)-R_A(k+1)-\frac1{k^9}>0
\qquad(k\ge5),
\]
because the numerator, after substituting \(k=m+5\), has all coefficients positive. Similarly,
\[
\frac1{k^9}-\left(R_B(k)-R_B(k+1)\right)>0
\qquad(k\ge1),
\]
because the denominator-cleared numerator, after substituting \(k=m+1\), has all coefficients positive.

Therefore, for \(n\ge5\),
\[
A_9(n)<\zeta_n(9)^{-1}<B_9(n).
\]

This proves the inverse-window part of the reusable template and promotes the corresponding result.

The replay script also verifies the no-integer gap.

Write \(P_9(n)\) for the polynomial part of \(A_9(n)\):
\[
P_9(n)=8n^8-32n^7+80n^6-128n^5+120n^4-64n^3
+\frac{624}{7}n^2-\frac{512}{7}n-\frac{3324}{7}.
\]
Modulo \(7\), the fractional numerator of \(P_9(n)\) is
\[
n^2-n+1\pmod 7,
\]
so the fractional part of \(P_9(n)\) is one of
\[
0,\quad \frac17,\quad \frac37,\quad \frac67.
\]

The correction \(A_9(n)-P_9(n)\) is positive for \(n\ge9\); after denominator clearing and substituting \(n=m+9\), the numerator has all coefficients positive. For \(n\ge174\), the upper correction is bounded by
\[
B_9(n)-P_9(n)
<
\frac{90304}{21n^2}
+\frac{90304}{21n^3}
+\frac{109940976}{245n^6}
+\frac{384388288}{245n^7}
<\frac17.
\]
Therefore \(A_9(n)\) and \(B_9(n)\) cannot straddle an integer for \(n\ge174\). Exact rational arithmetic verifies
\[
\lfloor A_9(n)\rfloor=\lfloor B_9(n)\rfloor
\qquad(9\le n\le173).
\]
Consequently,
\[
\left\lfloor\zeta_n(9)^{-1}\right\rfloor=\lfloor A_9(n)\rfloor
\qquad(n\ge9).
\]

For \(1\le n\le8\), the exact values are
\[
\begin{array}{c|rrrrrrrr}
n&1&2&3&4&5&6&7&8\\
\hline
\lfloor\zeta_n(9)^{-1}\rfloor&0&497&18093&224086&1543530&7360226&27295287&84349541.
\end{array}
\]
For \(n=1\), this follows from \(\zeta(9)>1\). For \(2\le n\le8\), the replay script uses exact rational partial sums through \(N=100\) and the integral tail bound
\[
\sum_{k=N+1}^{\infty}k^{-9}<\frac{1}{8N^8}
\]
to prove
\[
\frac1{M_n+1}<\zeta_n(9)<\frac1{M_n}.
\]

This proves the claim.
\end{proof}

\begin{lemma}[Yang Detemple R normalized \(n=0\) complete monotonicity certificate]
\label{res:yang-detemple-r-n0-cm-certificate}
The function \(F_0(x)=\left(24x^2+21/5\right)(\psi(x+1/2)-\log x)-1\) is completely monotone on \((0,\infty)\), with nonzero finite tail coefficient \(2071/33600\) at order \(x^{-4}\).
\end{lemma}
\begin{proof}
The standard asymptotic expansion of \(\psi\) gives
\[
R(x)=\frac{1}{24x^2}-\frac{7}{960x^4}+O(x^{-6}).
\]
For the normalized \(n=0\) expression
\[
F(x)=(a_0+24x^2)R(x)-b_0
\]
to have a finite nonzero \(x^{-4}\) first tail term, the constant and \(x^{-2}\) terms must vanish.  Hence
\[
b_0=1,
\]
and
\[
\frac{a_0}{24}+24\left(-\frac{7}{960}\right)=0,
\]
so
\[
a_0=\frac{21}{5}.
\]
Thus the normalized candidate is
\[
F_0(x)=\left(24x^2+\frac{21}{5}\right)R(x)-1.
\]

Using
\[
\psi(z)=\int_0^\infty\left(\frac{e^{-t}}{t}-\frac{e^{-zt}}{1-e^{-t}}\right)dt
\]
and
\[
\log x=\int_0^\infty\frac{e^{-t}-e^{-xt}}{t}\,dt,
\]
we get
\[
R(x)=\int_0^\infty e^{-xt}r(t)\,dt,
\qquad
r(t)=\frac1t-\frac{1}{2\sinh(t/2)}.
\]
The kernel has \(r(0)=0\) and \(r'(0)=1/24\).  Twice integrating by parts gives
\[
x^2R(x)-\frac1{24}=\int_0^\infty e^{-xt}r''(t)\,dt.
\]
Therefore
\[
F_0(x)=\int_0^\infty e^{-xt}K(t)\,dt,
\qquad
K(t)=24r''(t)+\frac{21}{5}r(t).
\]

It remains to prove \(K(t)\ge0\).

Put \(y=t/2\) and \(h(y)=1/y-\operatorname{csch}y\).  Since \(r(t)=h(y)/2\) and \(d/dt=(1/2)d/dy\),
\[
K(t)=3h''(y)+\frac{21}{10}h(y).
\]
With
\[
h''(y)=\frac{2}{y^3}-\operatorname{csch}y\,\coth^2y-\operatorname{csch}^3y
\]
and \(\coth^2y=1+\operatorname{csch}^2y\), this is equivalent to
\[
K(t)=
\frac{
S(y)
}{
10y^3\sinh^3y
},
\]
where
\[
S(y)=(60+21y^2)\sinh^3y-51y^3\sinh^2y-60y^3.
\]

Using
\[
\sinh^3y=\frac{\sinh3y-3\sinh y}{4},
\qquad
\sinh^2y=\frac{\cosh2y-1}{2},
\]
one obtains the power series
\[
S(y)=\sum_{m=4}^{\infty} c_my^{2m+1},
\]
where
\[
c_m=
\frac{
15(3^{2m+1}-3)
+\frac{21}{4}(3^{2m-1}-3)(2m)(2m+1)
-\frac{51}{2}2^{2m-2}(2m-1)(2m)(2m+1)
}{
(2m+1)!}.
\]
The first coefficients are already positive:
\[
c_4=\frac{2071}{2520},\qquad
c_5=\frac{1909}{10080}.
\]
For \(m\ge6\), the positive \(3^{2m-1}\)-terms dominate the \(2^{2m-2}\) term; explicitly, after multiplying by \((2m+1)!\), the ratio of the positive part to the negative part is increasing from \(m=6\) onward and is already \(>1\) at \(m=6\).  Thus \(c_m>0\) for every \(m\ge4\).

Hence \(S(y)>0\) for \(y>0\), and therefore \(K(t)>0\) for \(t>0\).

By Bernstein's theorem,
\[
F_0(x)=\left(24x^2+\frac{21}{5}\right)(\psi(x+1/2)-\log x)-1
\]
is completely monotone on \((0,\infty)\).

Since \(K(t)=\frac{2071}{201600}t^3+O(t^5)\), Watson's lemma gives
\[
F_0(x)=\frac{2071}{33600}x^{-4}+O(x^{-6}).
\]
This verifies the nonzero finite tail condition for \(n=0\).
\end{proof}

\begin{lemma}[complete Bernstein finite origin difference quotient is Stieltjes]
\label{res:cbf-difference-quotient-stieltjes}
If \(G\) is a complete Bernstein function with finite \(G(0+)\), then \((G(x)-G(0+))/x\) is a Stieltjes function.
\end{lemma}
\begin{proof}
Let \(q\ge1\) and \(b_i>a_i>0\). Define
\[
F(x)={}_{q+1}F_q(1,a_1,\ldots,a_q;b_1,\ldots,b_q;-x)
\]
and
\[
R(x)=
\frac{{}_{q+1}F_q(1,a_1+1,\ldots,a_q+1;b_1+1,\ldots,b_q+1;-x)}
{{}_{q+1}F_q(1,a_1,\ldots,a_q;b_1,\ldots,b_q;-x)}.
\]
Then \(R\) is a Stieltjes function on \((0,\infty)\), hence completely monotone. Moreover
\[
xR(x)
\]
is a complete Bernstein function up to the positive scalar \(\prod_i b_i/a_i\).

For \(q=1\), this gives the Gauss/Beta slice:
\[
\frac{{}_2F_1(1,a+1;b+1;-x)}{{}_2F_1(1,a;b;-x)}
\]
is Stieltjes and completely monotone for \(b>a>0\).

Karp--Sitnik's positive-density representation gives a compact-support Stieltjes form for \(F\):
\[
F(x)=\int_0^1 \frac{d\mu(s)}{1+sx},
\qquad \mu\ge0,
\qquad \mu([0,1])=F(0)=1.
\]
Thus \(F\) is a nonzero Stieltjes function.

By the standard Stieltjes/complete-Bernstein duality, if \(F\) is nonzero Stieltjes, then
\[
G(x)=\frac1{F(x)}
\]
is a complete Bernstein function. Since \(G(0)=1\), the complete Bernstein representation gives
\[
G(x)=1+\alpha x+\int_0^\infty \frac{x}{x+t}\,d\nu(t),
\qquad \alpha\ge0,\quad \nu\ge0.
\]
Therefore
\[
\frac{G(x)-1}{x}
=\alpha+\int_0^\infty \frac{1}{x+t}\,d\nu(t)
\]
is Stieltjes.

Karp--Sitnik's contiguous defect identity for \(\sigma=1,\delta=1\) is
\[
1-F(x)
=x\,{}_{q+1}F_q(1,a_1+1,\ldots,a_q+1;b_1+1,\ldots,b_q+1;-x)
\prod_{i=1}^q\frac{a_i}{b_i}.
\]
Hence
\[
R(x)=
\prod_{i=1}^q\frac{b_i}{a_i}\,
\frac{1-F(x)}{xF(x)}
=
\prod_{i=1}^q\frac{b_i}{a_i}\,
\frac{G(x)-1}{x}.
\]
A positive scalar multiple of a Stieltjes function is Stieltjes. Thus \(R\) is Stieltjes and completely monotone.

Finally,
\[
xR(x)=\prod_{i=1}^q\frac{b_i}{a_i}\,(G(x)-1).
\]
Since subtracting the finite value \(G(0)=1\) from a complete Bernstein function leaves the same nonnegative drift and Levy measure part, \(G(x)-1\) is complete Bernstein. Hence \(xR(x)\) is complete Bernstein up to the stated positive scalar.

This is a strong local theorem and a reusable Stieltjes/CBF bridge. The subsequent primary-source wording audit did not find an explicit Karp--Sitnik open problem asking for the CM/Stieltjes/CBF upgrade, so this result is bridge-only unless a separate source asks for exactly this strengthening.
\end{proof}

\begin{lemma}[Baskakov higher power seed r=4 alpha=1 has positive Laplace density and is completely m...]
\label{res:baskakov-r4-alpha1-laplace-density-seed}
For Abel--Gawronski--Neuschel's higher-power Baskakov function, \(f^{[4]}_1(x)=1/((1+x)^4-x^4)=\int_0^\infty e^{-xt}e^{-t/2}(1-\cos(t/2))\,dt\); hence \(f^{[4]}_1\) is completely monotone on \((0,\infty)\).
\end{lemma}
\begin{proof}
For \(\alpha=1\),
\[
\binom{-1}{k}=(-1)^k,
\qquad
\binom{-1}{k}^{4}=1.
\]
Therefore, for \(x>0\),
\[
f^{[4]}_1(x)
=(1+x)^{-4}\sum_{k=0}^{\infty}
\left(\frac{x}{1+x}\right)^{4k}
=\frac{1}{(1+x)^4-x^4}.
\]
Since
\[
(1+x)^4-x^4=(2x+1)(2x^2+2x+1),
\]
partial fractions give
\[
\frac{1}{(1+x)^4-x^4}
=\frac{2}{2x+1}-\frac{2x+1}{2x^2+2x+1}.
\]
Writing \(a=x+\frac12\), this is
\[
\frac{1}{a}-\frac{a}{a^2+(1/2)^2}.
\]
Using
\[
\int_0^\infty e^{-xt}e^{-t/2}\,dt=\frac{1}{x+1/2}
\]
and
\[
\int_0^\infty e^{-xt}e^{-t/2}\cos(t/2)\,dt
=\frac{x+1/2}{(x+1/2)^2+(1/2)^2},
\]
we obtain
\[
f^{[4]}_1(x)
=\int_0^\infty e^{-xt}e^{-t/2}\left(1-\cos\frac{t}{2}\right)\,dt.
\]
The density \(e^{-t/2}(1-\cos(t/2))\) is nonnegative on \([0,\infty)\). By Bernstein's theorem, \(f^{[4]}_1\) is completely monotone on \((0,\infty)\).

From the source definition
\[
p^{[c]}_{n,k}(x)=\binom{-n/c}{k}(-cx)^k(1+cx)^{-n/c-k},
\]
one has
\[
\psi^{[r]}_{n,c}(x)=\sum_{k=0}^{\infty}\left(p^{[c]}_{n,k}(x)\right)^r
=f^{[r]}_{n/c}(cx)
\]
for even \(r\). Hence, for \(c=n\),
\[
\psi^{[4]}_{n,n}(x)=f^{[4]}_1(nx).
\]
Positive scaling preserves complete monotonicity, so \(\psi^{[4]}_{n,n}\) is completely monotone for every positive integer \(n\).
\end{proof}

\begin{lemma}[Bouali f alpha phi alpha endpoint asymptotic]
\label{res:bouali-falpha-phi-asymptotic}
For \(f_\alpha(x)=x^{x(\psi(x)-\log x)-\alpha}\), \(\theta(x)=x(\log x-\psi(x))\), and \(\phi_\alpha(x)=-(\log f_\alpha)'(x)\), one has \(\phi_\alpha(x)=(\alpha+1/2)/x+(1-\log x)/(12x^2)+O(\log x/x^4)\) as \(x\to\infty\).
\end{lemma}
\begin{proof}
Status: endpoint/lower-region obstruction solved; interior gap remains open.

Bouali studies
\[
f_\alpha(x)=x^{x(\psi(x)-\log x)-\alpha}.
\]
The source proves logarithmic complete monotonicity, hence complete monotonicity, for \(\alpha\ge -1/4\). Remark 1.12 prints the open range as \((-1/4,-1/2]\), which is reversed as a literal interval. The local correction treats the intended unresolved lower side as a frontier and proves that the endpoint \(\alpha=-1/2\) cannot be included.

Let \(\theta(x)=x(\log x-\psi(x))\) and
\[
\phi_\alpha(x)=-(\log f_\alpha)'(x)=\frac{\theta(x)+\alpha}{x}+\theta'(x)\log x.
\]
The digamma expansion gives
\[
\phi_\alpha(x)=\frac{\alpha+1/2}{x}+\frac{1-\log x}{12x^2}+O\!\left(\frac{\log x}{x^4}\right).
\]
Thus \(\phi_\alpha(x)<0\) eventually whenever \(\alpha\le -1/2\). Since positive complete monotone functions must be nonincreasing, this proves
\[
\alpha\le -\frac12 \quad\Longrightarrow\quad f_\alpha\notin CM(0,\infty).
\]

This is useful theory growth because it packages a reusable Gamma/psi endpoint-obstruction template, but it is not application-eligible as a full solution. The interior interval
\[
-\frac12<\alpha<-\frac14
\]
remains open.

confirmed the corrected interior interval as source-open and found no later
gap. It isolated the reduction
\[
f_\alpha(x)=x^\beta f_{-1/4}(x),
\qquad
\beta=-(\alpha+1/4)\in(0,1/4),
\]
so the remaining problem is whether the non-CM factor \(x^\beta\) preserves
complete monotonicity for this specific endpoint function. This is useful
theory growth, but not a solved source open problem and not APP-countable.
\end{proof}

\begin{lemma}[Gomilko Tomilov finite Bernstein factor denominator fractional power closure subclass]
\label{res:gt-bf-factorized-denominator-power-closure}
Let \(0<\alpha<1\), let \(\psi\in BF\), and suppose \(x/\psi(x)=\prod_{j=1}^m f_j(x)\), where every \(f_j\) is a positive Bernstein function on \((0,\infty)\). Then \(\psi_\alpha(x)=[\psi(x^\alpha)]^{1/\alpha}\) is a Bernstein function.
\end{lemma}
\begin{proof}
Assume \(0<\alpha<1\), \(\psi\in BF\), and
\[
\frac{x}{\psi(x)}=\prod_{j=1}^m f_j(x),
\]
where every \(f_j\) is a positive Bernstein function on \((0,\infty)\). Then
\[
\psi_\alpha(x)=[\psi(x^\alpha)]^{1/\alpha}
\]
is a Bernstein function.

Put
\[
\beta=\frac1\alpha-1>0.
\]
For \(x>0\), direct differentiation gives
\[
\psi_\alpha'(x)
=x^{\alpha-1}\psi'(x^\alpha)\psi(x^\alpha)^{1/\alpha-1}
=\psi'(x^\alpha)\left(\frac{\psi(x^\alpha)}{x^\alpha}\right)^\beta.
\]
Using the factorization of \(x/\psi(x)\), this becomes
\[
\psi_\alpha'(x)
=\psi'(x^\alpha)\prod_{j=1}^m f_j(x^\alpha)^{-\beta}.
\]

Since \(\psi\in BF\), \(\psi'\) is completely monotone. Since \(x^\alpha\in BF\), composition closure gives \(\psi'(x^\alpha)\in CM\).

For each \(j\), \(g_j(x)=f_j(x^\alpha)\) is Bernstein by BF composition closure. If \(g\) is a positive Bernstein function and \(\beta>0\), then
\[
g(x)^{-\beta}
=\frac1{\Gamma(\beta)}
\int_0^\infty t^{\beta-1}e^{-t g(x)}\,dt.
\]
For each \(t>0\), \(e^{-t g(x)}\) is completely monotone because \(y\mapsto e^{-ty}\) is completely monotone and \(g\in BF\). Hence \(g^{-\beta}\in CM\). Therefore every \(f_j(x^\alpha)^{-\beta}\) is completely monotone.

Finite products of completely monotone functions are completely monotone, so \(\psi_\alpha'\in CM\). Also \(\psi_\alpha\ge0\). Thus \(\psi_\alpha\) is a Bernstein function.

The full source problem remains open locally in the arbitrary non-special case, especially \(1/2<\alpha<1\), when no finite Bernstein-factor representation of \(x/\psi(x)\) is known. The next loop should rotate or look for a genuinely broader structural theorem/counterexample, not repeat this finite-factor subclass.
\end{proof}

\begin{lemma}[KMS arrowhead spectrahedral determinant p minus 2 explicit Riesz kernel]
\label{res:kms-arrowhead-riesz-kernel-pminus2}
For \(p(x,y,z)=x(xy-z^2)\) on \(C=\{x>0,xy>z^2\}\), under pairing \(xu+yv+zw\), \(p^{-2}\) has nonnegative Riesz kernel \(K(u,v,w)=\frac{4\sqrt v}{15\pi}(u-w^2/(4v))^{5/2}\mathbf1_{\{v>0,u>w^2/(4v)\}}\).
\end{lemma}
\begin{proof}
Let
\[
p(x,y,z)=x(xy-z^2),
\qquad
C=\{(x,y,z):x>0,\ xy>z^2\}.
\]
Use the pairing \(xu+yv+zw\). Define
\[
K(u,v,w)=
\frac{4\sqrt v}{15\pi}
\left(u-\frac{w^2}{4v}\right)^{5/2}
\mathbf 1_{\{v>0,\ u>w^2/(4v)\}}.
\]
Then, for \((x,y,z)\in C\),
\[
\int_{\mathbb R^3}e^{-xu-yv-zw}K(u,v,w)\,du\,dv\,dw
=\frac{1}{x^2(xy-z^2)^2}
=p(x,y,z)^{-2}.
\]

Set
\[
s=u-\frac{w^2}{4v}.
\]
Then \(u=s+w^2/(4v)\), \(s>0\), \(v>0\), and the transform equals
\[
\frac{4}{15\pi}\int_0^\infty\int_{-\infty}^{\infty}\int_0^\infty
e^{-xs-xw^2/(4v)-yv-zw}\sqrt v\,s^{5/2}\,ds\,dw\,dv.
\]
The \(s\)-integral is
\[
\int_0^\infty e^{-xs}s^{5/2}\,ds
=\Gamma(7/2)x^{-7/2}
=\frac{15\sqrt\pi}{8}x^{-7/2}.
\]
The \(w\)-integral is
\[
\int_{-\infty}^{\infty}\exp\left(-\frac{xw^2}{4v}-zw\right)\,dw
=2\sqrt{\frac{\pi v}{x}}\exp\left(\frac{vz^2}{x}\right).
\]
Therefore the transform becomes
\[
x^{-4}\int_0^\infty v\exp\left[-v\left(y-\frac{z^2}{x}\right)\right]\,dv.
\]
Since \((x,y,z)\in C\), \(y-z^2/x>0\), and
\[
\int_0^\infty ve^{-av}\,dv=a^{-2}.
\]
Thus
\[
x^{-4}\left(y-\frac{z^2}{x}\right)^{-2}
=\frac{1}{x^2(xy-z^2)^2}
=p(x,y,z)^{-2}.
\]

Identify
\[
(x,y,z)\leftrightarrow
\begin{pmatrix}x&z\\ z&y\end{pmatrix},
\qquad
(u,v,w)\leftrightarrow
\begin{pmatrix}u&w/2\\ w/2&v\end{pmatrix}.
\]
Then \(\operatorname{tr}(AB)=xu+yv+zw\). Hence the dual cone is
\[
C^*=\{u\ge0,\ v\ge0,\ 4uv\ge w^2\}.
\]
The formula uses the coordinate \(w=2B_{12}\). In the raw symmetric-matrix coordinate \(s=B_{12}\), the density becomes
\[
\frac{8\sqrt v}{15\pi}
\left(u-\frac{s^2}{v}\right)^{5/2}
\mathbf 1_{\{v>0,\ u>s^2/v\}}.
\]
\end{proof}

\begin{counterexample}[finite exponential mixture has inverse branch not completely monotone]
\label{res:keady-cm-inverse-not-cm-example}
The strictly completely monotone function \(f(x)=99e^{-x}+e^{-10x}\) maps \((0,\infty)\) decreasingly onto \((0,100)\), but its inverse branch is not completely monotone on \((0,100)\).
\end{counterexample}
\begin{proof}
There exists a strictly completely monotone decreasing function whose inverse branch is not completely monotone on its natural interval. An explicit example is
\[
f(x)=99e^{-x}+e^{-10x},\qquad x>0.
\]
It maps \((0,\infty)\) decreasingly onto \((0,100)\). Its inverse branch \(g=f^{-1}:(0,100)\to(0,\infty)\) is not completely monotone.

For every \(n\ge0\),
\[
(-1)^n f^{(n)}(x)=99e^{-x}+10^n e^{-10x}>0.
\]
Hence \(f\) is strictly completely monotone, and \(f'<0\), so the inverse branch \(g=f^{-1}\) exists.

For \(y=f(x)\), inverse differentiation gives
\[
g'(y)=\frac1{f'(x)},\qquad
g''(y)=-\frac{f''(x)}{(f'(x))^3},
\]
and
\[
g'''(y)=\frac{3(f''(x))^2-f'''(x)f'(x)}{(f'(x))^5}.
\]
At \(x=0+\),
\[
f'(0+)=-109,\qquad f''(0+)=199,\qquad f'''(0+)=-1099.
\]
Therefore
\[
3(f''(0+))^2-f'''(0+)f'(0+)
=3\cdot199^2-(-1099)(-109)
=-988<0.
\]
Since \((f'(0+))^5<0\), this gives
\[
g'''(100-)>0.
\]
By continuity, \(g'''>0\) on a left-neighborhood of \(100\). Complete monotonicity of \(g\) would require \((-1)^3g'''\ge0\), i.e. \(g'''\le0\). This contradiction proves that \(g\) is not completely monotone.
\end{proof}

\begin{counterexample}[Szabo psi difference endpoint alpha=2 is not completely monotone for all 0<d<1]
\label{res:szabo-psi-difference-alpha2-not-cm}
For every \(0<d<1\), the function \(y\mapsto y^2\left[\psi(y+d)-\psi(y)-d/y\right]\) is not completely monotone on \((0,\infty)\).
\end{counterexample}
\begin{proof}
Let \(0<d<1\) and
\[
\phi_d(y)=\psi(y+d)-\psi(y)-\frac{d}{y},\qquad y>0.
\]
Then \(y^2\phi_d(y)\) is not completely monotone on \((0,\infty)\). Equivalently, for Szabo's original variables \(d=b-a\) and \(y=x+a\), the endpoint
\[
\alpha=2
\]
is not admissible in Open Problem 1.5 for any \(0<b-a<1\).

The standard digamma difference formula gives
\[
\psi(y+d)-\psi(y)=\int_0^\infty e^{-yt}\frac{1-e^{-dt}}{1-e^{-t}}\,dt
\]
and
\[
\frac{d}{y}=\int_0^\infty e^{-yt}d\,dt.
\]
Thus
\[
\phi_d(y)=\int_0^\infty e^{-yt}k_d(t)\,dt,\qquad
k_d(t)=\frac{1-e^{-dt}}{1-e^{-t}}-d.
\]
At the origin,
\[
k_d(t)=\frac{d(1-d)}2t+O(t^2),
\]
so \(k_d(0+)=0\) and \(k_d'(0+)=d(1-d)/2>0\).

In the sense of Laplace transforms of distributions,
\[
y^2\phi_d(y)
=\frac{d(1-d)}2+\int_0^\infty e^{-yt}k_d''(t)\,dt.
\]
Indeed, integrating by parts twice gives
\[
\mathcal L(k_d'')(y)=y^2\mathcal L(k_d)(y)-k_d'(0+),
\]
because \(k_d(0+)=0\).

As \(t\to\infty\),
\[
k_d(t)=1-d-e^{-dt}+O(e^{-t}),
\]
and hence
\[
k_d''(t)=-d^2e^{-dt}+O(e^{-t}).
\]
Since \(0<d<1\), the term \(e^{-dt}\) dominates \(e^{-t}\), so \(k_d''(t)<0\) for all sufficiently large \(t\).

If \(y^2\phi_d(y)\) were completely monotone on \((0,\infty)\), the Bernstein-Widder theorem would give a nonnegative representing measure. The inverse Laplace distribution computed above has an absolutely continuous tail with negative density, contradicting uniqueness of the Laplace transform. Therefore \(y^2\phi_d(y)\) is not completely monotone.

The shift \(y=x+a\) preserves the obstruction for the original function on \((-a,\infty)\): the representing distribution is multiplied by \(e^{-at}\), which does not change the eventual negative sign.
\end{proof}

\begin{lemma}[sqrt(c s power a+d s power b) two-term concavity loss criterion outside unit disk centered at (1,1)]
\label{res:sqrt-two-term-symbol-concavity-loss-criterion}
For \(g(s)=c s^a+d s^b\), \(c,d>0\), \(1<a<2\), and \(0<b<a-1\), the function \(h(s)=\sqrt{g(s)}\) has \(h''(s)>0\) somewhere if and only if \((a-1)^2+(b-1)^2>1\).
\end{lemma}
\begin{proof}
Primary source: Emilia Bazhlekova and Ivan Bazhlekov, "Subordination approach to multi-term time-fractional diffusion-wave equations", arXiv:1707.09828.

The source proves the propagation-function positivity package under
\[
1<\alpha\le2,\qquad \alpha-\alpha_m\le1,
\]
by showing that
\[
\sqrt{g(s)}\in CBF,\qquad g(s)=c s^\alpha+\sum_j c_j s^{\alpha_j}.
\]
It leaves open whether, and to what extent, this gap condition can be relaxed for the actual properties
\[
w\ge0,\qquad w_t\ge0,\qquad -w_x\ge0.
\]

For the two-term wave endpoint
\[
g(s)=c s^2+d s^b,\qquad c,d>0,\quad 0<b<1,
\]
the condition cannot be relaxed. Put \(h(s)=\sqrt{g(s)}\). Then
\[
h''(s)=\frac{d}{2\sqrt c}(b-1)(b-2)s^{b-3}+O(s^{2b-5})>0
\]
for all sufficiently large \(s\).

The source gives
\[
\mathcal L\{w_t(x,\cdot)\}(s)=e^{-xh(s)}.
\]
If \(w_t(x,\cdot)\ge0\), this Laplace transform must be completely monotone. But
\[
\frac{d^2}{ds^2}e^{-xh(s)}
=e^{-xh(s)}\left(x^2h'(s)^2-xh''(s)\right),
\]
which is negative at a large \(s_0\) after choosing \(x>0\) sufficiently small. Hence \(w_t\) fails positivity for some \(x\).

\[
g(s)=c s^a+d s^b,\qquad c,d>0,\quad 1<a<2,\quad 0<b<a-1,
\]
and set \(h(s)=\sqrt{g(s)}\), \(y=(c/d)s^{a-b}\). Then the sign of \(h''\) is the sign of
\[
N_{a,b}(y)
=a(a-2)y^2
+2(a^2-ab-a+b^2-b)y
+b(b-2).
\]
The discriminant is
\[
4(a-b)^2\left((a-1)^2+(b-1)^2-1\right).
\]
Since \(a>1\) and \(b<1\), \(h''\) is positive somewhere exactly when
\[
(a-1)^2+(b-1)^2>1.
\]
In that region, \(e^{-x\sqrt g}\) fails complete monotonicity for a suitable small \(x>0\), so \(w_t\) cannot be nonnegative for all \(x,t>0\). The source examples \((a,b)=(1.9,0.5)\) and \((1.8,0.3)\) lie in this outside-disk region.

The remaining two-term inner gap
\[
1<a<2,\qquad 0<b<a-1,\qquad (a-1)^2+(b-1)^2\le1
\]
is still open locally. There the square-root symbol is concave by this test, but no Bernstein-function representation or positivity theorem has been admitted.

\[
g(s)=s^{28/25}+s^{1/50},\qquad h(s)=\sqrt{g(s)}.
\]
Then
\[
1<a<2,\qquad 0<b<a-1,\qquad a-b=\frac{11}{10}>1,
\]
and
\[
(a-1)^2+(b-1)^2
=\left(\frac3{25}\right)^2+\left(-\frac{49}{50}\right)^2
=\frac{2437}{2500}<1.
\]
Thus this is not covered by the positive-\(h''\) outside-disk obstruction.

Direct differentiation gives
\[
h^{(5)}(1)
=-\frac{5570045943\sqrt2}{320000000000}<0.
\]
Since a Bernstein function has completely monotone derivative, it must satisfy \(h^{(5)}\ge0\). Therefore \(\sqrt{s^{28/25}+s^{1/50}}\) is not a Bernstein function.

The same example gives propagation failure, not only failure of the sufficient CBF route. With
\[
F_x(s)=e^{-xh(s)}=\mathcal L\{w_t(x,\cdot)\}(s),
\]
we have
\[
F_x^{(5)}(1)=-x h^{(5)}(1)+O(x^2)>0
\]
for all sufficiently small \(x>0\). This contradicts complete monotonicity of \(F_x\), which would be necessary if \(w_t(x,\cdot)\ge0\). Hence \(w_t\)-positivity fails for this inner-gap two-term symbol.

This is still a partial source answer, not a complete classification of the inner disk. The remaining open problem is now a residual parameter-region classification for Bernstein status and propagation positivity.
\end{proof}

\begin{lemma}[Baskakov alpha 1 even r rational seed line complete monotonicity frontier]
\label{res:baskakov-alpha1-even-r-frontier-open}
Determine whether \(f^{[2m]}_1(x)=1/((1+x)^{2m}-x^{2m})\) is completely monotone on \((0,\infty)\) for every integer \(m\ge2\).
\end{lemma}
\begin{proof}
For \(\alpha=1\),
\[
\binom{-1}{k}=(-1)^k,
\qquad
\binom{-1}{k}^{4}=1.
\]
Therefore, for \(x>0\),
\[
f^{[4]}_1(x)
=(1+x)^{-4}\sum_{k=0}^{\infty}
\left(\frac{x}{1+x}\right)^{4k}
=\frac{1}{(1+x)^4-x^4}.
\]
Since
\[
(1+x)^4-x^4=(2x+1)(2x^2+2x+1),
\]
partial fractions give
\[
\frac{1}{(1+x)^4-x^4}
=\frac{2}{2x+1}-\frac{2x+1}{2x^2+2x+1}.
\]
Writing \(a=x+\frac12\), this is
\[
\frac{1}{a}-\frac{a}{a^2+(1/2)^2}.
\]
Using
\[
\int_0^\infty e^{-xt}e^{-t/2}\,dt=\frac{1}{x+1/2}
\]
and
\[
\int_0^\infty e^{-xt}e^{-t/2}\cos(t/2)\,dt
=\frac{x+1/2}{(x+1/2)^2+(1/2)^2},
\]
we obtain
\[
f^{[4]}_1(x)
=\int_0^\infty e^{-xt}e^{-t/2}\left(1-\cos\frac{t}{2}\right)\,dt.
\]
The density \(e^{-t/2}(1-\cos(t/2))\) is nonnegative on \([0,\infty)\). By Bernstein's theorem, \(f^{[4]}_1\) is completely monotone on \((0,\infty)\).

From the source definition
\[
p^{[c]}_{n,k}(x)=\binom{-n/c}{k}(-cx)^k(1+cx)^{-n/c-k},
\]
one has
\[
\psi^{[r]}_{n,c}(x)=\sum_{k=0}^{\infty}\left(p^{[c]}_{n,k}(x)\right)^r
=f^{[r]}_{n/c}(cx)
\]
for even \(r\). Hence, for \(c=n\),
\[
\psi^{[4]}_{n,n}(x)=f^{[4]}_1(nx).
\]
Positive scaling preserves complete monotonicity, so \(\psi^{[4]}_{n,n}\) is completely monotone for every positive integer \(n\).
\end{proof}

\begin{lemma}[Bazhlekova Wright density topcap bridge normal form]
\label{res:bazhlekova-wright-density-topcap-bridge-normal-form}
For the normalized Bazhlekova symbol \(h(s)=s^\alpha(1+s^{-p})^{1/2}\), with \(0<\beta=\alpha-p/2<\alpha<1\) and the principal branch in the no-cover seed range, define \(\mathcal W_{\alpha,p}(x)=-\sum_{m\ge0}\binom{1/2}{m}x^m/\Gamma(pm-\alpha)\). Then \(\mathcal L^{-1}\{h'\}(t)=t^{-\alpha}\mathcal W_{\alpha,p}(t^p)\), and the Wright top-cap limit has the same sign as \(\mathcal W_{\alpha,p}(\lambda^{-1})\). Thus normalized derivative density positivity and top-cap positivity are the same Wright-function sign problem.
\end{lemma}
\begin{proof}
Put
\[
\alpha=\frac a2,\qquad p=a-b,\qquad \beta=\frac b2=\alpha-\frac p2,
\]
so that
\[
h(s)=s^{b/2}(1+s^{a-b})^{1/2}
=s^\alpha(1+s^{-p})^{1/2}.
\]
At the two no-cover seeds,
\[
(\alpha,p)=\left(\frac34,\frac{11}{10}\right),
\qquad
(\alpha,p)=\left(\frac{11}{20},\frac{21}{20}\right),
\]
and \(0<\beta<\alpha<1\), \(1<p<2\).

Define the Wright-normalized series
\[
\mathcal W_{\alpha,p}(x)
=
-\sum_{m=0}^{\infty}
\binom{1/2}{m}\frac{x^m}{\Gamma(pm-\alpha)}.
\]
The coefficient ratio and the gamma denominator show that this series is entire in \(x\).

For \(s>1\), the binomial expansion gives
\[
h'(s)
=
\sum_{m=0}^{\infty}
\binom{1/2}{m}(\alpha-pm)s^{\alpha-pm-1}.
\]
For each \(m\), analytic continuation of the elementary Laplace formula gives
\[
\mathcal L\left\{
\frac{t^{pm-\alpha}}{\Gamma(pm-\alpha)}
\right\}(s)
=(pm-\alpha)s^{\alpha-pm-1}.
\]
Therefore
\[
\mathcal L^{-1}\{h'\}(t)
=
-\sum_{m=0}^{\infty}
\binom{1/2}{m}
\frac{t^{pm-\alpha}}{\Gamma(pm-\alpha)}
=
t^{-\alpha}\mathcal W_{\alpha,p}(t^p).
\]
For \(1<p<2\) the Wright series has subexponential growth, so the Laplace transform is defined on \(s>0\). Equality for \(s>1\) extends to \(s>0\) by analytic continuation on the principal branch.

Consequently, at the no-cover seeds, complete monotonicity of \(h'\) is equivalent to
\[
\mathcal W_{\alpha,p}(x)\ge0
\qquad(x>0),
\]
because the inverse Laplace transform is the locally integrable density above.

The previous top-cap audit used the signed polynomials \(R_n(y)=(-1)^{n-1}Q_n(y)\). Its formal top-cap function was
\[
W^{\mathrm{top}}_{\alpha,p}(\lambda)
=
1-\frac{\Gamma(1-\alpha)}{\alpha}
\sum_{m\ge1}\binom{1/2}{m}
\frac{\lambda^{-m}}{\Gamma(pm-\alpha)}.
\]
Since
\[
-\frac1{\Gamma(-\alpha)}
=\frac{\alpha}{\Gamma(1-\alpha)},
\]
the two normalizations satisfy
\[
W^{\mathrm{top}}_{\alpha,p}(\lambda)
=
\frac{\Gamma(1-\alpha)}{\alpha}
\mathcal W_{\alpha,p}(\lambda^{-1}).
\]
The factor \(\Gamma(1-\alpha)/\alpha\) is positive. Thus top-cap positivity and normalized derivative density positivity are the same sign problem after \(x=\lambda^{-1}\).
\end{proof}

The following appendix result records the fixed source problem \textbf{APP-0037} without interrupting the main exposition.
\begin{corollary}[APP-0037: Du-Wang \(h\_3\) is increasing exactly on \(\frac12\le a\le1\)]
\label{res:du-wang-h3-open-problem-2-classification}
For \(0<a<2\), Du--Wang's function \(h_3\) is increasing on \((0,\infty)\) exactly for \(1/2\le a\le1\), and is not monotone for \(0<a<1/2\) or \(1<a<2\).

This resolves the fixed application label \textbf{APP-0037}: Classify Du-Wang \(h\_3\) monotonicity on \((0,\infty)\).
\emph{Source reference.} \cite{app-app-0037-source}.
\end{corollary}
\begin{proof}
For \(t>1/2\), define
\[
S_m(t)=\sum_{n=0}^{\infty}\frac1{(n+t)^m}.
\]
The polygamma identities give
\[
\psi''(t)=-2S_3(t),\qquad \psi'''(t)=6S_4(t).
\]
Thus the desired bound
\[
-\frac{2\psi''(t)}{\psi'''(t)}\ge t-\frac12
\]
is equivalent to
\[
2S_3(t)-3\left(t-\frac12\right)S_4(t)\ge0.
\]
Write \(y=t-\frac12>0\).  Using
\[
S_m(t)=\frac1{\Gamma(m)}
\int_0^\infty \frac{x^{m-1}e^{-tx}}{1-e^{-x}}\,dx
\]
and
\[
\frac{e^{-tx}}{1-e^{-x}}
=e^{-yx}\frac1{2\sinh(x/2)},
\]
set
\[
K(x)=\frac1{2\sinh(x/2)}.
\]
Then
\[
2S_3(t)-3yS_4(t)
=
\int_0^\infty e^{-yx}x^2K(x)\,dx
-\frac y2\int_0^\infty e^{-yx}x^3K(x)\,dx.
\]
Since \(x^3K(x)\sim x^2\) as \(x\downarrow0\) and decays exponentially at
infinity, integration by parts gives
\[
y\int_0^\infty e^{-yx}x^3K(x)\,dx
=
\int_0^\infty e^{-yx}(x^3K(x))'\,dx.
\]
Therefore
\[
2S_3(t)-3yS_4(t)
=
\int_0^\infty e^{-yx}
\left(x^2K(x)-\frac12(x^3K(x))'\right)dx.
\]
Now
\[
K'(x)=-\frac12K(x)\coth(x/2),
\]
so
\[
x^2K(x)-\frac12(x^3K(x))'
=
\frac12x^2K(x)\left(\frac x2\coth(x/2)-1\right).
\]
For \(z>0\), \(z\coth z>1\).  The integrand is therefore strictly positive,
and
\[
2S_3(t)-3\left(t-\frac12\right)S_4(t)>0.
\]
This proves
\[
-\frac{2\psi''(t)}{\psi'''(t)}>t-\frac12,
\qquad t>\frac12.
\]

This proves the claim.

Let \(1/2\le a\le1\), \(u>0\), and \(t=a+u\).  Since \(a\ge1/2\),
\[
u=t-a\le t-\frac12.
\]
The ratio lemma gives
\[
t-\frac12<-\frac{2\psi''(t)}{\psi'''(t)}.
\]
Because \(\psi'''(t)>0\),
\[
2\psi''(t)+u\psi'''(t)<0.
\]
Consequently
\[
H_a'(u)=-u^2\{2\psi''(t)+u\psi'''(t)\}>0.
\]
Together with \(H_a(0+)=2\log\Gamma(a)\ge0\), this gives
\[
H_a(u)\ge0
\qquad (u>0,\ 1/2\le a\le1).
\]

This proves the claim.

For \(x>0\), \(u=x\), and \(t=x+a\),
\[
h_3'(x)=\widetilde h_3'(x+a)=\frac{h_{31}(x+a)}{x^2}
=\frac{H_a(x)}{x^2}\ge0.
\]
The inequality is strict for \(x>0\) except possibly at the endpoint
\(a=1,u=0\), which is outside the open \(x\)-interval.  Thus \(h_3\) is
increasing on \((0,\infty)\) for every
\[
\frac12\le a\le1.
\]

This proves the claim.
the corresponding result.

The earlier the preceding theorem proves that
\(h_3\) is not monotone for
\[
0<a<\frac12
\qquad\text{or}\qquad
1<a<2.
\]
The middle-window theorem above proves increasing behavior for
\[
\frac12\le a\le1.
\]
Therefore Du--Wang Open Problem 2 is classified on \(0<a<2\):
\[
h_3 \text{ is increasing on }(0,\infty)
\quad\Longleftrightarrow\quad
\frac12\le a\le1,
\]
and it is not monotone on the two outer windows.

This proves the claim.
the corresponding result true by source-problem classification.
\end{proof}

The following appendix result records the fixed source problem \textbf{APP-0018} without interrupting the main exposition.
\begin{corollary}[APP-0018: Qi-Guo tau threshold has a supremum-corrected exact value]
\label{res:qg-tau-global-supremum}
Let \(a_*>0\) be the unique positive root of \(e^{a_*}=1+a_*+a_*^2\). Then \(\sup_{s\in\mathbb N,t>0}\tau(s,t)=a_*/(1+a_*+a_*^2)=0.298425607525639\ldots\), and the supremum is not attained.

This resolves the fixed application label \textbf{APP-0018}: Find the optimal uniform threshold for the tau expression in Open Problem 3.
\emph{Source reference.} \cite{app-app-0018-source}.
\end{corollary}
\begin{proof}
Writing \(u=t/(t+1)\) gives a two-parameter expression on \(0<u<1\).  The
Euler-Lagrange equations reduce the interior extremal sequence to the scaling
limit \(u=e^{-a/s}\).  The limiting profile is maximized exactly when
\(e^a=1+a+a^2\), and the resulting value is \(a/(1+a+a^2)\).  The source
one-variable monotonicity estimates bound the finite-\(s\) expressions by this
limit, while the scaling sequence approaches it.  Hence the number is a
supremum and is not attained at finite \(s,t\).
\end{proof}

The following appendix result records the fixed source problem \textbf{APP-0023} without interrupting the main exposition.
\begin{corollary}[APP-0023: Sokal generalized-Stieltjes lambda-derivatives have a triangular nonnegative compression]
\label{res:sokal-gs-lambda-derivative-generating-function}
For fixed \(n\), Sokal's generalized-Stieltjes tests satisfy the formal exponential generating function \(\sum_{k\ge0}F^{[\lambda]}_{n,k}(x)z^k/k!=(1-z)^{-(n+\lambda)}(-1)^n\sum_{j\ge0}x^j f^{(n+j)}(x)(z/(1-z))^j/j!\).

This resolves the fixed application label \textbf{APP-0023}: Explain the lambda-derivative structure of Sokal's generalized-Stieltjes Hankel-type expressions.
\emph{Source reference.} \cite{app-app-0023-source}.
\end{corollary}
\begin{proof}
The exponential generating function is
\[
 \mathcal F(z)
 =(1-z)^{-a}\sum_{j\ge0}G_j(x)
      \left({z\over 1-z}\right)^j{1\over j!}.
\]
Differentiating \(\ell\) times with respect to \(\lambda\) multiplies
\(\mathcal F\) by \((-\log(1-z))^\ell\).  Extracting coefficients gives the
displayed triangular formula.  Since
\(-\log(1-z)=\sum_{m\ge1}z^m/m\) has non-negative coefficients, all compression
coefficients are non-negative.
\end{proof}

\begin{lemma}[generalized Bernstein classes closed under product with order addition]
\label{res:gbf-product-closure}
If \(\lambda_1,\lambda_2>0\), \(f_1\in\mathcal B_{\lambda_1}\), and \(f_2\in\mathcal B_{\lambda_2}\), then \(f_1f_2\in\mathcal B_{\lambda_1+\lambda_2}\).
\end{lemma}
\begin{proof}
For \(\lambda>0\), the source defines \(f\in\mathcal B_\lambda\) when
\[
x^{1-\lambda}f'(x)
\]
is completely monotone. It also recalls generalized Stieltjes classes \(\mathcal S_\lambda\) and the product closure
\[
\mathcal S_{\lambda_1}\mathcal S_{\lambda_2}\subseteq \mathcal S_{\lambda_1+\lambda_2}.
\]

Proposition 5.1 states that if
\[
f_1\in\mathcal B_{\lambda_1},\qquad f_2\in\mathcal B_{\lambda_2},
\]
then
\[
f_1f_2\in\mathcal B_{\lambda_1+\lambda_2}.
\]

Local proof audit: by definition \(f_j'(x)/x^{\lambda_j-1}\) is completely monotone, and by the cited Corollary 2.1, \(f_j(x)/x^{\lambda_j}\) is completely monotone. Then
\[
\frac{(f_1f_2)'(x)}{x^{\lambda_1+\lambda_2-1}}
=
\frac{f_1'(x)}{x^{\lambda_1-1}}\frac{f_2(x)}{x^{\lambda_2}}
+
\frac{f_1(x)}{x^{\lambda_1}}\frac{f_2'(x)}{x^{\lambda_2-1}}.
\]
The right side is a positive sum of products of completely monotone functions, hence completely monotone.

The same source's Proposition 5.11 defines a class \(\mathcal C_\alpha\): \(f\in\mathcal C_\alpha\) when \(x^\alpha f(x)\) is the Laplace transform of a decreasing logarithmically convex function. It proves
\[
f\in\mathcal C_\alpha
\quad\Longrightarrow\quad
\frac1f\in\mathcal B_{\alpha+1}
\]
and
\[
\frac{1}{x^{\alpha+1}f(x)}
\]
is completely monotone.

Promote bridge nodes:

Demote/quarantine:
\end{proof}

The following appendix result records the fixed source problem \textbf{APP-0036} without interrupting the main exposition.
\begin{counterexample}[APP-0036: Baskakov even-line complete-monotonicity fails at \(r=8\)]
\label{res:not-baskakov-higher-power-even-conjecture}
Abel--Gawronski--Neuschel's even-power Baskakov complete-monotonicity conjecture is false; it fails at \(\alpha=1\), \(r=8\).

This resolves the fixed application label \textbf{APP-0036}: Decide the even-line Baskakov complete-monotonicity conjecture at \(r=8\).
\emph{Source reference.} \cite{app-app-0036-source}.
\end{counterexample}
\begin{proof}
Set
\[
D_m(x)=(1+x)^{2m}-x^{2m}.
\]
The roots are obtained from
\[
\frac{1+x}{x}=\omega_k,\qquad
\omega_k=e^{\pi i k/m},\qquad k=1,\ldots,2m-1.
\]
Thus
\[
\lambda_{k,m}
=\frac{1}{\omega_k-1}
=-\frac12-\frac{i}{2}\cot\frac{\pi k}{2m}.
\]
At this root,
\[
D_m'(\lambda_{k,m})
=2m\lambda_{k,m}^{2m-1}(\omega_k^{-1}-1),
\]
and direct simplification gives the real residue
\[
R_{k,m}
=\frac{1}{D_m'(\lambda_{k,m})}
=\frac{(-1)^{m+k}}{2m}
\left(2\sin\frac{\pi k}{2m}\right)^{2m-2}.
\]
Consequently the inverse-Laplace density is
\[
\rho_m(t)
=\sum_{k=1}^{2m-1}R_{k,m}e^{\lambda_{k,m}t}.
\]
Pairing conjugate roots gives
\[
\rho_m(t)
=\frac{e^{-t/2}}{2m}
\left[
2^{2m-2}
+2\sum_{k=1}^{m-1}
(-1)^{m+k}
\left(2\sin\frac{\pi k}{2m}\right)^{2m-2}
\cos\!\left(
\frac{t}{2}\cot\frac{\pi k}{2m}
\right)
\right].
\]

This proves the finite trigonometric density formula used by the diagnostic
candidate.

For \(m=2\), this gives
\[
\rho_2(t)=e^{-t/2}\left(1-\cos\frac t2\right)
=2e^{-t/2}\sin^2\frac t4\ge0,
\]
recovering the previously admitted \(r=4,\alpha=1\) seed.

For \(m=3\),
\[
\rho_3(t)
=\frac{e^{-t/2}}6
\left[
16+2\cos\left(\frac{\sqrt3\,t}{2}\right)
-18\cos\left(\frac{t}{2\sqrt3}\right)
\right].
\]
Set \(u=t/(2\sqrt3)\).  Since
\(\cos(3u)=4\cos^3u-3\cos u\),
\[
\rho_3(t)
=\frac{4}{3}e^{-t/2}(1-\cos u)^2(2+\cos u)\ge0.
\]
Thus \(r=6,\alpha=1\) is completely monotone by an explicit positive density.

For \(m=4\), write \(h_4(t)=e^{t/2}\rho_4(t)\).  The density formula gives
\[
h_4(t)
=8+2\cos\frac t2
-\frac{(2-\sqrt2)^3}{4}
\cos\left(\frac{1+\sqrt2}{2}t\right)
-\frac{(2+\sqrt2)^3}{4}
\cos\left(\frac{\sqrt2-1}{2}t\right).
\]
At \(t=10\pi\),
\[
\cos(t/2)=\cos(5\pi)=-1,
\]
and both other cosine terms equal \(-\cos(5\pi\sqrt2)\).  Since
\[
(2-\sqrt2)^3+(2+\sqrt2)^3=40,
\]
we obtain
\[
h_4(10\pi)=6+10\cos(5\pi\sqrt2).
\]
Now
\[
0<5\sqrt2-7<\frac14
\]
because \(7/5<\sqrt2<29/20\).  Hence, with
\(\delta=5\sqrt2-7\),
\[
\cos(5\pi\sqrt2)=\cos(7\pi+\pi\delta)=-\cos(\pi\delta)
<-\cos\frac\pi4=-\frac{\sqrt2}{2}<-\frac35.
\]
Therefore
\[
h_4(10\pi)<6+10\left(-\frac35\right)=0,
\]
and
\[
\rho_4(10\pi)=e^{-5\pi}h_4(10\pi)<0.
\]
Since the inverse-Laplace transform is unique for measures on
\([0,\infty)\), \(f^{[8]}_1(x)=1/((1+x)^8-x^8)\) cannot be completely monotone.
\end{proof}

The following appendix result records the fixed source problem \textbf{APP-0020} without interrupting the main exposition.
\begin{counterexample}[APP-0020: Qi's degree-four conjecture for h\_lambda is refuted at the endpoint]
\label{res:not-qi-hlambda-degree4-conjecture}
not(For \(\Psi(x)=[\psi'(x)]^2+\psi''(x)\) and \(h_\lambda(x)=\Psi(x)-(x^2+\lambda x+12)/(12x^4(x+1)^2)\), \(\deg_x^{\rm cm}h_\lambda=\deg_x^{\rm cm}(-h_\mu)=4\) if and only if \(\lambda\le0\) and \(\mu\ge4\).)

This resolves the fixed application label \textbf{APP-0020}: Test the conjectured complete-monotonic degree four for \(h\_\lambda\) and \(-h\_\mu\).
\emph{Source reference.} \cite{app-app-0020-source}.
\end{counterexample}
\begin{proof}
Near zero,
\[
 h_\lambda(x)=-{\lambda\over 12x^3}
 +\left({\lambda\over6}+{\pi^2\over3}-{37\over12}\right){1\over x^2}
 +O(x^{-1}).
\]
If \(\lambda<0\), then \(x^4h_\lambda(x)=(-\lambda/12)x+O(x^2)\), whose
derivative is positive near zero.  If \(\lambda=0\), the coefficient
\(\pi^2/3-37/12\) is positive, again giving a positive derivative for the
degree-four transform near zero.  Complete monotonicity would require the first
derivative to be non-positive.  The same small-\(x\) test applied to
\(-x^4h_\mu\) excludes the claimed \(\mu\)-range.  Hence the degree-four
conjecture fails.
\end{proof}

The following appendix result records the fixed source problem \textbf{APP-0013} without interrupting the main exposition.
\begin{corollary}[APP-0013: Yin \((p,k)\)-digamma sharp \(\alpha\)-necessity]
\label{res:yin-pk-digamma-alpha-necessity}
For \(p\in\mathbb N\), \(k>0\), if \(x^\alpha\left[\frac1k\log\frac{pkx}{x+k(p+1)}-\psi_{p,k}(x)\right]\) is completely monotone on \((0,\infty)\), then \(\alpha\le1\).

This resolves the fixed application label \textbf{APP-0013}: Yin \((p,k)\)-digamma sharp \(\alpha\)-necessity
\emph{Source reference.} \cite{app-app-0013-source}.
\end{corollary}
\begin{proof}
\emph{Setup.}
Fix \(p\in\mathbb N\) and \(k>0\). Define
\[
B_{p,k}(x)=
\frac1k\log\frac{pkx}{x+k(p+1)}-\psi_{p,k}(x),
\qquad
\delta_{p,k,\alpha}(x)=x^\alpha B_{p,k}(x).
\]

The source gives the finite formula
\[
\psi_{p,k}(x)=\frac1k\log(pk)-\sum_{n=0}^p\frac1{x+nk}.
\]
Therefore
\[
\begin{aligned}
B_{p,k}(x)
&=
\frac1k\log\frac{pkx}{x+k(p+1)}
-\frac1k\log(pk)
+\sum_{n=0}^p\frac1{x+nk}  \\
&=
\frac1k\log\frac{x}{x+k(p+1)}
+\sum_{n=0}^p\frac1{x+nk}.
\end{aligned}
\]

This normal form also gives positivity directly. Put \(t=x/k\). Then
\[
B_{p,k}(x)
=\frac1k\left[
\sum_{n=0}^p\frac1{t+n}
+\log\frac{t}{t+p+1}
\right].
\]
Since
\[
\log\frac{t+p+1}{t}=\int_0^{p+1}\frac{du}{t+u},
\]
and \(u\mapsto(t+u)^{-1}\) is strictly decreasing, the left Riemann sum dominates:
\[
\sum_{n=0}^p\frac1{t+n}
>
\int_0^{p+1}\frac{du}{t+u}.
\]
Hence \(B_{p,k}(x)>0\) for \(x>0\).

This proves the claim.

From the finite normal form,
\[
B_{p,k}(x)
=\frac1x
+\frac1k\log\frac{x}{x+k(p+1)}
+\sum_{n=1}^p\frac1{x+nk}.
\]

As \(x\to0^+\), the finite sum over \(n\ge1\) is \(O(1)\), and
\[
\frac1k\log\frac{x}{x+k(p+1)}=O(|\log x|).
\]
Thus
\[
B_{p,k}(x)=\frac1x+O(|\log x|)
\qquad (x\to0^+).
\]
Equivalently,
\[
xB_{p,k}(x)\to1.
\]
Therefore
\[
\delta_{p,k,\alpha}(x)
=x^\alpha B_{p,k}(x)
=x^{\alpha-1}(xB_{p,k}(x))
\sim x^{\alpha-1}.
\]

This proves the claim.

Assume that \(\delta_{p,k,\alpha}\) is completely monotone on \((0,\infty)\). Then
\[
\delta_{p,k,\alpha}(x)\ge0,\qquad
\delta_{p,k,\alpha}'(x)\le0.
\]
Since \(B_{p,k}(x)>0\), the function \(\delta_{p,k,\alpha}\) is actually positive for every \(x>0\).

Suppose, for contradiction, that \(\alpha>1\). The endpoint asymptotic gives
\[
\lim_{x\to0^+}\delta_{p,k,\alpha}(x)=0.
\]
Fix \(x_0>0\). Positivity gives \(\delta_{p,k,\alpha}(x_0)>0\). Because \(\delta_{p,k,\alpha}(x)\to0\) as \(x\to0^+\), choose \(0<y<x_0\) with
\[
\delta_{p,k,\alpha}(y)<\frac12\delta_{p,k,\alpha}(x_0).
\]
But \(\delta_{p,k,\alpha}'\le0\) means the function is nonincreasing as \(x\) increases, hence for \(0<y<x_0\),
\[
\delta_{p,k,\alpha}(y)\ge\delta_{p,k,\alpha}(x_0),
\]
a contradiction.

Therefore \(\alpha\le1\).

This proves the claim.
\end{proof}

The following appendix result records the fixed source problem \textbf{APP-0039} without interrupting the main exposition.
\begin{counterexample}[APP-0039: Ramanujan Turan-window complete-monotonicity interval is false]
\label{res:ramanujan-turan-window-negative-answer}
The Mishra--Swaminathan Ramanujan integral Turan-window complete-monotonicity problem has a negative answer: the proposed interval already fails for \(n=2\) and suitable \(\alpha\in(0,1/2)\).

This resolves the fixed application label \textbf{APP-0039}: Decide the Ramanujan Turan-window complete-monotonicity interval.
\emph{Source reference.} \cite{app-app-0039-source}.
\end{counterexample}
\begin{proof}
Put
\[
A_k(x)=(-1)^k I_R^{(k)}(x)
=\int_0^\infty e^{-xt}\frac{t^{k-1}}{\pi^2+\log^2 t}\,dt
\qquad (k\ge1).
\]
Then
\[
H_2(x;\alpha)=A_2(x)^2-\alpha A_1(x)A_3(x).
\]

We study \(x=e^{-L}\) as \(L\to\infty\). With \(y=xt\),
\[
A_k(e^{-L})
=
e^{kL}\int_0^\infty
\frac{e^{-y}y^{k-1}}
{\pi^2+(L+\log y)^2}\,dy.
\]

For fixed \(k=1,2,3\),
\[
A_k(e^{-L})
=
e^{kL}L^{-2}\Gamma(k)
\left(1-\frac{2\psi(k)}{L}+O(L^{-2})\right),
\]
where \(\psi\) is the digamma function. To justify the expansion, split the \(y\)-integral into the central region \(|\log y|\le L^{2/3}\) and its complement. On the central interval, with \(z=\log y\), expand uniformly:
\[
\frac{1}{\pi^2+(L+\log y)^2}
=
L^{-2}\left(1-\frac{2\log y}{L}+O\left(\frac{1+\log^2 y}{L^2}\right)\right),
\]
and integrate against \(e^{-y}y^{k-1}\). The lower complement \(0<y<e^{-L^{2/3}}\) is \(O(e^{-cL^{2/3}})\) for \(k=1,2,3\) after the change \(y=e^{-u}\), and the upper complement \(y>e^{L^{2/3}}\) is super-exponentially small because of \(e^{-y}\). These errors are \(o(L^{-N})\) for every fixed \(N\), which is more than needed for the displayed expansion.

Therefore
\[
\frac{A_2(e^{-L})^2}{A_1(e^{-L})A_3(e^{-L})}
=
\frac{\Gamma(2)^2}{\Gamma(1)\Gamma(3)}
\left(
1+\frac{2(\psi(1)+\psi(3)-2\psi(2))}{L}+O(L^{-2})
\right).
\]
Using
\[
\psi(1)=-\gamma,\qquad
\psi(2)=1-\gamma,\qquad
\psi(3)=\frac32-\gamma,
\]
we get
\[
\psi(1)+\psi(3)-2\psi(2)=-\frac12,
\]
so
\[
\frac{A_2(e^{-L})^2}{A_1(e^{-L})A_3(e^{-L})}
=
\frac12-\frac{1}{2L}+O(L^{-2}).
\]

Choose
\[
\alpha_L=\frac12-\frac{1}{4L}.
\]
For all sufficiently large \(L\), \(\alpha_L\in(0,1/2)\) and
\[
\frac{A_2(e^{-L})^2}{A_1(e^{-L})A_3(e^{-L})}<\alpha_L.
\]
Consequently
\[
H_2(e^{-L};\alpha_L)
=
A_2(e^{-L})^2-\alpha_L A_1(e^{-L})A_3(e^{-L})
<0.
\]

Complete monotonicity would imply nonnegativity at order zero. Hence \(H_2(\cdot;\alpha_L)\) is not completely monotone for these fixed parameters \(\alpha_L\in(0,1/2)\). The source's universal Turan-window statement is false already in the \(n=2\) upper window.

The new mechanism is a logarithmic-density moment-ratio asymptotic obstruction for quadratic Laplace-moment Turan gaps. It is distinct from the public APP-0014 Stieltjes/complete-Bernstein classification of \(I_R\).
\end{proof}

The following appendix result records the fixed source problem \textbf{APP-0031} without interrupting the main exposition.
\begin{counterexample}[APP-0031: Baricz gamma-quotient Bernstein test is false]
\label{res:baricz-gamma-quotient-bf-forall-negative-answer}
Baricz's question asking whether \(x\mapsto\Gamma(x)\Gamma(x-a+b)/(\Gamma(x-a)\Gamma(x+b))\) is a Bernstein function on \((a,\infty)\) for every \(a,b>0\) has a negative answer.

This resolves the fixed application label \textbf{APP-0031}: Decide whether Baricz's gamma-quotient is Bernstein for every \(a,b>0\).
\emph{Source reference.} \cite{app-app-0031-source}.
\end{counterexample}
\begin{proof}
Take \(a=2\), \(b=3\), and set \(y=x-2>0\). Then by the Gamma recurrence,
\[
\frac{\Gamma(x)\Gamma(x-a+b)}{\Gamma(x-a)\Gamma(x+b)}
=\frac{\Gamma(x)\Gamma(x+1)}{\Gamma(x-2)\Gamma(x+3)}
=\frac{(x-2)(x-1)}{(x+1)(x+2)}
=\frac{y(y+1)}{(y+3)(y+4)}.
\]

Let
\[
g(y)=\frac{y(y+1)}{(y+3)(y+4)},\qquad y>0.
\]

If \(g\) were a Bernstein function, then \(g'\) would be completely monotone. In particular \(g'''\ge0\) would be necessary.

Exact symbolic differentiation gives
\[
g'(y)=\frac{6(y^2+4y+2)}{(y+3)^2(y+4)^2},
\]
\[
g''(y)=
-\frac{12(y^3+6y^2+6y-10)}{(y+3)^3(y+4)^3},
\]
and
\[
g'''(y)=
\frac{36(y^4+8y^3+12y^2-40y-94)}{(y+3)^4(y+4)^4}.
\]

At \(y=1\),
\[
g'''(1)=-\frac{1017}{40000}<0.
\]

Therefore \(g'\) is not completely monotone, \(g\) is not Bernstein, and the Baricz for-all-\(a,b\) Bernstein question has a negative answer.
\end{proof}

\begin{proposition}[Bulboaca-Zayed gamma quotient critical window reduction]
\label{res:bz-gamma-quotient-critical-window-reduction}
There is a rational interval \(J\) around \(1.126207061\ldots\) such that \(\widetilde F'<0\) on \((-1,\inf J]\), \(\widetilde F'>0\) on \([\sup J,\infty)\), and \(\widetilde F'\) has exactly one zero inside \(J\).
\end{proposition}
\begin{proof}
\emph{Setup.}
Use the notation from the previous Bulboaca--Zayed pass:
\[
G(x)=\log\Gamma(x+1),\qquad
D(x)=\log\frac{x^2+6}{x+6},
\]
and
\[
F(x)=\frac{G(x)}{D(x)}.
\]
For \(x>1\), \(D(x)>0\), and
\[
D'(x)=p(x)=\frac{x^2+12x-6}{(x^2+6)(x+6)}.
\]
Let
\[
N_0(x)=\psi(x+1)D(x)-G(x)D'(x).
\]
The derivative-sign numerator recorded in
the corresponding result is
\[
N(x)=((x^2+6)(x+6))\,N_0(x),
\]
so \(N\) and \(N_0\) have the same sign on \(x>1\).

The source itself proves the decreasing part on \((-1,1)\), and the earlier
the preceding theorem proves \(F'(x)>0\) for
\(x\ge8\).  It remains to prove a single sign change of \(N_0\) on \((1,8]\).

Put
\[
P(x)=x^2+12x-6,\qquad Q(x)=(x^2+6)(x+6),
\]
so \(p=P/Q\), and define
\[
r(x)=\frac{\psi(x+1)}{p(x)}=\psi(x+1)\frac{Q(x)}{P(x)}.
\]
Then
\[
P(x)^2r'(x)=S(x),
\]
where
\[
S(x)=\psi'(x+1)Q(x)P(x)+\psi(x+1)K(x),
\]
and
\[
K(x)=Q'(x)P(x)-Q(x)P'(x)
=x^4+24x^3+48x^2-144x-468.
\]
Direct differentiation gives
\[
S'(x)=P(x)\{Q(x)\psi''(x+1)+2Q'(x)\psi'(x+1)\}+K'(x)\psi(x+1).
\]

For \(y=x+1\ge2\), use the standard elementary bounds
\[
\psi(y)>\log y-\frac1y,\qquad
\psi'(y)>\frac1y+\frac{1}{2y^2},\qquad
\psi''(y)>-\frac1{y^2}-\frac2{y^3}.
\]
The last inequality follows from the series
\[
\psi''(y)=-2\sum_{n=0}^{\infty}(n+y)^{-3}
\]
and an integral upper bound for the sum.

Substituting these bounds into \(S'\) yields
\[
S'(x)>
P(x)\frac{5x^4+30x^3+57x^2+12x-90}{(x+1)^3}
K'(x)\left(\log(x+1)-\frac1{x+1}\right).
\]
On \(x\ge1\), \(P(x)>0\).  The quartic numerator has value \(14\) at \(x=1\)
and positive derivative
\[
20x^3+90x^2+114x+12.
\]
Also
\[
K'(x)=4x^3+72x^2+96x-144
\]
has value \(28\) at \(x=1\) and positive derivative for \(x\ge1\), while
\(\log(x+1)-1/(x+1)>0\) for \(x\ge1\).  Hence
\[
S'(x)>0\qquad(1\le x\le8).
\]

At the endpoints,
\[
S(1)=343\left(\frac{\pi^2}{6}-1\right)-539(1-\gamma)<0,
\]
for example using \(\pi^2<987/100\) and \(\gamma<5773/10000\), which gives
\(-66003/10000<0\).  Also \(S(8)>0\), because \(K(8)=17836>0\),
\(\psi'(9)>0\), and \(\psi(9)=H_8-\gamma>0\).  Therefore \(S\) has exactly
one zero on \([1,8]\).  Since \(P^2>0\), \(r'\) has exactly one zero; \(r\)
decreases first and then increases.

Now define
\[
I(x)=D(x)r(x)-G(x).
\]
Then
\[
N_0(x)=p(x)I(x),\qquad I'(x)=D(x)r'(x).
\]
Because \(D(x)>0\) for \(x>1\), \(I\) decreases and then increases with the
same turning point as \(r\).  Moreover \(\lim_{x\downarrow1}I(x)=0\), so
\(I(x)<0\) immediately to the right of \(1\).

Finally,
\[
N_0(8)=(H_8-\gamma)\log5-\frac{11}{70}\log(40320)>0.
\]
For instance, the bounds
\[
\gamma<\frac{5773}{10000},\quad \log5>\frac{1609}{1000},
\quad \log(40320)<\frac{10605}{1000}
\]
give the lower bound
\[
\left(\frac{761}{280}-\frac{5773}{10000}\right)\frac{1609}{1000}
-\frac{11}{70}\frac{10605}{1000}
=\frac{124435951}{70000000}>0.
\]
Thus \(I\), and therefore \(N_0\) and \(N\), has exactly one zero on
\((1,8]\).  Denote it by \(\xi\).  The sign pattern is
\[
N(x)<0\quad(1<x<\xi),\qquad
N(\xi)=0,\qquad
N(x)>0\quad(\xi<x\le8).
\]

A high-precision numerical check places
\[
\xi=1.12620706105164346558\ldots,
\]
matching the source's reported value.  The numerical value is not used in the
proof.

The raw numerator has a removable zero at \(x=1\), since \(G(1)=D(1)=0\).  Its
quadratic leading coefficient is
\[
\lim_{x\downarrow1}\frac{N_0(x)}{(x-1)^2}
=
\frac{7(\pi^2/6-1)-11(1-\gamma)}{98}<0,
\]
again by \(\pi^2<987/100\) and \(\gamma<5773/10000\), which gives
\(-1347/980000<0\).  Therefore any normalization dividing out the removable
\((x-1)^2\) factor is negative at the left endpoint.

The single-crossing candidate is true:
\[
T\text{-BZ-gamma-quotient-N-single-crossing-one-eight}.
\]
Together with the source's theorem on \((-1,1)\) and the previous local
right-tail theorem for \(x\ge8\), this proves the critical-window reduction
\[
T\text{-BZ-gamma-quotient-critical-window-reduction}
\]
and then the source problem
\[
T\text{-Bulboaca-Zayed-gamma-quotient-full-monotonicity}.
\]

The finite-envelope candidate is not separately promoted: the proof above is a
clean analytic ratio-kernel argument rather than the finite rational cover
specified in that candidate.  The diagnostic obstruction-map candidate is also
not needed after the single-crossing proof.
\end{proof}

The following appendix result records the fixed source problem \textbf{APP-0011} without interrupting the main exposition.
\begin{corollary}[APP-0011: Exact \(n=2\) beta-window endpoint]
\label{res:q2-terminal-exact}
The exact admissible beta set \(\mathcal I_2\) for Qi--Lim--Nantomah Open Problem 4 is determined by a certified value or certified description of \(L_2=\sup_{0<x<1}Q_2(x)\), including the endpoint-inclusion convention.

This resolves the fixed application label \textbf{APP-0011}: Exact \(n=2\) beta-window endpoint
\emph{Source reference.} \cite{app-app-0011-source}.
\end{corollary}
\begin{proof}
This roll proves a true auxiliary finite-middle outside-cover node:

the corresponding result.

The replayable certificate is:

\begin{verbatim}
\end{verbatim}

It proves
\[
Q_2(x)<q_J
\]
on
\[
\left[\frac9{50},\frac{1409}{5000}\right]
\cup
\left[\frac{293}{1000},\frac9{10}\right].
\]

The left bridge \([9/50,221050/1000000]\) uses the already-derived identity
\[
3xZ_4(x)-Z_3(x)
=2x^{-3}+\sum_{k=1}^{\infty}\frac{2x-k}{(x+k)^4}.
\]
For \(x<1/2\), this gives
\[
R(x)>Z_3(1/x).
\]
The certificate proves
\[
Z_3(1/x)>x^{1157/500}
\]
on \([9/50,221050/1000000]\). Since
\[
\frac{1157}{500}<\frac{115727}{50000}<q_J,
\]
this gives \(Q_2(x)<q_J\) on the bridge.

The remaining two subintervals use the direct comparison
\[
\log R(x)>\frac{115727}{50000}\log x.
\]
Since \(\log x<0\), this implies
\[
Q_2(x)<\frac{115727}{50000}<q_J.
\]
The script uses power-of-two range reduction for the atanh logarithm enclosure:
\[
\log r=\log(2^k r)-k\log 2,
\]
so every logarithm is evaluated with an argument in \([1,2)\).

The certificate output was:

\begin{verbatim}
REMAINING_FINITE_MIDDLE_COVER_OK
ZETA_LEFT_SUBCOVER
interval_count 245
domain [180000/1000000,221050/1000000]
min_width 1 max_width 1283
worst_interval [180000/1000000,181282/1000000]
worst_margin_positive True
DIRECT_LEFT_SUBCOVER
interval_count 3739
domain [221050/1000000,281800/1000000]
min_width 1 max_width 238
worst_interval [280389/1000000,280392/1000000]
worst_margin_positive True
DIRECT_RIGHT_SUBCOVER
interval_count 3801
domain [293000/1000000,900000/1000000]
min_width 1 max_width 18969
worst_interval [293124/1000000,293126/1000000]
worst_margin_positive True
\end{verbatim}

Together with the earlier true covers \((0,9/50]\) and \([9/10,1)\), this covers all outside-compact points below \(q_J\).

This roll also proves:

the corresponding result.

The replayable certificate is:

\begin{verbatim}
\end{verbatim}

Let
\[
H(x)=\Lambda(x)+x\Lambda'(x).
\]
Using \(Z_s'(x)=-sZ_{s+1}(x)\) and
\[
\frac{d}{dx}Z_s(1/x)=s x^{-2} Z_{s+1}(1/x),
\]
the script builds a rational interval enclosure for \(H\) using \(Z_3,\ldots,Z_6\) at \(x\) and \(Z_3,\ldots,Z_5\) at \(1/x\). It proves
\[
H(x)>0
\qquad
\left(\frac{1409}{5000}\le x\le\frac{293}{1000}\right).
\]

Since
\[
G'(x)=\log x\,H(x)
\]
and \(\log x<0\) on the compact bracket, this gives
\[
G'(x)<0.
\]
The existing endpoint certificate proves
\[
G(287345/1000000)>0,
\qquad
G(287346/1000000)<0.
\]
Therefore \(G\) has a unique zero
\[
\xi\in J=\left[\frac{287345}{1000000},\frac{287346}{1000000}\right],
\]
with \(G>0\) to the left of \(J\) and \(G<0\) to the right of \(J\) on the compact bracket.

Because
\[
Q_2'(x)=\frac{G(x)}{x(\log x)^2},
\]
\(Q_2\) is increasing before \(\xi\) and decreasing after \(\xi\) on the compact bracket.

The certificate output was:

\begin{verbatim}
COMPACT_MONOTONICITY_CERTIFICATE_OK
endpoint_signs G_left_positive G_right_negative
h_interval_count 128
domain [281800/1000000,293000/1000000]
min_width 87 max_width 88
worst_interval [281887/1000000,281975/1000000]
worst_H_lower_positive True
\end{verbatim}

The true cover package is now:

\[
(0,9/50]\cup
\left[9/50,\frac{1409}{5000}\right]\cup
\left[\frac{293}{1000},9/10\right]\cup
[9/10,1).
\]

All points outside the compact bracket have \(Q_2(x)<q_J\), and \(q_J<Q_2(\xi)\) by the certified inner witness. On the compact bracket, \(Q_2\) has a unique maximum at the unique zero \(\xi\in J\). Thus
\[
L_2=Q_2(\xi).
\]

At \(\beta=Q_2(\xi)\), the defining inequality has equality at \(x=\xi\), so the lower endpoint is excluded. The known upper endpoint remains included. Therefore
\[
\mathcal I_2=(Q_2(\xi),3].
\]

This proves the terminal exact endpoint node.
\end{proof}

The following appendix result records the fixed source problem \textbf{APP-0019} without interrupting the main exposition.
\begin{corollary}[APP-0019: Wakrim W-symbol Bernstein range closes the beta gap]
\label{res:wakrim-w-operator-bf-gap-solved}
Wakrim's open Bernstein-status gap for the W-symbol in the range \(0<\beta\le1\) is resolved affirmatively: \(\Phi_{\alpha,\beta}\) is Bernstein for every \(0<\alpha<1\) and \(0<\beta\le1\).

This resolves the fixed application label \textbf{APP-0019}: Close the open Bernstein-symbol range for the W-operator when \(0<\beta\le1\).
\emph{Source reference.} \cite{app-app-0019-source}.
\end{corollary}
\begin{proof}
Wakrim's source proves the non-Bernstein obstruction for \(\beta>1\) and
leaves the range \(0<\beta\le1\) as the remaining delicate case.  We prove
that remaining Bernstein half in detail.  No complete-Bernstein or
subordination conclusion is asserted here.

Put
\[
  a=1-\alpha\in(0,1),\qquad y=s^a .
\]
Since \(\alpha=1-a\),
\[
 \Phi_{\alpha,\beta}(s)
 =s^{1-a}(1+a s^{-a})^{-\beta}
 =s^{1-a+a\beta}(s^a+a)^{-\beta}.
\]
Let
\[
  p=1-a+a\beta=\alpha+a\beta .
\]
Differentiating the last expression gives
\[
\begin{aligned}
 \Phi'_{\alpha,\beta}(s)
 &=s^{p-1}(s^a+a)^{-\beta-1}
     \{p(s^a+a)-\beta a s^a\}  \\
 &=s^{p-1}(s^a+a)^{-\beta-1}
     \{\alpha s^a+\alpha a+a^2\beta\}       \\
 &=s^{p-1}(s^a+a)^{-\beta}
     \left(\alpha+{a^2\beta\over s^a+a}\right).
\end{aligned}
\]
Because \(p-1=a(\beta-1)\), this is
\[
  \Phi'_{\alpha,\beta}(s)=H_{\alpha,\beta}(s^a),
\]
where
\[
 H_{\alpha,\beta}(y)
 =y^{\beta-1}(y+a)^{-\beta}
      \left(\alpha+{a^2\beta\over y+a}\right).
\]

We now check complete monotonicity of \(H_{\alpha,\beta}\) for
\(0\le\beta\le1\).  The first factor is
\[
  y^{\beta-1}=y^{-(1-\beta)},
\]
and this is interpreted as follows.  If \(0\le\beta<1\), put
\(r=1-\beta>0\).  Then
\[
  y^{\beta-1}=y^{-r}
  ={1\over\Gamma(r)}
    \int_0^\infty e^{-yt}t^{r-1}\,dt,
\]
so \(y^{\beta-1}\) is completely monotone by the Bernstein--Widder
criterion.  If \(\beta=1\), the same factor is the constant \(1\).  The shifted factor
\((y+a)^{-\beta}\) is completely monotone for \(\beta>0\) by the Laplace
representation
\[
 (y+a)^{-\beta}
 ={1\over\Gamma(\beta)}
   \int_0^\infty e^{-yt}e^{-at}t^{\beta-1}\,dt,
\]
and for \(\beta=0\) it is the constant \(1\).  The last factor
\[
  \alpha+{a^2\beta\over y+a}
\]
is a positive constant plus a nonnegative multiple of the completely monotone
resolvent \((y+a)^{-1}\).  By closure of completely monotone functions under
products and nonnegative sums, \(H_{\alpha,\beta}\) is completely monotone.

The endpoint \(\beta=1\) has no limiting ambiguity; explicitly
\[
 H_{\alpha,1}(y)
 ={\alpha\over y+a}+{a^2\over (y+a)^2},
\]
which is completely monotone.  The endpoint \(\beta=0\) also agrees with the
direct identity \(\Phi_{\alpha,0}(s)=s^\alpha\).

Finally, let \(g(s)=s^a\).  The composition theorem used here is the standard
one: if \(f\) is completely monotone on \((0,\infty)\) and \(g\) is a
Bernstein function with \(g((0,\infty))\subset(0,\infty)\), then \(f\circ g\)
is completely monotone.  Its hypotheses are satisfied in the present case:
\(H_{\alpha,\beta}\) is completely monotone by the factor check above,
\(g(s)=s^a>0\), and \(g\) is a Bernstein function for \(0<a<1\) since
\[
  g'(s)=a s^{a-1}=a s^{-(1-a)}
\]
is completely monotone.  Hence
\(\Phi'_{\alpha,\beta}=H_{\alpha,\beta}\circ g\) is completely monotone for
\(0\le\beta\le1\), so \(\Phi_{\alpha,\beta}\) is a Bernstein function.
Together with Wakrim's source obstruction for \(\beta>1\), this closes the
remaining case and gives the exact Bernstein range.
\end{proof}

The following appendix result records the fixed source problem \textbf{APP-0006} without interrupting the main exposition.
\begin{corollary}[APP-0006: Exact inverse-tail floor formula at s=8]
\label{res:zeta-tail-floor-next-case}
Find a new exact reciprocal zeta-tail floor formula beyond the staged \(s=7\) and \(s=8\) cases, or prove a reusable enclosure template that yields such formulas for further integer exponents.

This resolves the fixed application label \textbf{APP-0006}: Exact inverse-tail floor formula at \(s=8\)
\emph{Source reference.} \cite{app-app-0006-source}.
\end{corollary}
\begin{proof}
For integer \(s>1\), the Hurwitz-tail Euler--Maclaurin expansion has the form
\[
\zeta_n(s)
=n^{1-s}\left[
\frac1{s-1}+\frac{z}{2}
+\sum_{r\ge1}\frac{B_{2r}}{(2r)!}(s)_{2r-1}z^{2r}
\right],
\qquad z=\frac1n.
\]
Formally inverting the bracket gives
\[
\zeta_n(s)^{-1}\sim n^{s-1}G_s(z).
\]

The replay script
computes \(G_s\) symbolically and verifies two things:

1. For \(s=8\), the coefficients through \(z^{15}\) reproduce the staged \(A_8\) and \(B_8-A_8\) coefficients.
2. For \(s=9\), truncating \(n^8G_9(1/n)\) through \(z^{14}\) gives an \(A_9\), and adding the \(z^{15}\) term gives \(B_9\), with exact telescoping residuals of the expected signs.

The generated \(s=9\) lower approximant is
\[
\begin{aligned}
A_9(n)=&8n^8-32n^7+80n^6-128n^5+120n^4-64n^3
+\frac{624}{7}n^2-\frac{512}{7}n-\frac{3324}{7}\\
&+\frac{90304}{21n^2}+\frac{90304}{21n^3}
-\frac{10280944}{245n^4}-\frac{64846304}{735n^5}
+\frac{109940976}{245n^6}.
\end{aligned}
\]
The upper approximant is
\[
B_9(n)=A_9(n)+\frac{384388288}{245n^7}.
\]

Let \(R_A(n)=1/A_9(n)\) and \(R_B(n)=1/B_9(n)\). Exact denominator clearing gives:
\[
R_A(k)-R_A(k+1)-\frac1{k^9}>0
\qquad(k\ge5),
\]
because the numerator, after substituting \(k=m+5\), has all coefficients positive. Similarly,
\[
\frac1{k^9}-\left(R_B(k)-R_B(k+1)\right)>0
\qquad(k\ge1),
\]
because the denominator-cleared numerator, after substituting \(k=m+1\), has all coefficients positive.

Therefore, for \(n\ge5\),
\[
A_9(n)<\zeta_n(9)^{-1}<B_9(n).
\]

This proves the inverse-window part of the reusable template and promotes the corresponding result.

The replay script also verifies the no-integer gap.

Write \(P_9(n)\) for the polynomial part of \(A_9(n)\):
\[
P_9(n)=8n^8-32n^7+80n^6-128n^5+120n^4-64n^3
+\frac{624}{7}n^2-\frac{512}{7}n-\frac{3324}{7}.
\]
Modulo \(7\), the fractional numerator of \(P_9(n)\) is
\[
n^2-n+1\pmod 7,
\]
so the fractional part of \(P_9(n)\) is one of
\[
0,\quad \frac17,\quad \frac37,\quad \frac67.
\]

The correction \(A_9(n)-P_9(n)\) is positive for \(n\ge9\); after denominator clearing and substituting \(n=m+9\), the numerator has all coefficients positive. For \(n\ge174\), the upper correction is bounded by
\[
B_9(n)-P_9(n)
<
\frac{90304}{21n^2}
+\frac{90304}{21n^3}
+\frac{109940976}{245n^6}
+\frac{384388288}{245n^7}
<\frac17.
\]
Therefore \(A_9(n)\) and \(B_9(n)\) cannot straddle an integer for \(n\ge174\). Exact rational arithmetic verifies
\[
\lfloor A_9(n)\rfloor=\lfloor B_9(n)\rfloor
\qquad(9\le n\le173).
\]
Consequently,
\[
\left\lfloor\zeta_n(9)^{-1}\right\rfloor=\lfloor A_9(n)\rfloor
\qquad(n\ge9).
\]

For \(1\le n\le8\), the exact values are
\[
\begin{array}{c|rrrrrrrr}
n&1&2&3&4&5&6&7&8\\
\hline
\lfloor\zeta_n(9)^{-1}\rfloor&0&497&18093&224086&1543530&7360226&27295287&84349541.
\end{array}
\]
For \(n=1\), this follows from \(\zeta(9)>1\). For \(2\le n\le8\), the replay script uses exact rational partial sums through \(N=100\) and the integral tail bound
\[
\sum_{k=N+1}^{\infty}k^{-9}<\frac{1}{8N^8}
\]
to prove
\[
\frac1{M_n+1}<\zeta_n(9)<\frac1{M_n}.
\]

This proves the claim.
\end{proof}

\begin{lemma}[Baskakov diagonal r=4 c=n generalized sum is completely monotone]
\label{res:baskakov-diagonal-r4-c-equals-n-cm}
For every positive integer \(n\), the diagonal generalized Baskakov sum \(\psi^{[4]}_{n,n}(x)=\sum_{k=0}^{\infty}(p^{[n]}_{n,k}(x))^4\) is completely monotone on \((0,\infty)\).
\end{lemma}
\begin{proof}
For \(\alpha=1\),
\[
\binom{-1}{k}=(-1)^k,
\qquad
\binom{-1}{k}^{4}=1.
\]
Therefore, for \(x>0\),
\[
f^{[4]}_1(x)
=(1+x)^{-4}\sum_{k=0}^{\infty}
\left(\frac{x}{1+x}\right)^{4k}
=\frac{1}{(1+x)^4-x^4}.
\]
Since
\[
(1+x)^4-x^4=(2x+1)(2x^2+2x+1),
\]
partial fractions give
\[
\frac{1}{(1+x)^4-x^4}
=\frac{2}{2x+1}-\frac{2x+1}{2x^2+2x+1}.
\]
Writing \(a=x+\frac12\), this is
\[
\frac{1}{a}-\frac{a}{a^2+(1/2)^2}.
\]
Using
\[
\int_0^\infty e^{-xt}e^{-t/2}\,dt=\frac{1}{x+1/2}
\]
and
\[
\int_0^\infty e^{-xt}e^{-t/2}\cos(t/2)\,dt
=\frac{x+1/2}{(x+1/2)^2+(1/2)^2},
\]
we obtain
\[
f^{[4]}_1(x)
=\int_0^\infty e^{-xt}e^{-t/2}\left(1-\cos\frac{t}{2}\right)\,dt.
\]
The density \(e^{-t/2}(1-\cos(t/2))\) is nonnegative on \([0,\infty)\). By Bernstein's theorem, \(f^{[4]}_1\) is completely monotone on \((0,\infty)\).

From the source definition
\[
p^{[c]}_{n,k}(x)=\binom{-n/c}{k}(-cx)^k(1+cx)^{-n/c-k},
\]
one has
\[
\psi^{[r]}_{n,c}(x)=\sum_{k=0}^{\infty}\left(p^{[c]}_{n,k}(x)\right)^r
=f^{[r]}_{n/c}(cx)
\]
for even \(r\). Hence, for \(c=n\),
\[
\psi^{[4]}_{n,n}(x)=f^{[4]}_1(nx).
\]
Positive scaling preserves complete monotonicity, so \(\psi^{[4]}_{n,n}\) is completely monotone for every positive integer \(n\).
\end{proof}

The following appendix result records the fixed source problem \textbf{APP-0033} without interrupting the main exposition.
\begin{corollary}[APP-0033: Ramanujan antiderivative is complete Bernstein]
\label{res:ramanujan-antiderivative-complete-bernstein}
The antiderivative \(\widetilde I_R(x)=a+\int_0^\infty(1-e^{-xt})\,dt/[t(\pi^2+\log^2 t)]\) of the Ramanujan integral is a complete Bernstein function.

This resolves the fixed application label \textbf{APP-0033}: Resolve whether the Ramanujan integral antiderivative is a complete Bernstein function.
\emph{Source reference.} \cite{app-app-0033-source}.
\end{corollary}
\begin{proof}
The density
\[
\phi_0(t)=\frac{1}{t(\pi^2+\log^2 t)}
\]
is completely monotone on \((0,\infty)\). Consequently \(I_R\) is a Stieltjes function. Moreover, the antiderivative
\[
\widetilde I_R(x)=a+\int_0^\infty (1-e^{-xt})\phi_0(t)\,dt
\]
is a complete Bernstein function.

For \(u=\log t\),
\[
\int_0^1 e^{-au}\sin(\pi a)\,da
=\frac{\pi(1+e^{-u})}{u^2+\pi^2}.
\]
Substituting \(u=\log t\) gives
\[
\int_0^1 t^{-a}\sin(\pi a)\,da
=\frac{\pi(1+t^{-1})}{\pi^2+\log^2 t}.
\]
Therefore
\[
\phi_0(t)
=\frac{1}{t(\pi^2+\log^2 t)}
=\frac{1}{\pi(1+t)}\int_0^1 t^{-a}\sin(\pi a)\,da.
\]

For \(0<a<1\), \(t^{-a}\) is completely monotone because
\[
t^{-a}=\frac1{\Gamma(a)}\int_0^\infty e^{-ts}s^{a-1}\,ds.
\]
Also \(t\mapsto(1+t)^{-1}\) is completely monotone. Complete monotonicity is closed under products and positive mixtures, and \(\sin(\pi a)>0\) on \((0,1)\). Hence \(\phi_0\) is completely monotone.

If \(\phi\) is completely monotone and
\[
F(x)=\int_0^\infty e^{-xt}\phi(t)\,dt
\]
is finite for \(x>0\), then \(F\) is Stieltjes: write \(\phi(t)=\int_0^\infty e^{-ts}\,d\mu(s)\) and use Tonelli to obtain
\[
F(x)=\int_0^\infty\frac{d\mu(s)}{x+s}.
\]
Applying this to \(\phi_0\) proves that \(I_R\) is Stieltjes.

Finally, the source already proves that
\[
\widetilde I_R(x)=a+\int_0^\infty (1-e^{-xt})\phi_0(t)\,dt
\]
is a Bernstein function. The Levy density \(\phi_0\) is completely monotone, and it satisfies the Bernstein integrability condition
\[
\int_0^\infty (1\wedge t)\phi_0(t)\,dt<\infty.
\]
The standard complete-Bernstein criterion for Levy densities therefore gives that \(\widetilde I_R\) is a complete Bernstein function.
\end{proof}

The following appendix result records the fixed source problem \textbf{APP-0034} without interrupting the main exposition.
\begin{corollary}[APP-0034: Bulboaca-Zayed gamma-quotient monotonicity solved]
\label{res:bulboaca-zayed-gamma-quotient-full-monotonicity}
For the continuous extension \(\widetilde F\) of \(F(x)=\log\Gamma(x+1)/(\log(x^2+6)-\log(x+6))\) on \((-1,\infty)\), prove analytically that \(\widetilde F'\) has a unique zero \(x_m\simeq1.126207061\ldots\), with \(\widetilde F'<0\) on \((-1,x_m)\) and \(\widetilde F'>0\) on \((x_m,\infty)\).

This resolves the fixed application label \textbf{APP-0034}: Solve the Bulboaca-Zayed one-critical-point monotonicity statement for the gamma quotient.
\emph{Source reference.} \cite{app-app-0034-source}.
\end{corollary}
\begin{proof}
\emph{Setup.}
Use the notation from the previous Bulboaca--Zayed pass:
\[
G(x)=\log\Gamma(x+1),\qquad
D(x)=\log\frac{x^2+6}{x+6},
\]
and
\[
F(x)=\frac{G(x)}{D(x)}.
\]
For \(x>1\), \(D(x)>0\), and
\[
D'(x)=p(x)=\frac{x^2+12x-6}{(x^2+6)(x+6)}.
\]
Let
\[
N_0(x)=\psi(x+1)D(x)-G(x)D'(x).
\]
The derivative-sign numerator recorded in
the corresponding result is
\[
N(x)=((x^2+6)(x+6))\,N_0(x),
\]
so \(N\) and \(N_0\) have the same sign on \(x>1\).

The source itself proves the decreasing part on \((-1,1)\), and the earlier
the preceding theorem proves \(F'(x)>0\) for
\(x\ge8\).  It remains to prove a single sign change of \(N_0\) on \((1,8]\).

Put
\[
P(x)=x^2+12x-6,\qquad Q(x)=(x^2+6)(x+6),
\]
so \(p=P/Q\), and define
\[
r(x)=\frac{\psi(x+1)}{p(x)}=\psi(x+1)\frac{Q(x)}{P(x)}.
\]
Then
\[
P(x)^2r'(x)=S(x),
\]
where
\[
S(x)=\psi'(x+1)Q(x)P(x)+\psi(x+1)K(x),
\]
and
\[
K(x)=Q'(x)P(x)-Q(x)P'(x)
=x^4+24x^3+48x^2-144x-468.
\]
Direct differentiation gives
\[
S'(x)=P(x)\{Q(x)\psi''(x+1)+2Q'(x)\psi'(x+1)\}+K'(x)\psi(x+1).
\]

For \(y=x+1\ge2\), use the standard elementary bounds
\[
\psi(y)>\log y-\frac1y,\qquad
\psi'(y)>\frac1y+\frac{1}{2y^2},\qquad
\psi''(y)>-\frac1{y^2}-\frac2{y^3}.
\]
The last inequality follows from the series
\[
\psi''(y)=-2\sum_{n=0}^{\infty}(n+y)^{-3}
\]
and an integral upper bound for the sum.

Substituting these bounds into \(S'\) yields
\[
S'(x)>
P(x)\frac{5x^4+30x^3+57x^2+12x-90}{(x+1)^3}
K'(x)\left(\log(x+1)-\frac1{x+1}\right).
\]
On \(x\ge1\), \(P(x)>0\).  The quartic numerator has value \(14\) at \(x=1\)
and positive derivative
\[
20x^3+90x^2+114x+12.
\]
Also
\[
K'(x)=4x^3+72x^2+96x-144
\]
has value \(28\) at \(x=1\) and positive derivative for \(x\ge1\), while
\(\log(x+1)-1/(x+1)>0\) for \(x\ge1\).  Hence
\[
S'(x)>0\qquad(1\le x\le8).
\]

At the endpoints,
\[
S(1)=343\left(\frac{\pi^2}{6}-1\right)-539(1-\gamma)<0,
\]
for example using \(\pi^2<987/100\) and \(\gamma<5773/10000\), which gives
\(-66003/10000<0\).  Also \(S(8)>0\), because \(K(8)=17836>0\),
\(\psi'(9)>0\), and \(\psi(9)=H_8-\gamma>0\).  Therefore \(S\) has exactly
one zero on \([1,8]\).  Since \(P^2>0\), \(r'\) has exactly one zero; \(r\)
decreases first and then increases.

Now define
\[
I(x)=D(x)r(x)-G(x).
\]
Then
\[
N_0(x)=p(x)I(x),\qquad I'(x)=D(x)r'(x).
\]
Because \(D(x)>0\) for \(x>1\), \(I\) decreases and then increases with the
same turning point as \(r\).  Moreover \(\lim_{x\downarrow1}I(x)=0\), so
\(I(x)<0\) immediately to the right of \(1\).

Finally,
\[
N_0(8)=(H_8-\gamma)\log5-\frac{11}{70}\log(40320)>0.
\]
For instance, the bounds
\[
\gamma<\frac{5773}{10000},\quad \log5>\frac{1609}{1000},
\quad \log(40320)<\frac{10605}{1000}
\]
give the lower bound
\[
\left(\frac{761}{280}-\frac{5773}{10000}\right)\frac{1609}{1000}
-\frac{11}{70}\frac{10605}{1000}
=\frac{124435951}{70000000}>0.
\]
Thus \(I\), and therefore \(N_0\) and \(N\), has exactly one zero on
\((1,8]\).  Denote it by \(\xi\).  The sign pattern is
\[
N(x)<0\quad(1<x<\xi),\qquad
N(\xi)=0,\qquad
N(x)>0\quad(\xi<x\le8).
\]

A high-precision numerical check places
\[
\xi=1.12620706105164346558\ldots,
\]
matching the source's reported value.  The numerical value is not used in the
proof.

The raw numerator has a removable zero at \(x=1\), since \(G(1)=D(1)=0\).  Its
quadratic leading coefficient is
\[
\lim_{x\downarrow1}\frac{N_0(x)}{(x-1)^2}
=
\frac{7(\pi^2/6-1)-11(1-\gamma)}{98}<0,
\]
again by \(\pi^2<987/100\) and \(\gamma<5773/10000\), which gives
\(-1347/980000<0\).  Therefore any normalization dividing out the removable
\((x-1)^2\) factor is negative at the left endpoint.

The single-crossing candidate is true:
\[
T\text{-BZ-gamma-quotient-N-single-crossing-one-eight}.
\]
Together with the source's theorem on \((-1,1)\) and the previous local
right-tail theorem for \(x\ge8\), this proves the critical-window reduction
\[
T\text{-BZ-gamma-quotient-critical-window-reduction}
\]
and then the source problem
\[
T\text{-Bulboaca-Zayed-gamma-quotient-full-monotonicity}.
\]

The finite-envelope candidate is not separately promoted: the proof above is a
clean analytic ratio-kernel argument rather than the finite rational cover
specified in that candidate.  The diagnostic obstruction-map candidate is also
not needed after the single-crossing proof.
\end{proof}

The following appendix result records the fixed source problem \textbf{APP-0015} without interrupting the main exposition.
\begin{corollary}[APP-0015: Sibisi-Prabhakar measure is a stable-subordination push-forward]
\label{res:sibisi-prabhakar-q-measure-problem-solved-canonical}
In Sibisi's strict range \(0<\alpha<1\), \(\gamma>0\), \(\beta>\alpha\gamma\), the source problem of determining \(Q^\gamma_{\alpha,\beta}(\cdot\mid\lambda)\) with Laplace transform \(E^\gamma_{\alpha,\beta}(-\lambda x^\alpha)\) is solved at the canonical finite-measure level: \(Q\) is the stable subordination of the transform-normalized Pollard measure \(P^\gamma_{\alpha,\beta}\).

This resolves the fixed application label \textbf{APP-0015}: Give a canonical probabilistic representation of the Prabhakar Q-measure in the strict Pollard range.
\emph{Source reference.} \cite{app-app-0015-source}.
\end{corollary}
\begin{proof}
Conditioning on \(r\) and using the stable transform gives
\[
 \mathbb E e^{-x(\lambda r)^{1/\alpha}S_\alpha}=e^{-\lambda r x^\alpha}.
\]
Integration against \(P_{\alpha,\beta}^{\eta}\) gives the Pollard
representation of \(E_{\alpha,\beta}^{\eta}(-\lambda x^\alpha)\).  The total
mass is finite, equal to the value at \(x=0\), and no normalization to a
probability measure is needed for the transform identity.
\end{proof}

\section{Application ledger}
This appendix is the immutable public memory of solved source problems.  APP numbers and names are fixed labels; theorem numbers in the main text are allowed to change as the exposition improves.
\begin{enumerate}[leftmargin=*]
\item \textbf{APP-0001.} Alzer-Kwong convexity and concavity problem.  Source problem: Alzer-Kwong convexity and concavity problem. Source reference: \cite{app-app-0001-source}. Resolution: Alzer-Kwong convexity and concavity pattern for reciprocal zeta. Main-text result: ledger-only in this rendering.
\item \textbf{APP-0002.} Nantomah zeta positivity problem.  Source problem: Nantomah zeta positivity problem. Source reference: \cite{app-app-0002-source}. Resolution: Affirmative solution of Nantomah zeta positivity problem. Main-text result: ledger-only in this rendering.
\item \textbf{APP-0003.} Sroysang generalized Holder problem.  Source problem: Sroysang generalized Holder problem. Source reference: \cite{app-app-0003-source}. Resolution: Generalized Holder inequality for Gamma zeta. Main-text result: ledger-only in this rendering.
\item \textbf{APP-0004.} Complete monotonicity of \((\log\Gamma(x)+\log\Gamma(1/x))''\).  Source problem: Complete monotonicity of \((\log\Gamma(x)+\log\Gamma(1/x))''\). Source reference: \cite{app-app-0004-source}. Resolution: Complete monotonicity of reciprocal-Gamma curvature. Main-text result: \ref{res:baricz-vq-strict-q-logconvexity}.
\item \textbf{APP-0005.} Exact inverse-tail floor formula at s=7.  Source problem: Exact inverse-tail floor formula at \(s=7\). Source reference: \cite{app-app-0005-source}. Resolution: Exact inverse-tail floor formula at \(s=7\). Main-text result: \ref{res:zeta-tail-floor-next-case}.
\item \textbf{APP-0006.} Exact inverse-tail floor formula at s=8.  Source problem: Exact inverse-tail floor formula at \(s=8\). Source reference: \cite{app-app-0006-source}. Resolution: Exact inverse-tail floor formula at \(s=8\). Main-text result: \ref{res:zeta-tail-floor-next-case}.
\item \textbf{APP-0007.} Concavity or complete monotonicity of the polygamma product P0.  Source problem: Concavity or complete monotonicity of the polygamma product \(P\_0\). Source reference: \cite{app-app-0007-source}. Resolution: Complete monotonicity of reciprocal digamma product curvature. Main-text result: \ref{res:yin-pk-digamma-alpha-necessity}.
\item \textbf{APP-0008.} Higher-order monotonicity of polygamma products Pn.  Source problem: Higher-order monotonicity of polygamma products \(P\_n\). Source reference: \cite{app-app-0008-source}. Resolution: Counterexample to complete monotonicity of higher-order polygamma product curvature. Main-text result: \ref{res:baricz-vq-strict-q-logconvexity}.
\item \textbf{APP-0009.} Sharp reciprocal Gamma-product monotonicity threshold.  Source problem: Sharp reciprocal Gamma-product monotonicity threshold. Source reference: \cite{app-app-0009-source}. Resolution: APP-0009: Sharp threshold for reciprocal Gamma-product monotonicity. Main-text result: \ref{res:qg-tau-global-supremum}.
\item \textbf{APP-0010.} Nielsen \(k\)-beta derivative-ratio monotonicity.  Source problem: Nielsen \(k\)-beta derivative-ratio monotonicity. Source reference: \cite{app-app-0010-source}. Resolution: Parity law from a complete-monotone Laplace moment-ratio bridge. Main-text result: \ref{res:nielsen-k-beta-derivative-ratio-monotonicity}.
\item \textbf{APP-0011.} Exact \(n=2\) beta-window endpoint.  Source problem: Exact \(n=2\) beta-window endpoint. Source reference: \cite{app-app-0011-source}. Resolution: Exact certified endpoint for the remaining Qi--Lim--Nantomah beta-window case. Main-text result: \ref{res:q2-terminal-exact}.
\item \textbf{APP-0012.} Baricz \(V\_q\) strict \(q\)-log-convexity.  Source problem: Baricz \(V\_q\) strict \(q\)-log-convexity. Source reference: \cite{app-app-0012-source}. Resolution: Positive Laplace representation proving strict parameter log-convexity. Main-text result: \ref{res:baricz-vq-strict-q-logconvexity}.
\item \textbf{APP-0013.} Yin \((p,k)\)-digamma sharp \(\alpha\)-necessity.  Source problem: Yin \((p,k)\)-digamma sharp \(\alpha\)-necessity. Source reference: \cite{app-app-0013-source}. Resolution: Near-zero obstruction proving complete monotonicity forces \(\alpha\le1\). Main-text result: \ref{res:yin-pk-digamma-alpha-necessity}.
\item \textbf{APP-0014.} Ramanujan integral is Stieltjes and has a complete Bernstein primitive.  Source problem: Classify the Ramanujan integral by a standard positive-kernel transform rather than only by inequalities.. Source reference: \cite{app-app-0014-source}. Resolution: The density is completely monotone, so the integral is Stieltjes and its Bernstein primitive is complete Bernstein.. Main-text result: \ref{res:ramanujan-density-logsquare-cm}.
\item \textbf{APP-0015.} Sibisi-Prabhakar measure is a stable-subordination push-forward.  Source problem: Give a canonical probabilistic representation of the Prabhakar Q-measure in the strict Pollard range.. Source reference: \cite{app-app-0015-source}. Resolution: The Q-measure is the alpha-stable subordination of the transform-normalized Pollard measure.. Main-text result: \ref{res:sibisi-prabhakar-q-measure-problem-solved-canonical}.
\item \textbf{APP-0016.} From's Mills-ratio bound extends to every derivative level.  Source problem: Lift the published Mills-ratio bound to the full derivative hierarchy.. Source reference: \cite{app-app-0016-source}. Resolution: The level-\(L\) determinant inequality yields an explicit bound for every derivative quotient and hence for the original Mills ratio.. Main-text result: \ref{res:from-mills-explicit-bound-family-frontier}.
\item \textbf{APP-0017.} Baricz gamma-quotient Bernstein-for-all problem has a rational counterexample.  Source problem: Decide whether the gamma quotient is Bernstein throughout the full parameter range.. Source reference: \cite{app-app-0017-source}. Resolution: The parameters a=2, b=3 give a rational quotient whose derivative is not completely monotone.. Main-text result: \ref{res:baricz-gamma-quotient-a2b3-not-bf}.
\item \textbf{APP-0018.} Qi-Guo tau threshold has a supremum-corrected exact value.  Source problem: Find the optimal uniform threshold for the tau expression in Open Problem 3.. Source reference: \cite{app-app-0018-source}. Resolution: The exact threshold is the supremum \(a\_*/(1+a\_*+a\_*^2)\), where \(e^{a\_*}=1+a\_*+a\_*^2\); the value is not attained.. Main-text result: \ref{res:qg-tau-global-supremum}.
\item \textbf{APP-0019.} Wakrim W-symbol Bernstein range closes the beta gap.  Source problem: Close the open Bernstein-symbol range for the W-operator when \(0<\beta\le1\).. Source reference: \cite{app-app-0019-source}. Resolution: Wakrim's source supplies the \(\beta>1\) obstruction; the staged proof proves the remaining Bernstein range \(0\le\beta\le1\).. Main-text result: \ref{res:wakrim-w-operator-bf-gap-solved}.
\item \textbf{APP-0020.} Qi's degree-four conjecture for h\_lambda is refuted at the endpoint.  Source problem: Test the conjectured complete-monotonic degree four for \(h\_\lambda\) and \(-h\_\mu\).. Source reference: \cite{app-app-0020-source}. Resolution: The degree-four transforms fail the first derivative test near zero, so the conjecture is false.. Main-text result: \ref{res:not-qi-hlambda-degree4-conjecture}.
\item \textbf{APP-0021.} Szabo's cutoff problem has \(\alpha\_0=1\).  Source problem: Determine the exact cutoff \(\alpha\_0\) for \(y^\alpha H\_d(y)\).. Source reference: \cite{app-app-0021-source}. Resolution: The cutoff is \(\alpha\_0=1\): sufficiency is the source theorem and necessity follows from the singular expansion at zero.. Main-text result: \ref{res:dagum-c-endpoints-and-upper-gate}.
\item \textbf{APP-0022.} Chen-Choi Gurland gamma-ratio conjecture is logarithmically completely monotone.  Source problem: Prove logarithmic complete monotonicity of the Gurland gamma ratio conjectured in the source.. Source reference: \cite{app-app-0022-source}. Resolution: A Weierstrass product reduces log F to a sum of strictly completely monotone logarithmic second differences.. Main-text result: \ref{res:chen-choi-gurland-logf-strict-cm}.
\item \textbf{APP-0023.} Sokal generalized-Stieltjes lambda-derivatives have a triangular nonnegative compression.  Source problem: Explain the lambda-derivative structure of Sokal's generalized-Stieltjes Hankel-type expressions.. Source reference: \cite{app-app-0023-source}. Resolution: The exponential generating function gives a triangular nonnegative convolution formula for every lambda derivative.. Main-text result: \ref{res:sokal-gs-lambda-derivative-generating-function}.
\item \textbf{APP-0024.} Simon gamma quotient is Bernstein.  Source problem: Determine whether Simon's gamma quotient \(F\_\alpha(x)=\Gamma(x+\alpha)/(\Gamma(x)x^\alpha)\) is Bernstein for \(0<\alpha<1\).. Source reference: \cite{app-app-0024-source}. Resolution: A complement complete-monotonicity representation proves \(F\_\alpha\) is Bernstein.. Main-text result: \ref{res:baricz-gamma-quotient-bf-forall-negative-answer}.
\item \textbf{APP-0025.} Du-Wang \(h\_3\) monotonicity is classified.  Source problem: Classify Du-Wang's \(h\_3\) monotonicity on \((0,\infty)\) for \(0<a<2\).. Source reference: \cite{app-app-0025-source}. Resolution: A polygamma-ratio kernel proves increasing behavior exactly for \(1/2\le a\le1\), with nonmonotonicity outside.. Main-text result: \ref{res:du-wang-h3-open-problem-2-classification}.
\item \textbf{APP-0026.} Baskakov even-power complete-monotonicity conjecture is false.  Source problem: Decide the Abel-Gawronski-Neuschel complete-monotonicity conjecture for all even powers in the Baskakov family.. Source reference: \cite{app-app-0026-source}. Resolution: The \(\alpha=1,r=8\) inverse-Laplace density is negative at \(t=10\pi\), refuting the conjecture.. Main-text result: \ref{res:not-baskakov-higher-power-even-conjecture}.
\item \textbf{APP-0027.} Ma-Weigert derivative regions form a descending chain.  Source problem: Prove the Ma-Weigert derivative-region descending-chain assertion for log-functions.. Source reference: \cite{app-app-0027-source}. Resolution: Tail vanishing and integration show \(D\_{k+1}\subseteq D\_k\) for every \(k,n\).. Main-text result: \ref{res:ma-weigert-log-function-dk-chain-conjecture}.
\item \textbf{APP-0028.} Qi-Agarwal/Yin divisor-polygamma parity is corrected.  Source problem: Decide the Qi-Agarwal/Yin divisor-polygamma parity problem.. Source reference: \cite{app-app-0028-source}. Resolution: Odd divisor-polygamma sums are completely monotone, but the even non-CM clause is false because \(f\_2\) is completely monotone.. Main-text result: \ref{res:qa-divisor-polygamma-even-claim-refuted}.
\item \textbf{APP-0029.} Bulboaca-Zayed Gamma quotient has one critical point.  Source problem: Prove the Bulboaca-Zayed analytic one-critical-point pattern for their Gamma quotient.. Source reference: \cite{app-app-0029-source}. Resolution: A ratio-kernel single-crossing proof on \([1,8]\), plus source and right-tail intervals, gives the unique zero of \(F'\).. Main-text result: \ref{res:bulboaca-zayed-gamma-quotient-full-monotonicity}.
\item \textbf{APP-0030.} Ramanujan integral Turan window is false.  Source problem: Decide the Mishra-Swaminathan Ramanujan integral Turan-window complete-monotonicity problem.. Source reference: \cite{app-app-0030-source}. Resolution: A logarithmic-density moment-ratio asymptotic makes \(H\_2(e^{-L};1/2-1/(4L))<0\), refuting the interval.. Main-text result: \ref{res:ramanujan-turan-window-negative-answer}.
\item \textbf{APP-0031.} Baricz gamma-quotient Bernstein test is false.  Source problem: Decide whether Baricz's gamma-quotient is Bernstein for every \(a,b>0\).. Source reference: \cite{app-app-0031-source}. Resolution: A single instance \((a,b)=(2,3)\) gives \(g'''(1)<0\), so the universal claim is false.. Main-text result: \ref{res:baricz-gamma-quotient-bf-forall-negative-answer}.
\item \textbf{APP-0032.} From's all-\(L\) Mills-ratio bound family is explicit.  Source problem: Solve From's open all-\(L\) Mills-ratio bound problem from Remark 6.5.. Source reference: \cite{app-app-0032-source}. Resolution: A determinant-to-ratio reduction gives a single explicit quadratic bound for every \(L\), and alternating upper/lower inequalities follow uniformly.. Main-text result: \ref{res:from-mills-all-l-alternating-r-bound-family}.
\item \textbf{APP-0033.} Ramanujan antiderivative is complete Bernstein.  Source problem: Resolve whether the Ramanujan integral antiderivative is a complete Bernstein function.. Source reference: \cite{app-app-0033-source}. Resolution: The Ramanujan density is shown completely monotone, yielding complete-Bernstein structure for the antiderivative.. Main-text result: \ref{res:ramanujan-antiderivative-complete-bernstein}.
\item \textbf{APP-0034.} Bulboaca-Zayed gamma-quotient monotonicity solved.  Source problem: Solve the Bulboaca-Zayed one-critical-point monotonicity statement for the gamma quotient.. Source reference: \cite{app-app-0034-source}. Resolution: A ratio-kernel signature plus source interval/tail controls proves the claimed sign pattern and unique critical point.. Main-text result: \ref{res:bulboaca-zayed-gamma-quotient-full-monotonicity}.
\item \textbf{APP-0035.} Keady self-bijection inverse-CM question is negative.  Source problem: Ask whether a CM self-bijection \((0,\infty)\to(0,\infty)\) must have a CM inverse (Keady Q3).. Source reference: \cite{app-app-0035-source}. Resolution: A local CM decreasing CM bijection is built whose inverse violates the third-derivative sign condition.. Main-text result: \ref{res:keady-q3-self-bijection-inverse-cm-negative-answer}.
\item \textbf{APP-0036.} Baskakov even-line complete-monotonicity fails at \(r=8\).  Source problem: Decide the even-line Baskakov complete-monotonicity conjecture at \(r=8\).. Source reference: \cite{app-app-0036-source}. Resolution: The \(\alpha=1\), \(r=8\) inverse-Laplace density changes sign.. Main-text result: \ref{res:not-baskakov-higher-power-even-conjecture}.
\item \textbf{APP-0037.} Du-Wang \(h\_3\) is increasing exactly on \(\frac12\le a\le1\).  Source problem: Classify Du-Wang \(h\_3\) monotonicity on \((0,\infty)\).. Source reference: \cite{app-app-0037-source}. Resolution: \(h\_3\) is monotone exactly for \(1/2\le a\le1\).. Main-text result: \ref{res:du-wang-h3-open-problem-2-classification}.
\item \textbf{APP-0038.} Ma-Weigert derivative-sign regions form a chain.  Source problem: Prove the Ma-Weigert derivative-region descending chain.. Source reference: \cite{app-app-0038-source}. Resolution: Derivative-sign identities yield \(D\_{k+1}\subseteq D\_k\) for \(k\ge0\).. Main-text result: \ref{res:ma-weigert-log-function-dk-chain-conjecture}.
\item \textbf{APP-0039.} Ramanujan Turan-window complete-monotonicity interval is false.  Source problem: Decide the Ramanujan Turan-window complete-monotonicity interval.. Source reference: \cite{app-app-0039-source}. Resolution: A sharper asymptotic refutes the claimed interval.. Main-text result: \ref{res:ramanujan-turan-window-negative-answer}.
\item \textbf{APP-0040.} Yang-Tian Bessel-\(W\) power-Bernstein conjecture is true.  Source problem: Solve Yang and Tian's Bessel-\(W\) power-Bernstein conjecture on the source interval.. Source reference: \cite{app-app-0040-source}. Resolution: The conjecture is proved on the source interval.. Main-text result: \ref{res:yt-bessel-w-power-bernstein-conjecture}.
\item \textbf{APP-0041.} Qi \(h\_\lambda\) degree-four complete-monotonicity conjecture is false.  Source problem: Settle Qi's degree-4 complete-monotonicity conjecture for \(h\_\lambda\).. Source reference: \cite{app-app-0041-source}. Resolution: Small-\(x\) derivative tests for \(h\_\lambda\) and \(-h\_\mu\) give positive sign on conjectured ranges.. Main-text result: \ref{res:ma-weigert-log-function-dk-chain-conjecture}.
\item \textbf{APP-0042.} Modified Bessel square-root log-concavity fails.  Source problem: Is \(u\mapsto \sqrt{u}I\_\nu(u)\) strictly log-concave for all \(\nu\ge0\)?. Source reference: \cite{app-app-0042-source}. Resolution: A concrete counterexample at \((\nu,u)=(0,10)\) gives positive logarithmic second derivative.. Main-text result: \ref{res:not-bessel-i-sqrt-log-concavity-nu-ge-0}.
\item \textbf{APP-0044.} Bessel ratio quadratic lower bound is false.  Source problem: Prove \(q\_\nu(u)^2+(2\nu+1)q\_\nu(u)/u+(\nu+1/2)/u^2>1\) for all \(\nu\ge0,u>0\), with \(q\_\nu(u)=I\_{\nu+1}(u)/I\_\nu(u)\).. Source reference: \cite{app-app-0044-source}. Resolution: A concrete counterexample at \((\nu,u)=(0,10)\) gives the opposite strict inequality.. Main-text result: \ref{res:from-mills-all-l-moment-ratio-quadratic-bound}.
\end{enumerate}

\section{Other applications pending integration}
The following APP labels are preserved because their public numbering has been assigned, but they are not used as theorem-level consequences of the present Theory.  Each entry remains available for future ingestion once a Theory bridge is recorded.
\begin{enumerate}[leftmargin=*]
\item \textbf{APP-0043.} Riccati log-concavity inequality for \(I\_\nu\) is false.  Source problem: Prove the universal Riccati log-concavity inequality for \(r\_\nu(u)=I\_\nu'(u)/I\_\nu(u)\).. Source reference: \cite{app-app-0043-source}.  Recorded resolution: The inequality is violated at \((\nu,u)=(0,10)\) using explicit rational bounds..
\item \textbf{APP-0045.} Three-regime Bessel log-concavity certificate route is refuted.  Source problem: Give a three-regime certificate for universal \(I\_\nu\) Riccati log-concavity.. Source reference: \cite{app-app-0045-source}.  Recorded resolution: The source inequality is already false globally, so the proposed split-regime certificate cannot exist..
\end{enumerate}

\begin{thebibliography}{74}
\bibitem{app-app-0001-source} Horst Alzer and Man Kam Kwong, "On the concavity and convexity of \(1/\zeta\)", International Journal of Number Theory, Vol. 21, No. 8 (2025), 1825-1835. DOI: \href{https://doi.org/10.1142/S1793042125500897}{DOI}.
\bibitem{app-app-0002-source} Kwara Nantomah, "Open Problem on Riemann Zeta Function", ResearchGate problem note, October 2024. \href{https://www.researchgate.net/publication/384676538_Open_Problem_on_Riemann_Zeta_Function}{ResearchGate}.
\bibitem{app-app-0003-source} Banyat Sroysang, "Two Inequalities for the Riemann Zeta Functions", Mathematica Aeterna, Vol. 3, No. 1 (2013), 21-24. \href{https://arastirmax.com/en/system/files/dergiler/135290/makaleler/3/1/arastirmax-two-inequalities-riemann-zeta-functions.pdf}{PDF}.
\bibitem{app-app-0004-source} Feng Qi, Dongkyu Lim, and Kwara Nantomah, "Monotonicity and positivity of several functions involving ratios and products of polygamma functions", Journal of Inequalities and Applications 2025, article 5. DOI: \href{https://doi.org/10.1186/s13660-024-03245-8}{DOI}.
\bibitem{app-app-0005-source} Donggyun Kim and Kyunghwan Song, "The inverses of tails of the Riemann zeta function", Journal of Inequalities and Applications 2018, article 157. \href{https://link.springer.com/article/10.1186/s13660-018-1743-6}{https://link.springer.com/article/10.1186/s13660-018-1743-6}; Zhenjiang Pan and Zhengang Wu, "The inverse of tails of Riemann zeta function, Hurwitz zeta function and Dirichlet L-function", AIMS Mathematics 9(6), 16564-16585, 2024. DOI: \href{https://doi.org/10.3934/math.2024803}{DOI}.
\bibitem{app-app-0006-source} Donggyun Kim and Kyunghwan Song, "The inverses of tails of the Riemann zeta function", Journal of Inequalities and Applications 2018, article 157. \href{https://link.springer.com/article/10.1186/s13660-018-1743-6}{https://link.springer.com/article/10.1186/s13660-018-1743-6}; Zhenjiang Pan and Zhengang Wu, "The inverse of tails of Riemann zeta function, Hurwitz zeta function and Dirichlet L-function", AIMS Mathematics 9(6), 16564-16585, 2024. DOI: \href{https://doi.org/10.3934/math.2024803}{DOI}.
\bibitem{app-app-0007-source} Feng Qi, Dongkyu Lim, and Kwara Nantomah, "Monotonicity and positivity of several functions involving ratios and products of polygamma functions", Journal of Inequalities and Applications 2025, article 5. DOI: \href{https://doi.org/10.1186/s13660-024-03245-8}{DOI}.
\bibitem{app-app-0008-source} Feng Qi, Dongkyu Lim, and Kwara Nantomah, "Monotonicity and positivity of several functions involving ratios and products of polygamma functions", Journal of Inequalities and Applications 2025, article 5. DOI: \href{https://doi.org/10.1186/s13660-024-03245-8}{DOI}.
\bibitem{app-app-0009-source} Teodor Bulboaca and Hanaa M. Zayed, "Monotonic nature of the Gamma function", Journal of Inequalities and Applications 2026, article 27. DOI: \href{https://doi.org/10.1186/s13660-025-03425-0}{DOI}.
\bibitem{app-app-0010-source} Li Yin and Jumei Zhang, "On some properties of special functions involving \(k\)-gamma and \(k\)-digamma functions", \href{https://arxiv.org/abs/2502.15852}{arXiv:2502.15852}. \href{https://arxiv.org/abs/2502.15852}{arXiv}.
\bibitem{app-app-0011-source} Feng Qi, Dongkyu Lim, and Kwara Nantomah, "Monotonicity and positivity of several functions involving ratios and products of polygamma functions", Journal of Inequalities and Applications 2025, article 5. DOI: \href{https://doi.org/10.1186/s13660-024-03245-8}{DOI}.
\bibitem{app-app-0012-source} Arpad Baricz, "Turan type inequalities for some special functions", Ph.D. thesis, University of Debrecen, 2008. \href{https://dea.lib.unideb.hu/bitstreams/bdd4469f-1b1e-4996-9b90-73f5e1a6b0e8/download}{PDF}.
\bibitem{app-app-0013-source} Li Yin, "Complete monotonicity of a function involving the \((p,k)\)-digamma function", International Journal of Open Problems in Computer Mathematics 11(2), 103-108, 2018. \href{https://www.ijopcm.org/Vol/2018/2.9.pdf}{PDF}.
\bibitem{app-app-0014-source} Mishra and Swaminathan, Inequalities involving a Ramanujan Integral, \href{https://arxiv.org/abs/2511.07443}{arXiv:2511.07443}..
\bibitem{app-app-0015-source} Sibisi, A Probabilistic Perspective on Feller, Pollard and the Complete Monotonicity of the Mittag-Leffler Function, \href{https://arxiv.org/abs/2301.01466}{arXiv:2301.01466}..
\bibitem{app-app-0016-source} From, Some new upper and lower bounds for the Mills ratio, J. Math. Anal. Appl. 486 (2020), 123872..
\bibitem{app-app-0017-source} Baricz, Turan type inequalities for hypergeometric functions, Proc. Amer. Math. Soc. 136 (2008), 3223--3229..
\bibitem{app-app-0018-source} Qi and Guo, Complete Monotonicities of Functions Involving Gamma and Digamma Functions, RGMIA Res. Rep. Coll. 7 (2004), Article 8..
\bibitem{app-app-0019-source} Wakrim, The W-Operator: A Volterra Fractional Time Operator with Non-Bernstein Symbol, \href{https://arxiv.org/abs/2601.02876}{arXiv:2601.02876}..
\bibitem{app-app-0020-source} Qi, Completely monotonic degree of a function involving trigamma and tetragamma functions, AIMS Math. 5 (2020), 3391--3407..
\bibitem{app-app-0021-source} Szabo, Completely monotone functions in general and some applications, \href{https://arxiv.org/abs/2411.17670}{arXiv:2411.17670}..
\bibitem{app-app-0022-source} Chen and Choi, Completely monotonic functions related to Gurland's ratio for the gamma function, Math. Inequal. Appl. 20 (2017), 651--659..
\bibitem{app-app-0023-source} Sokal, Real-variables characterization of generalized Stieltjes functions, Expo. Math. 28 (2010), 179--185..
\bibitem{app-app-0024-source} Thomas Simon, Moment problems related to Bernstein functions, Ann. Fac. Sci. Toulouse Math. 29 (2020), 577--594..
\bibitem{app-app-0025-source} Peipei Du and Gendi Wang, Monotonicity, convexity, and inequalities for functions involving gamma function, \href{https://arxiv.org/abs/2205.12530}{arXiv:2205.12530}..
\bibitem{app-app-0026-source} Ulrich Abel, Wolfgang Gawronski, and Thorsten Neuschel, Complete Monotonicity and Zeros of Sums of Squared Baskakov Functions, \href{https://arxiv.org/abs/1411.7945}{arXiv:1411.7945}..
\bibitem{app-app-0027-source} Rourou Ma and Julian Weigert, Complete monotonicity of log-functions, \href{https://arxiv.org/abs/2505.04225}{arXiv:2505.04225}..
\bibitem{app-app-0028-source} Feng Qi and Ravi P. Agarwal, On complete monotonicity for several classes of functions related to ratios of gamma functions, Journal of Inequalities and Applications 2019, article 36..
\bibitem{app-app-0029-source} Teodor Bulboaca and Hanaa M. Zayed, Monotonic nature of the Gamma function, Journal of Inequalities and Applications 2026, article 27..
\bibitem{app-app-0030-source} D. Mishra and A. Swaminathan, Inequalities involving a Ramanujan integral, \href{https://arxiv.org/abs/2511.07443}{arXiv:2511.07443}..
\bibitem{app-app-0031-source} Arpad Baricz, Turan type inequalities for hypergeometric functions, Proc. Amer. Math. Soc. 136 (2008), 3223--3229..
\bibitem{app-app-0032-source} S. G. From, Some new upper and lower bounds for the Mills ratio, J. Math. Anal. Appl. 486 (2020), 123872..
\bibitem{app-app-0033-source} Deepshikha Mishra and A. Swaminathan, Inequalities involving a Ramanujan Integral, \href{https://arxiv.org/abs/2511.07443}{arXiv:2511.07443}..
\bibitem{app-app-0034-source} Teodor Bulboaca and Hanaa M. Zayed, Monotonic nature of the Gamma function, Journal of Inequalities and Applications 2026, article 27..
\bibitem{app-app-0035-source} G. Keady, N. Khajohnsaksumeth, and B. Wiwatanapataphee, On functions and inverses, both positive, decreasing and convex: and Stieltjes functions, Cogent Mathematics \& Statistics 5:1 (2018), 1477543..
\bibitem{app-app-0036-source} Ulrich Abel, Wolfgang Gawronski, and Thorsten Neuschel, Complete Monotonicity and Zeros of Sums of Squared Baskakov Functions, \href{https://arxiv.org/abs/1411.7945}{arXiv:1411.7945}..
\bibitem{app-app-0037-source} Peipei Du and Gendi Wang, Monotonicity, convexity, and inequalities for functions involving gamma function, \href{https://arxiv.org/abs/2205.12530}{arXiv:2205.12530}..
\bibitem{app-app-0038-source} Rourou Ma and Julian Weigert, Complete monotonicity of log-functions, \href{https://arxiv.org/abs/2505.04225}{arXiv:2505.04225}..
\bibitem{app-app-0039-source} Deepshikha Mishra and A. Swaminathan, Inequalities involving a Ramanujan integral, \href{https://arxiv.org/abs/2511.07443}{arXiv:2511.07443}..
\bibitem{app-app-0040-source} Zhen-Hang Yang and Jing-Feng Tian, Convexity of ratios of the modified Bessel functions of the first kind with applications, Revista Matematica Complutense, 2022..
\bibitem{app-app-0041-source} Feng Qi, "Completely monotonic degree of a function involving trigamma and tetragamma functions", AIMS Mathematics 5 (2020), 3391-3407, \href{https://doi.org/10.3934/math.2020219}{DOI: 10.3934/math.2020219}.
\bibitem{app-app-0042-source} Mihály Baricz, Pietro Ponnusamy, and Matti Vuorinen, Functional inequalities for modified Bessel functions, \href{https://arxiv.org/abs/1009.4814}{arXiv:1009.4814}, \href{https://doi.org/10.1016/j.exmath.2011.07.001}{DOI: 10.1016/j.exmath.2011.07.001}.
\bibitem{app-app-0043-source} Mihály Baricz, Pietro Ponnusamy, and Matti Vuorinen, Functional inequalities for modified Bessel functions, \href{https://arxiv.org/abs/1009.4814}{arXiv:1009.4814}, \href{https://doi.org/10.1016/j.exmath.2011.07.001}{DOI: 10.1016/j.exmath.2011.07.001}.
\bibitem{app-app-0044-source} Mihály Baricz, Pietro Ponnusamy, and Matti Vuorinen, Functional inequalities for modified Bessel functions, \href{https://arxiv.org/abs/1009.4814}{arXiv:1009.4814}, \href{https://doi.org/10.1016/j.exmath.2011.07.001}{DOI: 10.1016/j.exmath.2011.07.001}.
\bibitem{app-app-0045-source} Mihály Baricz, Pietro Ponnusamy, and Matti Vuorinen, Functional inequalities for modified Bessel functions, \href{https://arxiv.org/abs/1009.4814}{arXiv:1009.4814}, \href{https://doi.org/10.1016/j.exmath.2011.07.001}{DOI: 10.1016/j.exmath.2011.07.001}.
\bibitem{nantomah}
K. Nantomah,
\emph{Open Problem on Riemann Zeta Function},
ResearchGate problem note, October 2024.
\url{https://www.researchgate.net/publication/384676538_Open_Problem_on_Riemann_Zeta_Function}.
\bibitem{sroysang}
B. Sroysang,
\emph{Two Inequalities for the Riemann Zeta Functions},
Mathematica Aeterna 3, no. 1 (2013), 21--24.
\url{https://arastirmax.com/en/system/files/dergiler/135290/makaleler/3/1/arastirmax-two-inequalities-riemann-zeta-functions.pdf}.
\bibitem{alzer-kwong}
H. Alzer and M. K. Kwong,
\emph{On the concavity and convexity of \(1/\zeta\)},
International Journal of Number Theory 21, no. 8 (2025), 1825--1835.
\url{https://doi.org/10.1142/S1793042125500897}.
\bibitem{qi-lim-nantomah-polygamma}
F. Qi, D. Lim, and K. Nantomah,
\emph{Monotonicity and positivity of several functions involving ratios and products of polygamma functions},
Journal of Inequalities and Applications 2025, article 5.
\url{https://doi.org/10.1186/s13660-024-03245-8}.
\bibitem{bulboaca-zayed-gamma}
T. Bulboaca and H. M. Zayed,
\emph{Monotonic nature of the Gamma function},
Journal of Inequalities and Applications 2026, article 27.
\url{https://doi.org/10.1186/s13660-025-03425-0}.
\bibitem{kim-song-tail}
D. Kim and K. Song,
\emph{The inverses of tails of the Riemann zeta function},
Journal of Inequalities and Applications 2018, article 157.
\url{https://doi.org/10.1186/s13660-018-1743-6}.
\bibitem{pan-wu-tail}
Z. Pan and Z. Wu,
\emph{The inverse of tails of Riemann zeta function, Hurwitz zeta function and Dirichlet L-function},
AIMS Mathematics 9, no. 6 (2024), 16564--16585.
\url{https://doi.org/10.3934/math.2024803}.
\bibitem{yin-zhang-k-beta}
L. Yin and J. Zhang,
\emph{On some properties of special functions involving \(k\)-gamma and \(k\)-digamma functions},
arXiv:2502.15852, 2025.
\url{https://arxiv.org/abs/2502.15852}.
\bibitem{baricz-turan-thesis}
A. Baricz,
\emph{Turan type inequalities for some special functions},
Ph.D. thesis, University of Debrecen, 2008.
\url{https://dea.lib.unideb.hu/bitstreams/bdd4469f-1b1e-4996-9b90-73f5e1a6b0e8/download}.
\bibitem{yin-pk-digamma}
L. Yin,
\emph{Complete monotonicity of a function involving the \((p,k)\)-digamma function},
International Journal of Open Problems in Computer Mathematics 11(2) (2018), 103--108.
\url{https://www.ijopcm.org/Vol/2018/2.9.pdf}.
\bibitem{ma-weigert-log-functions}
R. Ma and J. Weigert,
\emph{Complete monotonicity of log-functions},
arXiv:2505.04225.
\url{https://arxiv.org/abs/2505.04225}.
\bibitem{qi-agarwal-gamma-ratios}
F. Qi and R. P. Agarwal,
\emph{On complete monotonicity for several classes of functions related to
ratios of gamma functions},
Journal of Inequalities and Applications 2019, article 36.
\url{https://doi.org/10.1186/s13660-019-1976-z}.
\bibitem{mishra-swaminathan-ramanujan}
D. Mishra and A. Swaminathan,
\newblock Inequalities involving a Ramanujan integral,
\newblock arXiv:2511.07443.
\bibitem{sibisi-prabhakar}
N. K. Sibisi,
\newblock A probabilistic perspective on Feller, Pollard and the complete
  monotonicity of the Mittag-Leffler function,
\newblock arXiv:2301.01466.
\bibitem{from-mills}
S. G. From,
\newblock Some new upper and lower bounds for the Mills ratio,
\newblock J. Math. Anal. Appl. 486 (2020), 123872.
\bibitem{baricz-turan}
A. Baricz,
\newblock Turan type inequalities for hypergeometric functions,
\newblock Proc. Amer. Math. Soc. 136 (2008), 3223--3229.
\bibitem{qi-guo-2004}
F. Qi and B.-N. Guo,
\newblock Complete monotonicities of functions involving gamma and digamma
  functions,
\newblock RGMIA Res. Rep. Coll. 7 (2004), Article 8.
\bibitem{wakrim-w}
M. Wakrim,
\newblock The W-operator: a Volterra fractional time operator with
  non-Bernstein symbol,
\newblock arXiv:2601.02876.
\bibitem{qi-hlambda}
F. Qi,
\newblock Completely monotonic degree of a function involving trigamma and
  tetragamma functions,
\newblock AIMS Math. 5 (2020), 3391--3407.
\bibitem{szabo-cm}
V. E. S. Szabo,
\newblock Completely monotone functions in general and some applications,
\newblock arXiv:2411.17670.
\bibitem{chen-choi-gurland}
C.-P. Chen and J. Choi,
\newblock Completely monotonic functions related to Gurland's ratio for the
  gamma function,
\newblock Math. Inequal. Appl. 20 (2017), 651--659.
\bibitem{sokal-gs}
A. D. Sokal,
\newblock Real-variables characterization of generalized Stieltjes functions,
\newblock Expo. Math. 28 (2010), 179--185.
\bibitem{simon-moment}
T. Simon,
\newblock Moment problems related to Bernstein functions,
\newblock Ann. Fac. Sci. Toulouse Math. 29 (2020), 577--594.
\url{https://doi.org/10.5802/afst.1640}.
\bibitem{du-wang-h3}
P. Du and G. Wang,
\newblock Monotonicity, convexity, and inequalities for functions involving
  gamma function,
\newblock arXiv:2205.12530.
\url{https://arxiv.org/abs/2205.12530}.
\bibitem{abel-gawronski-neuschel}
U. Abel, W. Gawronski, and T. Neuschel,
\newblock Complete monotonicity and zeros of sums of squared Baskakov
  functions,
\newblock Appl. Math. Comput. 258 (2015), 130--137.
\url{https://arxiv.org/abs/1411.7945}.
\bibitem{apostol}
T. M. Apostol,
\emph{Introduction to Analytic Number Theory},
Undergraduate Texts in Mathematics, Springer, 1976.
\url{https://doi.org/10.1007/978-1-4757-5579-4}.
\bibitem{davenport}
H. Davenport,
\emph{Multiplicative Number Theory},
3rd ed., revised by H. L. Montgomery, Graduate Texts in Mathematics, vol. 74, Springer, 2000.
\url{https://doi.org/10.1007/978-1-4757-5927-3}.
\bibitem{titchmarsh}
E. C. Titchmarsh,
\emph{The Theory of the Riemann Zeta-function},
2nd ed., revised by D. R. Heath-Brown, Oxford University Press, 1986.
\url{https://global.oup.com/academic/product/the-theory-of-the-riemann-zeta-function-9780198533696}.
\bibitem{cover-thomas}
T. M. Cover and J. A. Thomas,
\emph{Elements of Information Theory},
2nd ed., Wiley-Interscience, 2006.
\url{https://doi.org/10.1002/047174882X}.
\end{thebibliography}

\end{document}
